
\Data\ID{1003108601}
\Updated{2023-09-10 11:09:33}
\Level{A}
\SubArea{Rational functions}
\LANGen{
\Question{
Peter drove from Ostrava to Warsaw. He drove at an average speed of \( 104 \)  kilometres per hour and reached Warsaw in \( 4 \) hours. Select the function that describes  the dependence of Peter’s driving time \( t \) on the average speed \( v \) of the car. (The driving time \( t \) is given in hours and the average speed \( v \) is given in kilometres per hour.)}
{\( t=\frac{416}v\text{ ,}\ v\in(0;\infty) \)}{\( t=\frac{26}v\text{ ,}\ v\in(0;\infty) \)}{\( t=\frac v{26}\text{ ,}\ v\in(0;\infty) \)}{\( t=\frac{104}v\text{ ,}\ v\in(0;\infty) \)}{}{}}
\LANGpl{
\Question{
Piotr podróżował z Ostrawy  do Warszawy. Jechał ze średnią prędkością \( 104 \)  kilometrów na godzinę i dotarł do Warszawy w  \( 4 \) godziny. Wybierz funkcję, która przedstawia zależność czasu jazdy \( t \) i średniej prędkości \( v \) samochodu. (Czas podróży \( t \) podany jest w godzinach, a prędkość \( v \) w kilometrach na godzinę.)}
{\( t=\frac{416}v\text{ ,}\ v\in(0;\infty) \)}{\( t=\frac{26}v\text{ ,}\ v\in(0;\infty) \)}{\( t=\frac v{26}\text{ ,}\ v\in(0;\infty) \)}{\( t=\frac{104}v\text{ ,}\ v\in(0;\infty) \)}{}{}}
\LANGcs{
\Question{
Petr jel autem z Ostravy do Varšavy. Jel průměrnou rychlostí  \( 104 \)  kilometrů za hodinu a do Varšavy  dojel za  \( 4 \) hodiny. Vyberte funkci, která vyjadřuje závislost času jízdy \( t \) na průměrné rychlosti \( v \) Petrova auta. Čas \( t \) je udáván v hodinách a průměrná rychlost \( v \) v kilometrech za hodinu.}
{\( t=\frac{416}v\text{ ,}\ v\in(0;\infty) \)}{\( t=\frac{26}v\text{ ,}\ v\in(0;\infty) \)}{\( t=\frac v{26}\text{ ,}\ v\in(0;\infty) \)}{\( t=\frac{104}v\text{ ,}\ v\in(0;\infty) \)}{}{}}
\LANGsk{
\Question{
Peter cestoval autom z Ostravy do Varšavy priemernou rýchlosťou  \( 104 \)  kilometrov za hodinu a do Varšavy prišiel za  \( 4 \) hodiny. Vyberte funkciu, ktorá vyjadruje závislosť času jazdy \( t \) na priemernej rýchlosti \( v \) Petrovho auta. Čas \( t \) je udávame v hodinách a priemernú rýchlosť \( v \) v kilometroch za hodinu.}
{\( t=\frac{416}v\text{ ,}\ v\in(0;\infty) \)}{\( t=\frac{26}v\text{ ,}\ v\in(0;\infty) \)}{\( t=\frac v{26}\text{ ,}\ v\in(0;\infty) \)}{\( t=\frac{104}v\text{ ,}\ v\in(0;\infty) \)}{}{}}
\LANGes{
\Question{
Peter condujo de Ostrava a Varsovia. Condujo a una velocidad media de \( 104 \)  kilómetros por hora y llegó a Varsovia en \( 4 \) horas. Elige la función que describe la dependencia del tiempo de conducción de Peter \( t \) respecto de la velocidad media \( v \) del coche.  (El tiempo de conducción \( t \) se da en horas y la velocidad \( v \) se da en kilómetros por hora).}
{\( t=\frac{416}v\text{ ,}\ v\in(0;\infty) \)}{\( t=\frac{26}v\text{ ,}\ v\in(0;\infty) \)}{\( t=\frac v{26}\text{ ,}\ v\in(0;\infty) \)}{\( t=\frac{104}v\text{ ,}\ v\in(0;\infty) \)}{}{}}
\End

\Data\ID{1003108603}
\Updated{2021-04-09 07:04:12}
\Level{A}
\SubArea{Rational functions}
\LANGen{
\Question{
The fuel consumption of Skoda Fabia \( 1.4 \) MPi/\( 44\,\mathrm{kW} \) stated by the manufacturer varies between  \( 5.5\,\mathrm{l} \) / \( 100\,\mathrm{km} \) (out of a town) to  \( 9.6\,\mathrm{l} \) / \( 100\,\mathrm{km} \) (in a town). Suppose the car’s fuel tank of capacity \( 45\,\mathrm{l} \) is filled completely. Choose the function describing relation between the distance \( p \) in \( \mathrm{km} \) that the car can travel without tanking on the fuel consumption \( s \).}
{\( f\colon p=\frac{4\:500}s;\ s\in[5.5;9.6] \)}{\( h\colon p=\frac{45}s;\ s\in[5.5;9.6] \)}{\( r\colon p=\frac s{0.45};\ s\in[5.5;9.6] \)}{\( g\colon p=45\cdot s;\ s\in[5.5;9.6] \)}{}{}}
\LANGpl{
\Question{
Zużycie paliwa przez samochód marki Skoda Fabia \( 1.4 \) MPi/\( 44\,\mathrm{kW} \) określone przez producenta waha się miedzy  \( 5{,}5\,\mathrm{l} \) / \( 100\,\mathrm{km} \) (poza miastem) a \( 9{,}6\,\mathrm{l} \) / \( 100\,\mathrm{km} \) (w mieście). Załóżmy, że bak o pojemności \( 45\,\mathrm{l} \) uzupełniono do pełna. Wybierz funkcję, która określa zależność pomiędzy odległością  \( p \) w \( \mathrm{km} \) jaką może pokonać samochód bez tankowania a zużyciem  paliwa \( s \).}
{\( f\colon p=\frac{4\:500}s;\ s\in\langle5{,}5;9{,}6\rangle \)}{\( h\colon p=\frac{45}s;\ s\in\langle5{,}5;9{,}6\rangle \)}{\( r\colon p=\frac s{0{,}45};\ s\in\langle5{,}5;9{,}6\rangle \)}{\( g\colon p=45\cdot s;\ s\in\langle5{,}5;9{,}6\rangle \)}{}{}}
\LANGcs{
\Question{
Výrobcem udávaná spotřeba paliva Škody Fabia \( 1.4 \) MPi/\( 44\,\mathrm{kW} \) se pohybuje v rozmezí od \( 5{,}5\,\mathrm{l} \) / \( 100\,\mathrm{km} \) (mimo město) do  \( 9{,}6\,\mathrm{l} \) / \( 100\,\mathrm{km} \) (ve městě). Objem její palivové nádrže je \( 45\,\mathrm{l} \). Předpokládejme, že Fabia má nádrž zcela naplněnou. Vyberte funkci, která vyjadřuje závislost počtu kilometrů \( p \), které ujede toto auto bez tankování, na spotřebě  paliva \( s \).}
{\( f\colon p=\frac{4\:500}s;\ s\in\langle5{,}5;9{,}6\rangle \)}{\( h\colon p=\frac{45}s;\ s\in\langle5{,}5;9{,}6\rangle \)}{\( r\colon p=\frac s{0{,}45};\ s\in\langle5{,}5;9{,}6\rangle \)}{\( g\colon p=45\cdot s;\ s\in\langle5{,}5;9{,}6\rangle \)}{}{}}
\LANGsk{
\Question{
Výrobcom udávaná spotreba paliva Škody Fabia \( 1.4 \) MPi/\( 44\,\mathrm{kW} \) sa pohybuje v rozmedzí od \( 5{,}5\,\mathrm{l} \) / \( 100\,\mathrm{km} \) (mimo mesto) do  \( 9{,}6\,\mathrm{l} \) / \( 100\,\mathrm{km} \) (v meste). Objem palivovej nádrže je \( 45\,\mathrm{l} \). Prepokladajme, že nádrž Škody Fabia je úplne plná. Vyberte funkciu, ktorá vyjadruje závislosť počtu kilometrov \( p \), ktoré prejde toto auto bez tankovania, na spotrebe  paliva \( s \).}
{\( f\colon p=\frac{4\:500}s;\ s\in\langle5{,}5;9{,}6\rangle \)}{\( h\colon p=\frac{45}s;\ s\in\langle5{,}5;9{,}6\rangle \)}{\( r\colon p=\frac s{0{,}45};\ s\in\langle5{,}5;9{,}6\rangle \)}{\( g\colon p=45\cdot s;\ s\in\langle5{,}5;9{,}6\rangle \)}{}{}}
\LANGes{
\Question{
El consumo de combustible del Skoda Fabia modelo \( 1.4 \) MPi/\( 44\,\mathrm{kW} \) según indica el fabricante varía entre \( 5.5\,\mathrm{l} \) / \( 100\,\mathrm{km} \) (fuera de ciudad) a  \( 9.6\,\mathrm{l} \) / \( 100\,\mathrm{km} \) (en ciudad). Imagina que el tanque de combustible del coche con una capacidad de \( 45\,\mathrm{l} \) está completamente lleno. Elige la función que describa la relación entre la distancia \( p \) en \( \mathrm{km} \)que el coche puede viajar sin llenar el tanque respecto al consumo de combustible \( s \).}
{\( f\colon p=\frac{4\:500}s;\ s\in[5.5;9.6] \)}{\( h\colon p=\frac{45}s;\ s\in[5.5;9.6] \)}{\( r\colon p=\frac s{0.45};\ s\in[5.5;9.6] \)}{\( g\colon p=45\cdot s;\ s\in[5.5;9.6] \)}{}{}}
\End

\Data\ID{1003109501}
\Updated{2022-04-23 05:04:42}
\Level{A}
\SubArea{Rational functions}
\LANGen{
\Question{
Let \( f(x)=-\frac1{2x} \). Find the false statement.}
{The function \( f \) is increasing.}{The range of \( f \) is \( \left(-\infty;0\right)\cup\left(0;\infty\right) \).}{The function \( f \) is odd.}{The function \( f \) is not bounded.}{}{}}
\LANGpl{
\Question{
Niech  \( f(x)=-\frac1{2x} \). Które z poniższych zdań jest fałszywe?}
{Funkcja \( f \) jest rosnąca.}{Zakres funkcji \( f \) to \( \left(-\infty;0\right)\cup\left(0;\infty\right) \).}{Funkcja  \( f \) jest nieparzysta.}{Funkcja \( f \) nie jest ograniczona.}{}{}}
\LANGcs{
\Question{
Funkce \( f \) je dána předpisem \( f(x)=-\frac1{2x} \). Vyberte nepravdivý výrok.}
{Funkce \( f \) je rostoucí.}{Oborem hodnot funkce \( f \) je množina \( \left(-\infty;0\right)\cup\left(0;\infty\right) \).}{Funkce \( f \) je lichá.}{Funkce \( f \) není omezená.}{}{}}
\LANGsk{
\Question{
Daná je funkcia predpisom \( f(x)=-\frac1{2x} \). Vyberte nepravdivé tvrdenie o funkcii \( f \).}
{Funkcia \( f \) je klesajúca.}{Obor hodnôt funkcie \( f \) je \( \left(-\infty;0\right)\cup\left(0;\infty\right) \).}{Funkcia \( f \) je nepárna.}{Funkcia \( f \) je neohraničená.}{}{}}
\LANGes{
\Question{
Sea\( f(x)=-\frac1{2x} \). Determina la proposición  falsa.}
{La función\( f \) es creciente.}{El rango de la función \( f \) es \( \left(-\infty;0\right)\cup\left(0;\infty\right) \).}{La función\( f \) es impar.}{La función \( f \) no está acotada.}{}{}}
\End

\Data\ID{1003109502}
\Updated{2020-07-15 03:07:21}
\Level{A}
\SubArea{Rational functions}
\LANGen{
\Question{
Let \( f(x)=-\frac2x\text{, }x\in[-2;0)\cup(0;\infty) \). Find the true statement.}
{The function \( f \) is injective (one-to-one).}{The function \( f \) has minimum at \( x=-2 \).}{The range of \( f \) is \( [0;1) \).}{The function \( f \) is odd.}{}{}}
\LANGpl{
\Question{
Niech \( f(x)=-\frac2x\text{, }x\in\langle-2;0)\cup(0;\infty) \). Które stwierdzenie jest prawdziwe?}
{Funkcja \( f \) jest iniekcyjna (jeden do jednego).}{Funkcja \( f \) osiąga minimum w \( x=-2 \).}{Zakres funkcji \( f \) to \( \langle0;1) \).}{Funkcja  \( f \) jest nieparzysta.}{}{}}
\LANGcs{
\Question{
Funkce \( f \) je dána předpisem \( f(x)=-\frac2x\text{, }x\in\langle-2;0)\cup(0;\infty) \). Vyberte pravdivý výrok.}
{Funkce \( f \) je prostá.}{Funkce \( f \) má minimum v bodě \( x=-2 \).}{Oborem hodnot funkce \( f \) je interval \( \langle0;1) \).}{Funkce \( f \) je lichá.}{}{}}
\LANGsk{
\Question{
Funkcia \(f\) je daná predpisom \( f(x)=-\frac2x\text{, }x\in\langle-2;0)\cup(0;\infty) \). Vyberte pravdivý výrok.}
{Funkcia \( f \) je prostá.}{Funkcia \( f \) má minimum v bode \( x=-2 \).}{Obor hodnôt funkcie \( f \) je \( \langle0;1) \).}{Funkcia \( f \) je nepárna.}{}{}}
\LANGes{
\Question{
Sea \( f(x)=-\frac2x\text{, }x\in[-2;0)\cup(0;\infty) \).  Determina la proposición verdadera.}
{La función \( f \) es inyectiva (uno a uno).}{La función\( f \) tiene su mínimo en \( x=-2 \).}{El rango de la función \( f \) es \( [0;1) \).}{La función\( f \) es impar.}{}{}}
\End

\Data\ID{1103108602}
\Updated{2021-04-09 07:04:22}
\Level{A}
\SubArea{Rational functions}
\LANGen{
\Question{
In a simple electric circuit a voltage source and a variable resistor  with resistance \( R \) in the range  \(  [1\Omega;10\Omega] \) are connected.  Suppose the source gives fixed voltage of \( 5\,\mathrm{V} \). From the graphs given below select the one that describes the dependence of the electric current \( I \) on the resistance \( R \) in this circuit. (Note: The relationship between electric currant, voltage and resistance is described by Ohm’s law: \( U=RI \).)}
{}{}{}{}{}{}}
\LANGpl{
\Question{
W prostym obwodzie elektrycznym podłączone jest źródło napięcia i rezystor o oporze \( R \) w zakresie   \( \langle1\Omega;10\Omega\rangle \). Załóżmy, że źródło podaje stałe napięcie \( 5\,\mathrm{V} \). Z poniższych wykresów wybierz jeden, który przedstawia zależność pomiędzy prądem elektrycznym \( I \) a oporem \( R \) w tym obwodzie. (Wskazówka: Związek pomiędzy obwodem elektrycznym, napięciem i oporem jest określony jest przez prawo Ohma: \( U=RI \).)}
{}{}{}{}{}{}}
\LANGcs{
\Question{
V jednoduchém elektrickém obvodu je ke zdroji napětí připojen rezistor s proměnným odporem \( R \) o rozsahu \( \langle1\Omega;10\Omega\rangle \). Napětí \( U \) na svorkách zdroje je konstantní o velikosti \( 5\,\mathrm{V} \). Vyberte graf funkce vyjadřující závislost proudu \( I \) v tomto obvodu na odporu rezistoru \( R \). (Nápověda: Vztah mezi elektrickým odporem, napětím a proudem udává Ohmův zákon: \( U=RI \).)}
{}{}{}{}{}{}}
\LANGsk{
\Question{
V jednoduchom elektrickom obvode je pripojený k zdroju napätia rezistor s priemerným odporom \( R \) s rozsahom \( \langle1\Omega;10\Omega\rangle \). Napätie \( U \) na svorkách zdroja je konštantné s veľkosťou \( 5\,\mathrm{V} \). Vyberte graf funkcie, ktorý vyjadruje závislosť prúdu \( I \) v tomto obvode na odpore rezistora \( R \). (Pomôcka: Vzťah medzi elektrickým odporom, napätím a prúdom udáva Ohmov zákon: \( U=RI \).)}
{}{}{}{}{}{}}
\LANGes{
\Question{
Una fuente de voltaje y una resistencia variable con resistencia \(R \) en el rango \([1 \ Omega; 10 \ Omega] \) se conectan en un circuito eléctrico simple. Imagina que la fuente proporciona un voltaje fijo de \(5 \, \mathrm {V} \). De las gráficas que figuran a continuación, elige la que describe la dependencia de la corriente eléctrica \(I \) respecto de la resistencia \(R \) en este circuito. (Nota: la relación entre la corriente eléctrica, el voltaje y la resistencia se describe mediante la ley de Ohm: \(U = RI \)).}
{}{}{}{}{}{}}

\IMGanswer{1103108602a.pdf}
\IMGanswer{1103108602b.pdf}
\IMGanswer{1103108602c.pdf}
\IMGanswer{1103108602d.pdf}\End

\Data\ID{1103108604}
\Updated{2021-04-09 07:04:05}
\Level{A}
\SubArea{Rational functions}
\IMG{1103108604.pdf}
\LANGen{
\Question{
The hall’s floor needs a new tiling. All tiles used will be of the same size. The picture shows the graph of the function describing dependence of the number \( p \) of tiles needed in the hall on the area \( S \) of one tile. What is the area of the hall’s floor?}
{\( 10.5\,\mathrm{m}^2 \)}{\( 1\:050\,\mathrm{m}^2 \)}{\( 2\:100\,\mathrm{m}^2 \)}{\( 42\,\mathrm{m}^2 \)}{}{}}
\LANGpl{
\Question{
Podłoga w przedpokoju wymaga nowej posadzki. Wszystkie zużyte płytki będą tego samego rozmiaru. Rysunek przedstawia wykres funkcji określającej zależność pomiędzy liczbą  płytek  \( p \) potrzebnych do wyłożenia posadzki, a powierzchnią \( S \) jednej płytki. Jaka jest powierzchnia posadzki w przedpokoju?}
{\( 10{,}5\,\mathrm{m}^2 \)}{\( 1\:050\,\mathrm{m}^2 \)}{\( 2\:100\,\mathrm{m}^2 \)}{\( 42\,\mathrm{m}^2 \)}{}{}}
\LANGcs{
\Question{
Podlahu chodby je třeba vydláždit stejnými dlaždicemi. Na obrázku je graf funkce udávající závislost počtu dlaždic \( p \) potřebných pro vydláždění chodby na obsahu dlaždice \( S \). Jaký je obsah podlahy chodby?}
{\( 10{,}5\,\mathrm{m}^2 \)}{\( 1\:050\,\mathrm{m}^2 \)}{\( 2\:100\,\mathrm{m}^2 \)}{\( 42\,\mathrm{m}^2 \)}{}{}}
\LANGsk{
\Question{
Podlahu chodby je potrebné vydláždiť rovnakými dlaždicami. Na obrázku je graf funkcie, ktorý vyjadruje závislosť počtu dlaždíc \( p \) potrebných pre vydláždenie chodby na obsahu dlaždice \( S \). Aký je obsah podlahy chodby?}
{\( 10{,}5\,\mathrm{m}^2 \)}{\( 1\:050\,\mathrm{m}^2 \)}{\( 2\:100\,\mathrm{m}^2 \)}{\( 42\,\mathrm{m}^2 \)}{}{}}
\LANGes{
\Question{
El suelo de una sala necesita nuevas baldosas. Todas las baldosas utilizadas serán del mismo tamaño. La imagen muestra la gráfica de la función que describe la dependencia del número \(p \) de baldosas necesarias en la sala respecto a el área \(S \) de una baldosa. ¿Cuál es el área del suelo de la sala?}
{\( 10.5\,\mathrm{m}^2 \)}{\( 1\:050\,\mathrm{m}^2 \)}{\( 2\:100\,\mathrm{m}^2 \)}{\( 42\,\mathrm{m}^2 \)}{}{}}
\End

\Data\ID{1103124503}
\Updated{2021-04-09 07:04:14}
\Level{A}
\SubArea{Rational functions}
\IMG{1103124503.pdf}
\LANGen{
\Question{
The picture shows graphs of functions:
\[ \begin{aligned}
f(x)&=\frac2x\text{, }x\in\left[\frac12;4\right], \\ 
g(x)&=\frac{-3}x\text{, }x\in\left[\frac12;4\right], \\
h(x)&=\frac4x\text{, }x\in\left[\frac12;4\right]. 
\end{aligned}
\]
Choose the correct statement.}
{The function \( f \) is graphed in blue and the function \( h \) is graphed in green.}{The function \( g \) is graphed in red and the function \( f \) is graphed in green.}{The function \( f \) is graphed in green and the function \( h \) is graphed in blue.}{The function \( g \) is graphed in green and the function \( f \) is graphed in blue.}{}{}}
\LANGpl{
\Question{
Rysunek przedstawia wykresy funkcji
\[ \begin{aligned}
f(x)&=\frac2x\text{, }x\in\left\langle\frac12;4\right\rangle, \\ 
g(x)&=\frac{-3}x\text{, }x\in\left\langle\frac12;4\right\rangle, \\
h(x)&=\frac4x\text{, }x\in\left\langle\frac12;4\right\rangle. 
\end{aligned}
\]
Wybierz prawdziwe stwierdzenie.}
{Wykres funkcji \( f \) przedstawiono jest na niebiesko, a wykres funkcji \( h \)  na zielono.}{Wykres funkcji  \( g \) przedstawiono na czerwono, a  funkcji \( f \) na zielono.}{Wykres funkcji \( f \) przedstawiono na zielono, a funkcji \( h \) na niebiesko.}{Wykres funkcji \( g \) przedstawiono na zielono, a funkcji \( f \) na niebiesko.}{}{}}
\LANGcs{
\Question{
Na obrázku jsou grafy následujících funkcí:
\[ \begin{aligned}
f(x)&=\frac2x\text{, }x\in\left\langle\frac12;4\right\rangle, \\ 
g(x)&=\frac{-3}x\text{, }x\in\left\langle\frac12;4\right\rangle, \\
h(x)&=\frac4x\text{, }x\in\left\langle\frac12;4\right\rangle. 
\end{aligned}
\]
Vyberte pravdivý výrok.}
{Graf funkce \( f \) má modrou barvu a graf funkce \( h \) má zelenou barvu.}{Graf funkce \( g \) má červenou barvu a graf funkce \( f \) zelenou barvu.}{Graf funkce \( f \) má zelenou barvu a graf funkce \( h \) má modrou barvu.}{Graf funkce  \( g \) má zelenou barvu a graf funkce \( f \) má modrou barvu.}{}{}}
\LANGsk{
\Question{
Na obrázku sú grafy funkcií:
\[ \begin{aligned}
f(x)&=\frac2x\text{, }x\in\left\langle\frac12;4\right\rangle, \\ 
g(x)&=\frac{-3}x\text{, }x\in\left\langle\frac12;4\right\rangle, \\
h(x)&=\frac4x\text{, }x\in\left\langle\frac12;4\right\rangle. 
\end{aligned}
\]
Vyberte pravdivý výrok.}
{Graf funkcie \( f \) má modrú farbu a graf funkcie \( h \) má zelenú farbu.}{Graf funkcie \( g \) má červenú farbu a graf funkcie \( f \) má zelenú farbu.}{Graf funkcie \( f \) má zelenú farbu a graf funkcie \( h \) má modrú farbu.}{Graf funkcie \( g \) má zelenú farbu a graf funkcie \( f \) má modrú farbu.}{}{}}
\LANGes{
\Question{
La imagen muestra las gráficas de las funciones:
\[ \begin{aligned}
f(x)&=\frac2x\text{, }x\in\left[\frac12;4\right], \\ 
g(x)&=\frac{-3}x\text{, }x\in\left[\frac12;4\right], \\
h(x)&=\frac4x\text{, }x\in\left[\frac12;4\right]. 
\end{aligned}
\]
Elige la proposición verdadera.}
{La función \( f \) está representada en azul y la función \( h \) está representada en verde.}{La función \( g \) está representada en rojo y la función\( f \) está representada en verde.}{La función \( f \) está representada en verde y la función \( h \) está representada en azul.}{La función\( g \) está representada en verde y la función\( f \) está representada en azul.}{}{}}
\End

\Data\ID{2000003701}
\Updated{2021-12-26 08:12:28}
\Level{A}
\SubArea{Rational functions}
\LANGen{
\Question{
Group of mountaineers would climb the top of the mountain at an ascent speed of \(400\,\mathrm{m}\) per day in \(10\) days. However, due to the weather, they have to conquer the peak in \(8\) days. How many more meters per day do they have to cover?}
{\(100\) meters more}{\(80\) meters more}{\(120\) meters more}{\(90\) meters more}{}{}}
\LANGpl{
\Question{
Grupa alpinistów wspinała się na szczyt góry z prędkością \(400\,\mathrm{m}\) dziennie w ciagu  \(10\) dni. Jednak ze względu na pogodę muszą zdobyć szczyt w \(8\) dni. Ile więcej metrów dziennie muszą pokonać?}
{\(100\) metrów więcej}{\(80\) metrów więcej}{\(120\) metrów więcej}{\(90\) metrów więcej}{}{}}
\LANGcs{
\Question{
Skupina horolezců by zdolala vrcholek hory za \(10\) dní při rychlosti výstupu \(400\,\mathrm{m}\)  za den. Kvůli počasí však musí zdolat vrchol za \(8\) dní. O kolik metrů musí zdolat denně více?}
{o \(100\) metrů více}{o \(80\) metrů více}{o \(120\) metrů více}{o \(90\) metrů více}{}{}}
\LANGsk{
\Question{
Skupina horolezcov by zdolala vrchol hory za \(10\) dní pri rýchlosti výstupu \(400\,\mathrm{m}\)  za deň. Kvôli počasiu však musí zdolať vrchol za \(8\) dní. Koľko metrov viac musí denne zdolať?}
{\(100\) metrov viac}{\(80\) metrov viac}{\(120\) metrov viac}{\(90\) metrov viac}{}{}}
\LANGes{
\Question{
Un grupo de montañeros subiría la cima de una montaña en \(10\) días a una velocidad de ascenso de \(400\,\mathrm{m}\). Sin embargo, debido al mal tiempo, tienen que conquistar la cima en \(8\) días. ¿Cuántos metros más tienen que recorrer cada día?}
{\(100\) metros más}{\(80\) metros más}{\(120\) metros más}{\(90\) metros más}{}{}}
\End

\Data\ID{2000003702}
\Updated{2022-02-15 07:02:37}
\Level{A}
\SubArea{Rational functions}
\LANGen{
\Question{
Four workers assembled a garden pool in \(5\) hours. How long would it take eight workers to do the same job?}
{\(2\,\mathrm{h}\,30\,\mathrm{min}\)}{\(2\,\mathrm{h}\,40\,\mathrm{min}\)}{\(2\,\mathrm{h}\,20\,\mathrm{min}\)}{\(2\,\mathrm{h}\,45\,\mathrm{min}\)}{}{}}
\LANGpl{
\Question{
Czterech pracowników zmontowało basen ogrodowy w \(5\) godzin. Ile czasu zajęłoby ośmiu pracownikom wykonanie tej samej pracy?}
{\(2\,\mathrm{h}\,30\,\mathrm{min}\)}{\(2\,\mathrm{h}\,40\,\mathrm{min}\)}{\(2\,\mathrm{h}\,20\,\mathrm{min}\)}{\(2\,\mathrm{h}\,45\,\mathrm{min}\)}{}{}}
\LANGcs{
\Question{
Čtyři dělníci by smontovali zahradní bazén za  \(5\) hodin. Jak dlouho by stejná práce trvala osmi dělníkům?}
{\(2\,\mathrm{h}\,30\,\mathrm{min}\)}{\(2\,\mathrm{h}\,40\,\mathrm{min}\)}{\(2\,\mathrm{h}\,20\,\mathrm{min}\)}{\(2\,\mathrm{h}\,45\,\mathrm{min}\)}{}{}}
\LANGsk{
\Question{
Štyria robotníci by zmontovali záhradný bazén za  \(5\) hodín. Za ako dlho by rovnakú prácu spravili ôsmi robotníci?}
{\(2\,\mathrm{h}\,30\,\mathrm{min}\)}{\(2\,\mathrm{h}\,40\,\mathrm{min}\)}{\(2\,\mathrm{h}\,20\,\mathrm{min}\)}{\(2\,\mathrm{h}\,45\,\mathrm{min}\)}{}{}}
\LANGes{
\Question{
Cuatro trabajadores montaron una piscina de jardín en \(5\) horas. ¿Cuánto tardarían ocho trabajadores en hacer el mismo trabajo?}
{\(2\,\mathrm{h}\,30\,\mathrm{min}\)}{\(2\,\mathrm{h}\,40\,\mathrm{min}\)}{\(2\,\mathrm{h}\,20\,\mathrm{min}\)}{\(2\,\mathrm{h}\,45\,\mathrm{min}\)}{}{}}
\End

\Data\ID{2000003703}
\Updated{2022-02-20 08:02:39}
\Level{A}
\SubArea{Rational functions}
\LANGen{
\Question{
The pool is filled with eight equally powerful pumps in \(6\) hours. During the maintenance of some of the pumps, it took \(24\) hours to fill the pool. How many pumps were running?}
{\(2\) pumps}{\(3\) pumps}{\(4\) pumps}{\(6\) pumps}{}{}}
\LANGpl{
\Question{
Basen jest wypełniany ośmioma równie mocnymi pompami w ciągu \(6\) godzin. Podczas konserwacji niektórych pomp napełnianie basenu trwało \(24\) godziny. Ile pomp pracowało?}
{\(2\) pompy}{\(3\) pompy}{\(4\) pompy}{\(6\) pomp}{}{}}
\LANGcs{
\Question{
Bazén se napustí osmi stejně výkonnými čerpadly za  \(6\) hodin. Při údržbě některých čerpadel trvalo naplnění bazénu \(24\) hodin. Kolik čerpadel bylo v chodu?}
{\(2\) čerpadla}{\(3\) čerpadla}{\(4\) čerpadla}{\(6\) čerpadel}{}{}}
\LANGsk{
\Question{
Bazén sa ôsmimi rovnako výkonnými čerpadlami napustí  za  \(6\) hodín. Pri údržbe niektorých čerpadiel trvalo naplnenie bazéna \(24\) hodín. Koľko čerpadiel bolo v prevádzke?}
{\(2\) čerpadlá}{\(3\) čerpadlá}{\(4\) čerpadlá}{\(6\) čerpadiel}{}{}}
\LANGes{
\Question{
Una piscina se llena con ocho bombas de llenado igualmente potentes en \(6\) horas. Durante el mantenimiento de algunas de las bombas, la piscina se llenó en \(24\) horas. ¿Cuántas bombas estaban funcionando?}
{\(2\) bombas}{\(3\) bombas}{\(4\) bombas}{\(6\) bombas}{}{}}
\End

\Data\ID{2000003704}
\Updated{2022-02-15 07:02:37}
\Level{A}
\SubArea{Rational functions}
\LANGen{
\Question{
A car going at a speed \(60\,\mathrm{km/h}\)  covers the distance from city \(A\) to city B in \(30\) minutes. If the distance has to be covered in \(20\) minutes, by how many \(\mathrm{km/h}\) must the driver increase the speed when leaving \(A\).}
{by \(30\,\mathrm{km/h}\)}{by \(20\,\mathrm{km/h}\)}{by \(40\,\mathrm{km/h}\)}{by \(45\,\mathrm{km/h}\)}{}{}}
\LANGpl{
\Question{
Samochód jadący z prędkością  \(60\,\mathrm{km/h}\)  pokonuje odległość z miasta \(A\) do miasta B w ciągu \(30\) minut. Jeśli odległość ma zostać pokonana w ciągu \(20\) minut, o ile \(\mathrm{km/h}\) kierowca musi zwiększyć prędkość opuszczając miasto \(A\).}
{o \(30\,\mathrm{km/h}\)}{o \(20\,\mathrm{km/h}\)}{o \(40\,\mathrm{km/h}\)}{o \(45\,\mathrm{km/h}\)}{}{}}
\LANGcs{
\Question{
Automobil jedoucí rychlostí  \(60\,\mathrm{km/h}\)  urazí obvykle vzdálenost z města  \(A\) do města \(B\)  za \(30\) minut.  O kolik \(\mathrm{km/h}\) musí řidič zvýšit rychlost, má-li být po odjezdu z \(A\) ve městě \(B\) už za \(20\) minut?}
{o \(30\,\mathrm{km/h}\)}{o \(20\,\mathrm{km/h}\)}{o \(40\,\mathrm{km/h}\)}{o \(45\,\mathrm{km/h}\)}{}{}}
\LANGsk{
\Question{
Auto idúce rýchlosťou  \(60\,\mathrm{km/h}\)  zvyčajne prejde vzdialenosť z mesta  \(A\) do mesta \(B\)  za \(30\) minút.  O koľko \(\mathrm{km/h}\) musí vodič zvýšiť rýchlosť, ak má byť po odjazde z \(A\) v meste \(B\) už za \(20\) minút?}
{o \(30\,\mathrm{km/h}\)}{o \(20\,\mathrm{km/h}\)}{o \(40\,\mathrm{km/h}\)}{o \(45\,\mathrm{km/h}\)}{}{}}
\LANGes{
\Question{
Un coche que va a una velocidad de \(60\,\mathrm{km/h}\) recorre la distancia de la ciudad \(A\) a la ciudad B en \(30\) minutos. Si la distancia tiene que ser recorrida en \(20\) minutos, ¿en cuántos \(\mathrm{km/h}\) el conductor debe aumentar la velocidad al salir de \(A\)?}
{en \(30\,\mathrm{km/h}\)}{en \(20\,\mathrm{km/h}\)}{en \(40\,\mathrm{km/h}\)}{en \(45\,\mathrm{km/h}\)}{}{}}
\End

\Data\ID{2000003705}
\Updated{2022-02-15 07:02:37}
\Level{A}
\SubArea{Rational functions}
\LANGen{
\Question{
A car going at a speed \(60\,\mathrm{km/h}\)  covers the distance from city \(A\) to city \(B\) in \(30\) minutes. If the distance has to be covered in \(20\) minutes, how many times does the driver have to increase the speed when leaving \(A\).}
{\(1.5\) times}{\(1.\overline{3}\) times}{\(1.\overline{6}\) times}{\(1.2\) times}{}{}}
\LANGpl{
\Question{
Samochód jadący z prędkością \(60\,\mathrm{km/h}\)  pokonuje odległość z miasta \(A\) do miasta \(B\) w ciągu \(30\) minut. Jeśli odległość ma zostać pokonana w ciągu \(20\) minut, ile razy kierowca musi zwiększać prędkość opuszczając miasto \(A\).}
{\(1{,}5\) razy}{\(1{,}\overline{3}\) razy}{\(1{,}\overline{6}\) razy}{\(1{,}2\) razy}{}{}}
\LANGcs{
\Question{
Automobil jedoucí rychlostí \(60\,\mathrm{km/h}\) urazí vzdálenost z města \(A\) do města \(B\) za \(30\) minut. Kolikrát musí řidič zvýšit rychlost, má-li být v \(B\) za \(20\)  minut od výjezdu z \(A\)?}
{\(1{,}5\) krát}{\(1{,}\overline{3}\) krát}{\(1{,}\overline{6}\) krát}{\(1{,}2\) krát}{}{}}
\LANGsk{
\Question{
Auto idúce rýchlosťou \(60\,\mathrm{km/h}\) prejde vzdialenosť z mesta \(A\) do mesta \(B\) za \(30\) minút. Koľkokrát musí vodič zvýšiť rýchlosť, ak má byť v \(B\) za \(20\)  minút od odchodu z mesta  \(A\)?}
{\(1{,}5\) krát}{\(1{,}\overline{3}\) krát}{\(1{,}\overline{6}\) krát}{\(1{,}2\) krát}{}{}}
\LANGes{
\Question{
Un coche que va a una velocidad de \(60\,\mathrm{km/h}\)  recorre la distancia de la ciudad \(A\) a la ciudad \(B\) en \(30\) minutos. Si la distancia tiene que ser recorrida en \(20\) minutos, ¿cuántas veces tiene que aumentar el conductor  la velocidad al salir de \(A\)?}
{\(1.5\) veces}{\(1.\overline{3}\) veces}{\(1.\overline{6}\) veces}{\(1.2\) veces}{}{}}
\End

\Data\ID{2000003706}
\Updated{2022-02-20 09:02:41}
\Level{A}
\SubArea{Rational functions}
\LANGen{
\Question{
The length of a rectangle is extended twice its original length. How must its width be changed so that the area of the rectangle remains the same?}
{the width is reduced to half (of its original width)}{the width is increased by half (of its original width)}{the width is reduced by a quarter (of its original width)}{the width is increased to double  (of its original width)}{}{}}
\LANGpl{
\Question{
Długość prostokąta jest dwukrotnie wydłużana w stosunku do pierwotnej długości. Jak zmienić jego szerokość, aby powierzchnia prostokąta pozostała taka sama?}
{szerokość musi zostać zmniejszona do połowy (pierwotnej szerokości)}{szerokość musi zostać zwiększona o połowę (pierwotnej szerokości)}{szerokość musi zostać zmniejszona o jedną czwartą (pierwotnej szerokości)}{szerokość musi zostać zwiększona dwukrotnie (w stosunku do oryginalnej szerokości)}{}{}}
\LANGcs{
\Question{
Prodlouží-li se délka obdélníka dvakrát, jak se musí změnit jeho šířka, má-li mít stále stejný obsah?}
{šířka bude poloviční než původní šířka}{šířka se zvětší o polovinu původní šířky}{šířka se zmenší o čtvrtinu původní šířky}{šířka bude dvakrát větší než původní šířka}{}{}}
\LANGsk{
\Question{
Ak sa predĺži dĺžka obdĺžnika dvakrát, ako sa musí zmeniť jeho šírka, ak má mať stále rovnaký obsah?}
{šírka bude polovičná ako pôvodná šírka}{šírka sa zväčší o polovicu pôvodnej šírky}{šírka sa zmenší o štvrtinu pôvodnej šírky}{šírka bude dvakrát väčšia než pôvodná šírka}{}{}}
\LANGes{
\Question{
Aumentamos la longitud de un rectángulo el doble de su longitud original. ¿Cómo  debe cambiar su ancho para que el área del rectángulo permanezca igual?}
{el ancho se reduce a la mitad (de su ancho original)}{el ancho se incrementa a la mitad (de su ancho original)}{el ancho se reduce en un cuarto (de su ancho original)}{el ancho se incrementa al doble (de su ancho original)}{}{}}
\End

\Data\ID{2000018801}
\Updated{2023-01-23 10:01:39}
\Level{A}
\SubArea{Rational functions}
\LANGen{
\Question{
Consider a triangle with area of 
\(5\, \mathrm{cm}^{2}\). Find the formula which relates the length of its side \(a\) to the length of the height \(v_a\)  , where \(v_a\) is the height to the side \(a\).}
{\(v_a = \frac{10} 
{a}\)}{\(v_a = \frac{5} 
{a}\)}{\(v_a =5
{a}\)}{\(v_a = \frac{5} 
{2a}\)}{}{}}
\LANGpl{
\Question{
Dany jest trójkąt o powierzchni 
\(5\, \mathrm{cm}^{2}\).  Znajdź wzór, który wiąże długość jego boku \(a\)  z długością wysokości \(v_a\), gdzie \(v_a\) to wysokość do boku \(a\).}
{\(v_a = \frac{10} 
{a}\)}{\(v_a = \frac{5} 
{a}\)}{\(v_a =5
{a}\)}{\(v_a = \frac{5} 
{2a}\)}{}{}}
\LANGcs{
\Question{
Obsah trojúhelníku je  
\(5\, \mathrm{cm}^{2}\). Určete předpis funkce, která vyjadřuje závislost mezi velikostí jeho strany  \(a\) a velikostí výšky \(v_a\)  na tuto stranu.}
{\(v_a = \frac{10} 
{a}\)}{\(v_a = \frac{5} 
{a}\)}{\(v_a =5
{a}\)}{\(v_a = \frac{5} 
{2a}\)}{}{}}
\LANGsk{
\Question{
Obsah trojuholníka je  
\(5\, \mathrm{cm}^{2}\). Určte funkciu, ktorá vyjadruje závislosť medzi veľkosťou jeho strany  \(a\) a veľkosťou výšky \(v_a\)  na túto stranu.}
{\(v_a = \frac{10} 
{a}\)}{\(v_a = \frac{5} 
{a}\)}{\(v_a =5
{a}\)}{\(v_a = \frac{5} 
{2a}\)}{}{}}
\LANGes{
\Question{
Dado un triángulo de área 
\(5\, \mathrm{cm}^{2}\). Halla la fórmula que relaciona la longitud de su lado \(a\) con la longitud de la altura \(v_a\)  , donde \(v_a\) es la altura respecto del lado \(a\).}
{\(v_a = \frac{10} 
{a}\)}{\(v_a = \frac{5} 
{a}\)}{\(v_a =5
{a}\)}{\(v_a = \frac{5} 
{2a}\)}{}{}}
\End

\Data\ID{2000018802}
\Updated{2022-11-02 09:11:09}
\Level{A}
\SubArea{Rational functions}
\LANGen{
\Question{
Given the function \(f(x)= \frac{6}
 {x} \),
evaluate \( \frac{f(-3)}{ f(2)}\).}
{\(-\frac23\)}{\(-6\)}{\(-\frac32\)}{\(-\frac16\)}{}{}}
\LANGpl{
\Question{
Dana jest funkcja \(f(x)= \frac{6}
 {x} \).
Oblicz \( \frac{f(-3)}{ f(2)}\).}
{\(-\frac23\)}{\(-6\)}{\(-\frac32\)}{\(-\frac16\)}{}{}}
\LANGcs{
\Question{
Je dána funkce  \(f(x)= \frac{6}
 {x} \).
Vyčíslete \( \frac{f(-3)}{ f(2)}\).}
{\(-\frac23\)}{\(-6\)}{\(-\frac32\)}{\(-\frac16\)}{}{}}
\LANGsk{
\Question{
Je daná funkcia \(f(x)= \frac{6}
 {x} \).
Vypočítajte \( \frac{f(-3)}{ f(2)}\).}
{\(-\frac23\)}{\(-6\)}{\(-\frac32\)}{\(-\frac16\)}{}{}}
\LANGes{
\Question{
Dada la función \(f(x)= \frac{6}
 {x} \),
evalúa \( \frac{f(-3)}{ f(2)}\).}
{\(-\frac23\)}{\(-6\)}{\(-\frac32\)}{\(-\frac16\)}{}{}}
\End

\Data\ID{2000018803}
\Updated{2022-11-02 09:11:09}
\Level{A}
\SubArea{Rational functions}
\LANGen{
\Question{
Given the function \(f(x) = \frac{5} 
{x}\),
find the function \(g\) such
that the graphs of \(f\) 
and \(g\) are symmetric
about the line \(y = -x\).}
{\(g(x) = \frac{5} 
{x}\)}{\(g(x) =5
{x}\)}{\(g(x) = -\frac{5} 
{x}\)}{\(g(x) = -{5} 
{x}\)}{}{}}
\LANGpl{
\Question{
Mając funkcję \(f(x) = \frac{5} 
{x}\),
znajdź funkcję \(g\),  tak aby wykresy funkcji \(f\) 
i \(g\) były symetryczne względem linii \(y = -x\).}
{\(g(x) = \frac{5} 
{x}\)}{\(g(x) =5
{x}\)}{\(g(x) = -\frac{5} 
{x}\)}{\(g(x) = -{5} 
{x}\)}{}{}}
\LANGcs{
\Question{
Je dána funkce \(f(x) = \frac{5} 
{x}\).
Vyberte předpis funkce \(g\) tak,
aby grafy funkcí \(f\) 
a \(g\) byly osově souměrné
podle přímky \(y = -x\).}
{\(g(x) = \frac{5} 
{x}\)}{\(g(x) =5
{x}\)}{\(g(x) = -\frac{5} 
{x}\)}{\(g(x) = -{5} 
{x}\)}{}{}}
\LANGsk{
\Question{
Je daná funkcia \(f(x) = \frac{5} 
{x}\).
Vyberte predpis funkcie \(g\) tak,
aby grafy funkcií \(f\) 
a \(g\) boli osovo súmerné
podľa priamky \(y = -x\).}
{\(g(x) = \frac{5} 
{x}\)}{\(g(x) =5
{x}\)}{\(g(x) = -\frac{5} 
{x}\)}{\(g(x) = -{5} 
{x}\)}{}{}}
\LANGes{
\Question{
Dada la función \(f(x) = \frac{5} 
{x}\),
halla la función \(g\) tal que las gráficas de \(f\) 
y \(g\) sean simétricas respecto de la recta \(y = -x\).}
{\(g(x) = \frac{5} 
{x}\)}{\(g(x) =5
{x}\)}{\(g(x) = -\frac{5} 
{x}\)}{\(g(x) = -{5} 
{x}\)}{}{}}
\End

\Data\ID{2000018804}
\Updated{2022-11-02 09:11:13}
\Level{A}
\SubArea{Rational functions}
\LANGen{
\Question{
Given the function \(f(x) = -\frac{4} 
{x}\),
find the function \(g\) such
that the graphs of \(f\) 
and \(g\) are symmetric
about the \(y\)-axis.}
{\(g(x) = \frac{4} 
{x}\)}{\(g(x) =4
{x}\)}{\(g(x) = -\frac{4} 
{x}\)}{\(g(x) = -{4} 
{x}\)}{}{}}
\LANGpl{
\Question{
Dana jest funkcja \(f(x) = -\frac{4} 
{x}\).
Znajdź funkcje \(g\)  taką, że wykresy funkcji \(f\) 
i \(g\) są symetryczne względem osi \(y\).}
{\(g(x) = \frac{4} 
{x}\)}{\(g(x) =4
{x}\)}{\(g(x) = -\frac{4} 
{x}\)}{\(g(x) = -{4} 
{x}\)}{}{}}
\LANGcs{
\Question{
Je dána funkce \(f(x) = -\frac{4} 
{x}\).
Vyberte předpis funkce \(g\), jejíž
graf je souměrný s grafem funkce  \(f\) 
podle osy \(y\).}
{\(g(x) = \frac{4} 
{x}\)}{\(g(x) =4
{x}\)}{\(g(x) = -\frac{4} 
{x}\)}{\(g(x) = -{4} 
{x}\)}{}{}}
\LANGsk{
\Question{
Je daná funkcia \(f(x) = -\frac{4} 
{x}\).
Vyberte predpis funkcie \(g\), ktorej
graf je súmerný s grafom funkcie  \(f\) 
podľa osy \(y\).}
{\(g(x) = \frac{4} 
{x}\)}{\(g(x) =4
{x}\)}{\(g(x) = -\frac{4} 
{x}\)}{\(g(x) = -{4} 
{x}\)}{}{}}
\LANGes{
\Question{
Dada la función \(f(x) = -\frac{4} 
{x}\),
halla la función \(g\) tal que las gráficas de \(f\) 
y \(g\) sean simétricas respecto del eje \(y\).}
{\(g(x) = \frac{4} 
{x}\)}{\(g(x) =4
{x}\)}{\(g(x) = -\frac{4} 
{x}\)}{\(g(x) = -{4} 
{x}\)}{}{}}
\End

\Data\ID{2000018805}
\Updated{2022-11-02 09:11:13}
\Level{A}
\SubArea{Rational functions}
\LANGen{
\Question{
The test driver drove from Ostrava to Warsaw at an average speed of \(66\, \mathrm{km}/\mathrm{h}\) and the journey took him \(6\) hours. After him, the same route took several other drivers. (Each driver took a different driving time.) Choose the function giving the average speed \(v\) of each of these drivers as a function of the total driving time \(t\) from Ostrava to Warsaw.}
{\( v=\frac{396}t,\ \ t\in(0;\infty) \)}{\( v=\frac{66}t,\ \ t\in(0;\infty) \)}{\( v=66 t,\ \ t\in(0;\infty) \)}{\( v=\frac{t}{396},\ \ t\in(0;\infty) \)}{}{}}
\LANGpl{
\Question{
Kierowca testowy jechał z Ostrawy do Warszawy ze średnią prędkością \(66\, \mathrm{km}/\mathrm{h}\), a podróż zajęła mu \(6\) godzin. Po nim tą samą trasą jechało kilku innych kierowców. (Każdy kierowca miał inny czas jazdy.)  Wybierz funkcję podającą średnią prędkość \(v\) każdego z tych kierowców  jako funkcję łącznego czasu jazdy \(t\) z Ostrawy do Warszawy.}
{\( v=\frac{396}t,\ \ t\in(0;\infty) \)}{\( v=\frac{66}t,\ \ t\in(0;\infty) \)}{\( v=66 t,\ \ t\in(0;\infty) \)}{\( v=\frac{t}{396},\ \ t\in(0;\infty) \)}{}{}}
\LANGcs{
\Question{
Testovací jezdec jel z Ostravy do Varšavy průměrnou rychlostí  \(66\, \mathrm{km}/\mathrm{h}\) a cesta mu trvala  \(6\) hodin. Po něm stejnou trasu jelo ještě několik dalších vozidel (každému trvala cesta jinak dlouho). Vyberte funkci, která udává průměrnou rychlost \(v\) těchto vozidel v závislosti na celkovém čase \(t\) při jízdě z Ostravy do Varšavy.}
{\( v=\frac{396}t,\ \ t\in(0;\infty) \)}{\( v=\frac{66}t,\ \ t\in(0;\infty) \)}{\( v=66 t,\ \ t\in(0;\infty) \)}{\( v=\frac{t}{396},\ \ t\in(0;\infty) \)}{}{}}
\LANGsk{
\Question{
Testovací jazdec šiel z Ostravy do Varšavy priemernou rýchlosťou  \(66\, \mathrm{km}/\mathrm{h}\) a cesta mu trvala  \(6\) hodín. Po ňom rovnakú trasu išlo ešte niekoľko ďalších vozidiel (každému trvala cesta inak dlho). Vyberte funkciu, ktorá zobrazuje priemernú rýchlosť \(v\) týchto vozidiel v závislosti na celkovom čase \(t\) pri jazde z Ostravy do Varšavy.}
{\( v=\frac{396}t,\ \ t\in(0;\infty) \)}{\( v=\frac{66}t,\ \ t\in(0;\infty) \)}{\( v=66 t,\ \ t\in(0;\infty) \)}{\( v=\frac{t}{396},\ \ t\in(0;\infty) \)}{}{}}
\LANGes{
\Question{
Un piloto de prueba condujo desde Ostrava hasta Varsovia a una velocidad media de \(66\, \mathrm{km}/\mathrm{h}\) y el viaje duró \(6\) horas. Después de él, varios conductores tomaron la misma ruta. (Cada conductor tardó un tiempo diferente.) Elige la función que da la velocidad media \(v\) de cada uno de estos conductores en función del tiempo total de conducción \(t\) de Ostrava a Varsovia.}
{\( v=\frac{396}t,\ \ t\in(0;\infty) \)}{\( v=\frac{66}t,\ \ t\in(0;\infty) \)}{\( v=66 t,\ \ t\in(0;\infty) \)}{\( v=\frac{t}{396},\ \ t\in(0;\infty) \)}{}{}}
\End

\Data\ID{2010009905}
\Updated{2022-08-25 10:08:41}
\Level{A}
\SubArea{Rational functions}
\LANGen{
\Question{
Let \( f(x)=\frac{-3}{x} \). Find the false statement.}
{The function \(f\) is bounded above.}{The range of \( f \) is \( \left(-\infty;0\right)\cup\left(0;\infty\right) \).}{The function \( f \) is increasing on \( \left(-\infty;0\right) \).}{The function \( h \) defined by \(h(x)=-f(x)\) is an odd function.}{}{}}
\LANGpl{
\Question{
Dana jest funkcja\( f(x)=\frac{-3}{x} \). Wskaż błędne stwierdzenie.}
{Funkcja \(f\) jest ograniczona z góry.}{Zbiór wartości funkcji \( f \) to \( \left(-\infty;0\right)\cup\left(0;\infty\right) \).}{Funkcja \( f \) rośnie w \( \left(-\infty;0\right) \).}{Funkcja \( h \) określona jako \(h(x)=-f(x)\) jest funkcją nieparzystą.}{}{}}
\LANGcs{
\Question{
Je dána funkce \( f(x)=\frac{-3}{x} \). Vyberte nepravdivý výrok.}
{Funkce \(f\) je shora omezená.}{Oborem hodnot funkce \( f \) je množina \( \left(-\infty;0\right)\cup\left(0;\infty\right) \).}{Funkce \( f \) je rostoucí na intervalu \( \left(-\infty;0\right) \).}{Funkce \( h \) definovaná vztahem \(h(x)=-f(x)\) je lichá.}{}{}}
\LANGsk{
\Question{
Je daná funkcia \( f(x)=\frac{-3}{x} \). Vyberte nepravdivý výrok.}
{Funkcia \(f\) je ohraničená zhora.}{Oborom hodnôt funkcie \( f \) je \( \left(-\infty;0\right)\cup\left(0;\infty\right) \).}{Funkcia  \( f \) rastie na \( \left(-\infty;0\right) \).}{Funkcia \( h \) definovaná vzťahom \(h(x)=-f(x)\) je nepárna funkcia.}{}{}}
\LANGes{
\Question{
Sea \( f(x)=\frac{-3}{x} \). Determina la proposición falsa.}
{La función \(f\) está acotada superiormente.}{El rango de la función  \( f \) es \( \left(-\infty;0\right)\cup\left(0;\infty\right) \).}{La función \( f \) es creciente en \( \left(-\infty;0\right) \).}{La función \( h \) definida por \(h(x)=-f(x)\) es impar.}{}{}}
\End

\Data\ID{2100006702}
\Updated{2023-01-23 10:01:12}
\Level{A}
\SubArea{Rational functions}
\LANGen{
\Question{
Identify which of the given graphs is the graph of the function \(f(x)=\frac{-1}{3x}\).}
{}{}{}{}{}{}}
\LANGpl{
\Question{
Wskaż wykres, który jest wykresem funkcji  \(f(x)=\frac{-1}{3x}\).}
{}{}{}{}{}{}}
\LANGcs{
\Question{
Na kterém obrázku je graf funkce \(f(x)=\frac{-1}{3x}\)?}
{}{}{}{}{}{}}
\LANGsk{
\Question{
Ktorý z grafov, je grafom funkcie \(f(x)=\frac{-1}{3x}\)?}
{}{}{}{}{}{}}
\LANGes{
\Question{
Identifica una posible gráfica de la función \(f(x)=\frac{-1}{3x}\).}
{}{}{}{}{}{}}

\IMGanswer{2000006701_6-figure1.pdf}
\IMGanswer{2000006701_6-figure2.pdf}
\IMGanswer{2000006701_6-figure3.pdf}
\IMGanswer{2000006701_6-figure4.pdf}\End

\Data\ID{2100006703}
\Updated{2023-01-23 10:01:12}
\Level{A}
\SubArea{Rational functions}
\LANGen{
\Question{
Identify which of the given graphs is the graph of the function \(f(x)=\frac{2}{5x}\).}
{}{}{}{}{}{}}
\LANGpl{
\Question{
Wybierz wykres, który jest wykresem funkcji \(f(x)=\frac{2}{5x}\).}
{}{}{}{}{}{}}
\LANGcs{
\Question{
Určete, na kterém obrázku je graf funkce \(f(x)=\frac{2}{5x}\).}
{}{}{}{}{}{}}
\LANGsk{
\Question{
Zistite, ktorý z daných grafov je grafom funkcie \(f(x)=\frac{2}{5x}\).}
{}{}{}{}{}{}}
\LANGes{
\Question{
Identifica una posible gráfica de la función \(f(x)=\frac{2}{5x}\).}
{}{}{}{}{}{}}

\IMGanswer{2000006701_6-figure5.pdf}
\IMGanswer{2000006701_6-figure6.pdf}
\IMGanswer{2000006701_6-figure7.pdf}
\IMGanswer{2000006701_6-figure8.pdf}\End

\Data\ID{9000002910}
\Updated{2023-09-10 11:09:22}
\Level{A}
\SubArea{Rational functions}
\LANGen{
\Question{
Consider a rectangle with area of 
\(5\, \mathrm{cm}^{2}\). Find the formula which relates side \(a\) to the side \(b\)  of this rectangle.}
{\(b = \frac{5} 
{a}\),
\(a\in (0;\infty )\)}{\(b = 5a\),
\(a\in (0;\infty )\)}{\(b = \frac{10} 
 {a} \),
\(a\in (0;\infty )\)}{\(b = \frac{25} 
 {a} \),
\(a\in (0;\infty )\)}{}{}}
\LANGpl{
\Question{
Rozważ prostokąt o stałym polu powierzchni
\(5\, \mathrm{cm}^{2}\). Wyznacz wzór odnoszący się do obydwu boków tego prostokąta.}
{\(b = \frac{5} 
{a}\),
\(a\in (0;\infty )\)}{\(b = 5a\),
\(a\in (0;\infty )\)}{\(b = \frac{10} 
 {a} \),
\(a\in (0;\infty )\)}{\(b = \frac{25} 
 {a} \),
\(a\in (0;\infty )\)}{}{}}
\LANGcs{
\Question{
Obsah obdélníku je \(5\, \mathrm{cm}^{2}\).
Zapište funkci, která vyjadřuje závislost mezi velikostmi jeho stran.}
{\(b = \frac{5} 
{a}\),
\(a\in (0;\infty )\)}{\(b = 5a\),
\(a\in (0;\infty )\)}{\(b = \frac{10} 
 {a} \),
\(a\in (0;\infty )\)}{\(b = \frac{25} 
 {a} \),
\(a\in (0;\infty )\)}{}{}}
\LANGsk{
\Question{
Obsah obdĺžnika je \(5\, \mathrm{cm}^{2}\).
Určte funkciu, ktorá vyjadruje závislosť medzi veľkosťami jeho strán.}
{\(b = \frac{5} 
{a}\),
\(a\in (0;\infty )\)}{\(b = 5a\),
\(a\in (0;\infty )\)}{\(b = \frac{10} 
 {a} \),
\(a\in (0;\infty )\)}{\(b = \frac{25} 
 {a} \),
\(a\in (0;\infty )\)}{}{}}
\LANGes{
\Question{
Considera un rectángulo cuya área es de
\(5\, \mathrm{cm}^{2}\).  Determina la fórmula que relaciona ambos lados del rectángulo.}
{\(b = \frac{5} 
{a}\),
\(a\in (0;\infty )\)}{\(b = 5a\),
\(a\in (0;\infty )\)}{\(b = \frac{10} 
 {a} \),
\(a\in (0;\infty )\)}{\(b = \frac{25} 
 {a} \),
\(a\in (0;\infty )\)}{}{}}
\End

\Data\ID{9000003101}
\Updated{2021-04-09 07:04:21}
\Level{A}
\SubArea{Rational functions}
\IMG{000031_01-figure0.pdf}
\LANGen{
\Question{
Identify a possible analytic expression for the function graphed in the picture.}
{\(y = \frac{1} 
{2x}\)}{\(y = \frac{2} 
{x}\)}{\(y = -\frac{2}
{x}\)}{\(y = -\frac{1}
{2x}\)}{}{}}
\LANGpl{
\Question{
Określ prawdopodobny wzór funkcji, której wykres przedstawiony jest na rysunku poniżej.}
{\(y = \frac{1} 
{2x}\)}{\(y = \frac{2} 
{x}\)}{\(y = -\frac{2}
{x}\)}{\(y = -\frac{1}
{2x}\)}{}{}}
\LANGcs{
\Question{
K danému grafu funkce přiřaďte správný funkční předpis.}
{\(y = \frac{1} 
{2x}\)}{\(y = \frac{2} 
{x}\)}{\(y = -\frac{2}
{x}\)}{\(y = -\frac{1}
{2x}\)}{}{}}
\LANGsk{
\Question{
Predpis funkcie, ktorej graf vidíte na obrázku, je:}
{\(y = \frac{1} 
{2x}\)}{\(y = \frac{2} 
{x}\)}{\(y = -\frac{2}
{x}\)}{\(y = -\frac{1}
{2x}\)}{}{}}
\LANGes{
\Question{
Identifica una posible expresión analítica para la gráfica de la función de la imagen.}
{\(y = \frac{1} 
{2x}\)}{\(y = \frac{2} 
{x}\)}{\(y = -\frac{2}
{x}\)}{\(y = -\frac{1}
{2x}\)}{}{}}
\End

\Data\ID{9000003102}
\Updated{2021-04-09 07:04:41}
\Level{A}
\SubArea{Rational functions}
\IMG{000031_02-figure0.pdf}
\LANGen{
\Question{
Identify a possible analytic expression for the function graphed in the picture.}
{\(y = -\frac{2}
{x}\)}{\(y = \frac{2} 
{x}\)}{\(y = \frac{1} 
{2x}\)}{\(y = -\frac{1}
{2x}\)}{}{}}
\LANGpl{
\Question{
Określ prawdopodobny wzór funkcji, której wykres przedstawiony jest na rysunku poniżej.}
{\(y = -\frac{2}
{x}\)}{\(y = \frac{2} 
{x}\)}{\(y = \frac{1} 
{2x}\)}{\(y = -\frac{1}
{2x}\)}{}{}}
\LANGcs{
\Question{
K danému grafu funkce přiřaďte správný funkční předpis.}
{\(y = -\frac{2}
{x}\)}{\(y = \frac{2} 
{x}\)}{\(y = \frac{1} 
{2x}\)}{\(y = -\frac{1}
{2x}\)}{}{}}
\LANGsk{
\Question{
Predpis funkcie, ktorej graf vidíte na obrázku, je:}
{\(y = -\frac{2}
{x}\)}{\(y = \frac{2} 
{x}\)}{\(y = \frac{1} 
{2x}\)}{\(y = -\frac{1}
{2x}\)}{}{}}
\LANGes{
\Question{
Identifica una posible expresión analítica para la gráfica de la función de la imagen.}
{\(y = -\frac{2}
{x}\)}{\(y = \frac{2} 
{x}\)}{\(y = \frac{1} 
{2x}\)}{\(y = -\frac{1}
{2x}\)}{}{}}
\End

\Data\ID{9000008001}
\Updated{2021-04-09 07:04:18}
\Level{A}
\SubArea{Rational functions}
\LANGen{
\Question{
Consider the function \(f(x)= -\frac{4}
{x}\) 
and the points \(A = [1;-4]\),
\(B = [-2;2]\),
\(C = [4;1]\),
\(D = [2;2]\).
How many of these points are on the graph of the function
\(f\)?}
{\(2\)}{\(1\)}{\(3\)}{\(4\)}{}{}}
\LANGpl{
\Question{
Rozważ funkcję \(f\colon y = -\frac{4}
{x}\) 
i punkty  \(A = [1;-4]\),
\(B = [-2;2]\),
\(C = [4;1]\),
\(D = [2;2]\).
Ile z tych punktów należy do wykresu funkcji 
\(f\)?}
{\(2\)}{\(1\)}{\(3\)}{\(4\)}{}{}}
\LANGcs{
\Question{
Je dána funkce \(f\colon y = -\frac{4}
{x}\) 
a body \(A = [1;-4]\),
\(B = [-2;2]\),
\(C = [4;1]\),
\(D = [2;2]\).
Kolik z uvedených bodů leží na grafu funkce
\(f\)?}
{\(2\)}{\(1\)}{\(3\)}{\(4\)}{}{}}
\LANGsk{
\Question{
Daná je funkcia \(f\colon y = -\frac{4}
{x}\) 
a body \(A = [1;-4]\),
\(B = [-2;2]\),
\(C = [4;1]\),
\(D = [2;2]\).
Koľko z daných bodov leží na grafe funkcie
\(f\)?}
{\(2\)}{\(1\)}{\(3\)}{\(4\)}{}{}}
\LANGes{
\Question{
Considera la función \(f(x)= -\frac{4}
{x}\) 
y los puntos \(A = [1;-4]\),
\(B = [-2;2]\),
\(C = [4;1]\),
\(D = [2;2]\).
¿Cuántos de estos puntos pertenecen a la gráfica de la función?
\(f\)?}
{\(2\)}{\(1\)}{\(3\)}{\(4\)}{}{}}
\End

\Data\ID{9000008002}
\Updated{2020-07-15 03:07:34}
\Level{A}
\SubArea{Rational functions}
\LANGen{
\Question{
Consider the point \(A = [-1;-3]\) and
the function \(f(x) = \frac{k} 
{x}\) with a nonzero
real parameter \(k\in \mathbb{R}\setminus \{0\}\). Identify
the value of the parameter \(k\) 
which ensures that the point \(A\) 
is on the graph of \(f\).}
{\(3\)}{\(1\)}{\(- 1\)}{\(- 3\)}{}{}}
\LANGpl{
\Question{
Rozważ punkt  \(A = [-1;-3]\) i
funkcję  \(f\colon y = \frac{k} 
{x}\) z niezerowym rzeczywistym parametrem
 \(k\in \mathbb{R}\setminus \{0\}\).  Wyznacz wartość tego parametru \(k\), który gwarantuje, że punkt  \(A\) 
należy do wykresu funkcji \(f\).}
{\(3\)}{\(1\)}{\(- 1\)}{\(- 3\)}{}{}}
\LANGcs{
\Question{
Je dána funkce \(f\colon y = \frac{k} 
{x}\) 
a bod \(A = [-1;-3]\). Pro
které \(k\in \mathbb{R}\setminus \{0\}\) graf funkce
\(f\) prochází
bodem \(A\)?}
{\(3\)}{\(1\)}{\(- 1\)}{\(- 3\)}{}{}}
\LANGsk{
\Question{
Daná je funkcia \(f\colon y = \frac{k} 
{x}\) 
a bod \(A = [-1;-3]\). Pre
ktoré \(k\in \mathbb{R}\setminus \{0\}\) graf funkcie
\(f\) prechádza
bodom \(A\)?}
{\(3\)}{\(1\)}{\(- 1\)}{\(- 3\)}{}{}}
\LANGes{
\Question{
Considera el punto \(A = [-1;-3]\) y la función \(f(x) = \frac{k} 
{x}\) con un parámetro real distinto de cero \(k\in \mathbb{R}\setminus \{0\}\). Identifica
el valor del parámetro \(k\) 
que asegura que el punto \(A\) 
está en la gráfica de la función \(f\).}
{\(3\)}{\(1\)}{\(- 1\)}{\(- 3\)}{}{}}
\End

\Data\ID{9000008003}
\Updated{2020-07-15 03:07:34}
\Level{A}
\SubArea{Rational functions}
\LANGen{
\Question{
Given the function \(f(x) = \frac{6} 
{x}\),
solve the following equation. 
\[ 
 f(x) = 2
\]}
{\(3\)}{\(2\)}{\(- 2\)}{\(- 3\)}{}{}}
\LANGpl{
\Question{
Dana jest funkcja \(f\colon y = \frac{6} 
{x}\). Rozwiąż następujące równanie.
\[ 
 f(x) = 2
\]}
{\(3\)}{\(2\)}{\(- 2\)}{\(- 3\)}{}{}}
\LANGcs{
\Question{
Je dána funkce \(f\colon y = \frac{6} 
{x}\).
Pro které \(x\in D(f)\) 
nabývá funkce \(f\) 
hodnoty \(2\)?}
{\(3\)}{\(2\)}{\(- 2\)}{\(- 3\)}{}{}}
\LANGsk{
\Question{
Daná je funkcia \(f\colon y = \frac{6} 
{x}\).
Pre ktoré \(x\in D(f)\) 
nadobúda funkcia \(f\) 
hodnotu \(2\)?}
{\(3\)}{\(2\)}{\(- 2\)}{\(- 3\)}{}{}}
\LANGes{
\Question{
Dada la función \(f(x) = \frac{6} 
{x}\),
resuelve la siguiente ecuación.
\[ 
 f(x) = 2
\]}
{\(3\)}{\(2\)}{\(- 2\)}{\(- 3\)}{}{}}
\End

\Data\ID{9000008004}
\Updated{2020-07-15 03:07:34}
\Level{A}
\SubArea{Rational functions}
\LANGen{
\Question{
Given the function \(f(x) = -\frac{8}
{x}\),
evaluate \(f(-4)\).}
{\(2\)}{\(- 4\)}{\(4\)}{\(32\)}{}{}}
\LANGpl{
\Question{
Dana jest funkcja \(f\colon y = -\frac{8}
{x}\). Oblicz \(f(-4)\).}
{\(2\)}{\(- 4\)}{\(4\)}{\(32\)}{}{}}
\LANGcs{
\Question{
Je dána funkce \(f\colon y = -\frac{8}
{x}\).
Výraz \(f(-4)\) 
nabývá hodnoty:}
{\(2\)}{\(- 4\)}{\(4\)}{\(32\)}{}{}}
\LANGsk{
\Question{
Daná je funkcia \(f\colon y = -\frac{8}
{x}\).
Určte hodnotu výrazu \(f(-4)\).}
{\(2\)}{\(- 4\)}{\(4\)}{\(32\)}{}{}}
\LANGes{
\Question{
Dada la función \(f(x) = -\frac{8}
{x}\),
calcula \(f(-4)\).}
{\(2\)}{\(- 4\)}{\(4\)}{\(32\)}{}{}}
\End

\Data\ID{9000008005}
\Updated{2023-09-10 11:09:17}
\Level{A}
\SubArea{Rational functions}
\LANGen{
\Question{
Given the function \(f(x)= -\frac{10}
 {x} \),
evaluate \(f(-5)\cdot f(2)\).}
{\(- 10\)}{\(2.5\)}{\(1\)}{\(2.5\)}{}{}}
\LANGpl{
\Question{
Dana jest funkcja \(f\colon y = -\frac{10}
 {x} \). 
Oblicz \(f(-5)\cdot f(2)\).}
{\(- 10\)}{\(2.5\)}{\(1\)}{\(2.5\)}{}{}}
\LANGcs{
\Question{
Je dána funkce \(f\colon y = -\frac{10}
 {x} \).
Výraz \(f(-5)\cdot f(2)\) 
nabývá hodnoty:}
{\(- 10\)}{\(2{,}5\)}{\(1\)}{\(2{,}5\)}{}{}}
\LANGsk{
\Question{
Daná je funkcia \(f\colon y = -\frac{10}
 {x} \). Určte hodnotu výrazu \(f(-5)\cdot f(2)\).}
{\(- 10\)}{\(2.5\)}{\(1\)}{\(2.5\)}{}{}}
\LANGes{
\Question{
Dada la función \(f(x)= -\frac{10}
 {x} \),
calcula \(f(-5)\cdot f(2)\).}
{\(- 10\)}{\(2.5\)}{\(1\)}{\(2.5\)}{}{}}
\End

\Data\ID{9000008006}
\Updated{2020-07-15 03:07:34}
\Level{A}
\SubArea{Rational functions}
\LANGen{
\Question{
Given the functions \(f(x) = \frac{2} 
{x}\) 
and \(g(x) = \frac{4} 
{x}\),
identify a true statement.}
{\(f(2) = g(4)\)}{\(f\left (\frac{1}
{2}\right ) = g(2)\)}{\(f(1) > g(2)\)}{\(f(4) < g(10)\)}{}{}}
\LANGpl{
\Question{
Dane są funkcje  \(f\colon y = \frac{2} 
{x}\) 
and \(g\colon y = \frac{4} 
{x}\). Która odpowiedź jest prawidłowa?}
{\(f(2) = g(4)\)}{\(f\left (\frac{1}
{2}\right ) = g(2)\)}{\(f(1) > g(2)\)}{\(f(4) < g(10)\)}{}{}}
\LANGcs{
\Question{
Jsou dány funkce \(f\colon y = \frac{2} 
{x}\) 
a \(g\colon y = \frac{4} 
{x}\).
Které z následujících tvrzení je správné?}
{\(f(2) = g(4)\)}{\(f\left (\frac{1}
{2}\right ) = g(2)\)}{\(f(1) > g(2)\)}{\(f(4) < g(10)\)}{}{}}
\LANGsk{
\Question{
Dané sú funkcie \(f\colon y = \frac{2} 
{x}\) 
a \(g\colon y = \frac{4} 
{x}\),
vyberte pravdivé tvrdenie.}
{\(f(2) = g(4)\)}{\(f\left (\frac{1}
{2}\right ) = g(2)\)}{\(f(1) > g(2)\)}{\(f(4) < g(10)\)}{}{}}
\LANGes{
\Question{
Dadas la funciones \(f(x) = \frac{2} 
{x}\) 
y \(g(x) = \frac{4} 
{x}\),
identifica la proposición verdadera.}
{\(f(2) = g(4)\)}{\(f\left (\frac{1}
{2}\right ) = g(2)\)}{\(f(1) > g(2)\)}{\(f(4) < g(10)\)}{}{}}
\End

\Data\ID{9000008007}
\Updated{2020-07-15 03:07:34}
\Level{A}
\SubArea{Rational functions}
\LANGen{
\Question{
Given the functions \(f(x) = -\frac{3}
{x}\) 
and \(g(x) = 6\),
solve \(f(x) = g(x)\).}
{\(-\frac{1} 
{2}\)}{\(- 2\)}{\(3\)}{\(6\)}{}{}}
\LANGpl{
\Question{
Dane są funkcje \(f\colon y = -\frac{3}
{x}\) 
i \(g\colon y = 6\). Oblicz \(f(x) = g(x)\).}
{\(-\frac{1} 
{2}\)}{\(- 2\)}{\(3\)}{\(6\)}{}{}}
\LANGcs{
\Question{
Jsou dány funkce \(f\colon y = -\frac{3}
{x}\) 
a \(g\colon y = 6\). Pro
které \(x\in D(f)\) 
platí, že \(f(x) = g(x)\)?}
{\(-\frac{1} 
{2}\)}{\(- 2\)}{\(3\)}{\(6\)}{}{}}
\LANGsk{
\Question{
Dané sú funkcie \(f\colon y = -\frac{3}
{x}\) 
a \(g\colon y = 6\). Pre
ktoré \(x\in D(f)\) 
platí \(f(x) = g(x)\)?}
{\(-\frac{1} 
{2}\)}{\(- 2\)}{\(3\)}{\(6\)}{}{}}
\LANGes{
\Question{
Dadas las funciones \(f(x) = -\frac{3}
{x}\) 
y \(g(x) = 6\),
resuelve \(f(x) = g(x)\).}
{\(-\frac{1} 
{2}\)}{\(- 2\)}{\(3\)}{\(6\)}{}{}}
\End

\Data\ID{9000008009}
\Updated{2023-09-10 11:09:11}
\Level{A}
\SubArea{Rational functions}
\LANGen{
\Question{
Given the function \(f(x) = \frac{5} 
{x}\),
find the function \(g\) such
that the graphs of \(f\) 
and \(g\) are symmetric
about the line \(y = x\).}
{\(g(x) = \frac{5} 
{x}\)}{\(g(x) = \frac{1} 
{x}\)}{\(g(x)= -\frac{2}
{x}\)}{\(g(x) = -\frac{5}
{x}\)}{}{}}
\LANGpl{
\Question{
Dana jest funkcja  \(f\colon y = \frac{5} 
{x}\). Wyznacz funkcję \(g\), taką, że 
wykresy funkcji \(f\) 
i \(g\) są symetryczne 
względem linii \(y = x\).}
{\(g\colon y = \frac{5} 
{x}\)}{\(g\colon y = \frac{1} 
{x}\)}{\(g\colon y = -\frac{2}
{x}\)}{\(g\colon y = -\frac{5}
{x}\)}{}{}}
\LANGcs{
\Question{
Je dána funkce \(f\colon y = \frac{5} 
{x}\).
Předpis funkce \(g\), jejíž
graf je souměrný podle osy \(I\). 
a \(III\). kvadrantu
s grafem funkce \(f\),
je:}
{\(g\colon y = \frac{5} 
{x}\)}{\(g\colon y = \frac{1} 
{x}\)}{\(g\colon y = -\frac{2}
{x}\)}{\(g\colon y = -\frac{5}
{x}\)}{}{}}
\LANGsk{
\Question{
Daná je funkcia \(f\colon y = \frac{5} 
{x}\).
Určte predpis funkcie \(g\), ktorej
graf je osovo súmerný podľa osi \(y = x\)
s grafom funkcie \(f\).}
{\(g\colon y = \frac{5} 
{x}\)}{\(g\colon y = \frac{1} 
{x}\)}{\(g\colon y = -\frac{2}
{x}\)}{\(g\colon y = -\frac{5}
{x}\)}{}{}}
\LANGes{
\Question{
Dada la función \(f(x) = \frac{5} 
{x}\),
encuentra la función  \(g\) para que las gráficas de \(f\) 
y \(g\) sean simétricas
respecto al eje \(y = x\).}
{\(g(x) = \frac{5} 
{x}\)}{\(g(x) = \frac{1} 
{x}\)}{\(g(x)= -\frac{2}
{x}\)}{\(g(x) = -\frac{5}
{x}\)}{}{}}
\End

\Data\ID{9000008010}
\Updated{2023-06-20 01:06:18}
\Level{A}
\SubArea{Rational functions}
\LANGen{
\Question{
Given the function 
\(f\colon y = -\frac{3}
{x}\).
find the function \(g\) such
that the 
graphs of \(f\) 
and \(g\) are 
symmetric
about the \(x\)-axis.}
{\(g(x) = \frac{3} 
{x}\)}{\(g(x) = -\frac{3}
{x}\)}{\(g(x)= -\frac{1}
{x}\)}{\(g(x) = \frac{2} 
{x}\)}{}{}}
\LANGpl{
\Question{
Dana jest funkcja  \(f\colon y = -\frac{3}
{x}\). Znajdź funkcję  \(g\), taką by wykresy
funkcji  \(f\) 
i \(g\) były symetryczne 
względem osi  \(x\) układu współrzędnych.}
{\(g\colon y = \frac{3} 
{x}\)}{\(g\colon y = -\frac{3}
{x}\)}{\(g\colon y = -\frac{1}
{x}\)}{\(g\colon y = \frac{2} 
{x}\)}{}{}}
\LANGcs{
\Question{
Je dána funkce \(f\colon y = -\frac{3}
{x}\).
Předpis funkce \(g\), jejíž
graf je souměrný podle osy \(x\) 
s grafem funkce \(f\),
je:}
{\(g\colon y = \frac{3} 
{x}\)}{\(g\colon y = -\frac{3}
{x}\)}{\(g\colon y = -\frac{1}
{x}\)}{\(g\colon y = \frac{2} 
{x}\)}{}{}}
\LANGsk{
\Question{
Daná je funkcia \(f\colon y = -\frac{3}{x}\). 
Určte predpis funkcie \(g\), ktorej graf je osovo súmerný podľa osi \(x\) s grafom funkcie \(f\).}
{\(g\colon y = \frac{3} 
{x}\)}{\(g\colon y = -\frac{3}
{x}\)}{\(g\colon y = -\frac{1}
{x}\)}{\(g\colon y = \frac{2} 
{x}\)}{}{}}
\LANGes{
\Question{
Dada la función \(f(x) = -\frac{3}
{x}\),
encuentra la función  \(g\) para que las gráficas de \(f\) 
y \(g\)  sean simétricas
respecto al eje  \(x\).}
{\(g(x) = \frac{3} 
{x}\)}{\(g(x) = -\frac{3}
{x}\)}{\(g(x)= -\frac{1}
{x}\)}{\(g(x) = \frac{2} 
{x}\)}{}{}}
\End

\Data\ID{9000009908}
\Updated{2020-07-15 03:07:34}
\Level{A}
\SubArea{Rational functions}
\LANGen{
\Question{
Consider a function 
\[ 
 f(x) = \frac{-3} 
 {x}
\]
defined on the domain \(\mathrm{Dom}(f) =\mathbb{R}\setminus \{ - 1.0\}\).
Find the range of this function.}
{\(\mathbb{R}\setminus \{0.3\}\)}{\(\mathbb{R}\setminus \{0\}\)}{\(\mathbb{R}\setminus \{ - 3.0\}\)}{\(\mathbb{R}\)}{}{}}
\LANGpl{
\Question{
Rozważ funkcję 
\[ 
 f\colon y = \frac{-3} 
 {x},
\] której dziedziną jest  \(\mathrm{Dom}(f) =\mathbb{R}\setminus \{ - 1{,}0\}\).
Wyznacz zakres tej funkcji.}
{\(\mathbb{R}\setminus \{0{,}3\}\)}{\(\mathbb{R}\setminus \{0\}\)}{\(\mathbb{R}\setminus \{ - 3{,}0\}\)}{\(\mathbb{R}\)}{}{}}
\LANGcs{
\Question{
Je dána funkce \( f\colon y = \frac{-3} 
 {x} \),
\(D(f) =\mathbb{R}\setminus \{ - 1{,}0\}\).
Jaký je obor hodnot této funkce?}
{\(\mathbb{R}\setminus \{0{,}3\}\)}{\(\mathbb{R}\setminus \{0\}\)}{\(\mathbb{R}\setminus \{ - 3{,}0\}\)}{\(\mathbb{R}\)}{}{}}
\LANGsk{
\Question{
Daná je funkcia \(f\colon y = \frac{-3} 
 {x} \), ktorej
\(D(f) =\mathbb{R}\setminus \{ - 1{,}0\}\).
Určte obor hodnôt funkcie \(f\).}
{\(\mathbb{R}\setminus \{0{,}3\}\)}{\(\mathbb{R}\setminus \{0\}\)}{\(\mathbb{R}\setminus \{ - 3{,}0\}\)}{\(\mathbb{R}\)}{}{}}
\LANGes{
\Question{
Considera la función
\[ 
 f(x) = \frac{-3} 
 {x}
\]
definida en el dominio \(\mathrm{Dom}(f) =\mathbb{R}\setminus \{ - 1.0\}\).
Determina el rango de la función.}
{\(\mathbb{R}\setminus \{0.3\}\)}{\(\mathbb{R}\setminus \{0\}\)}{\(\mathbb{R}\setminus \{ - 3.0\}\)}{\(\mathbb{R}\)}{}{}}
\End

\Data\ID{9000009910}
\Updated{2021-04-10 09:04:35}
\Level{A}
\SubArea{Rational functions}
\LANGen{
\Question{
A body is deformed continuously in a machine press. The density
\(\rho \) is inversely proportional
to the volume \(V \) of the body,
i.e. there exists a constant \(k\) 
such that 
\[ 
 \rho = \frac{k} 
{V }.
\]
Find the constant \(k\) 
(including the correct unit) if it is known that the density was
\(\rho = 25\: \frac{\mathrm{kg}}
{\mathrm{m}^{3}} \) when the body
had volume \(V = 2\, \mathrm{dm}^{3}\).}
{\(50\, \mathrm{g}\)}{\(12.5\, \mathrm{g}\)}{\(12.5\, \mathrm{m}\)}{\(50\, \mathrm{m}\)}{}{}}
\LANGpl{
\Question{
Ciało jest odkształcane w sposób ciągły przez prasę maszyny. Gęstość tego ciała 
\(\rho \) jest odwrotnie proporcjonalna
do objętości \(V \) tego ciała. Oznacza to, że istnieje stała \(k\) 
taka, że
\[ 
 \rho = \frac{k} 
{V }.
\]
Wyznacz stałą \(k\) jeśli gęstość ciała  wyniosła 
\(\rho = 25\: \frac{\mathrm{kg}}
{\mathrm{m}^{3}} \), a jego objętość była równa \(V = 2\, \mathrm{dm}^{3}\).}
{\(50\, \mathrm{g}\)}{\(12.5\, \mathrm{g}\)}{\(12.5\, \mathrm{m}\)}{\(50\, \mathrm{m}\)}{}{}}
\LANGcs{
\Question{
Dané těleso umístíme do lisu, kde plynule zmenšuje svůj objem. Jeho
průměrná hustota je tomuto objemu nepřímo úměrná. Určete koeficient
nepřímé úměrnosti (včetně jednotky), víme-li, že při objemu
\(2\, \mathrm{dm}^{3}\) má těleso
průměrnou hustotu \(25 \:\frac{\mathrm{kg}}
{\mathrm{m}^{3}} \).}
{\(50\, \mathrm{g}\)}{\(12{,}5\, \mathrm{g}\)}{\(12{,}5\, \mathrm{m}\)}{\(50\, \mathrm{m}\)}{}{}}
\LANGsk{
\Question{
Teleso vložené do lisu plynule zmenšuje svoj objem. Jeho
priemerná hustota je tomuto objemu nepriamo úmerná. Určte koeficient
nepriamej úmernosti (vrátane jednotiek), ak vieme, že pri objeme
\(2\, \mathrm{dm}^{3}\) má teleso
priemernú hustotu \(25 \:\frac{\mathrm{kg}}
{\mathrm{m}^{3}} \).}
{\(50\, \mathrm{g}\)}{\(12{,}5\, \mathrm{g}\)}{\(12{,}5\, \mathrm{m}\)}{\(50\, \mathrm{m}\)}{}{}}
\LANGes{
\Question{
Un cuerpo se deforma continuamente en una  prensa. La densidad
\(\rho \)  es inversamente proporcional
al volumen
\(V \) del cuerpo,
es decir, existe una constante \(k\) 
tal que
\[ 
 \rho = \frac{k} 
{V }.
\]
Determina la constatnte \(k\) 
(la unidad correcta incluida) si se sabe que la densidad era
\(\rho = 25\: \frac{\mathrm{kg}}
{\mathrm{m}^{3}} \) cuando el cuerpo
tenía un volumen de \(V = 2\, \mathrm{dm}^{3}\).}
{\(50\, \mathrm{g}\)}{\(12.5\, \mathrm{g}\)}{\(12.5\, \mathrm{m}\)}{\(50\, \mathrm{m}\)}{}{}}
\End

\Data\ID{9100003201}
\Updated{2022-05-29 12:05:06}
\Level{A}
\SubArea{Rational functions}
\LANGen{
\Question{
Identify a possible graph of the function
\(f(x) = \frac{1} 
{2x}\).}
{}{}{}{}{}{}}
\LANGpl{
\Question{
Który z przedstawionych poniżej wykresów jest wykresem funkcji 
\(f(x) = \frac{1} 
{2x}\)?}
{}{}{}{}{}{}}
\LANGcs{
\Question{
K funkci určené předpisem \(f(x) = \frac{1} 
{2x}\) 
najděte její graf.}
{}{}{}{}{}{}}
\LANGsk{
\Question{
Určte, ktorý z nasledujúcich grafov je graf funkcie
\(f(x)= \frac{1} 
{2x}\).}
{}{}{}{}{}{}}
\LANGes{
\Question{
Identifica una posible gráfica de la función
\(f(x) = \frac{1} 
{2x}\).}
{}{}{}{}{}{}}

\IMGanswer{000032_01-figure0.pdf}
\IMGanswer{000032_01-figure1.pdf}
\IMGanswer{000032_01-figure2.pdf}
\IMGanswer{000032_01-figure3.pdf}\End

\Data\ID{9100003202}
\Updated{2022-05-29 12:05:06}
\Level{A}
\SubArea{Rational functions}
\LANGen{
\Question{
Identify a possible graph of the function
\(f(x) = -\frac{2}
{x}\).}
{}{}{}{}{}{}}
\LANGpl{
\Question{
Który z przedstawionych poniżej wykresów jest wykresem funkcji 
\(f(x) = -\frac{2}
{x}\)?}
{}{}{}{}{}{}}
\LANGcs{
\Question{
K funkci určené předpisem \(f(x) = -\frac{2}
{x}\) 
najděte její graf.}
{}{}{}{}{}{}}
\LANGsk{
\Question{
Určte, ktorý z nasledujúcich grafov je graf funkcie
\(f(x) = -\frac{2}
{x}\).}
{}{}{}{}{}{}}
\LANGes{
\Question{
Identifica una posible gráfica de la función
\(f(x) = -\frac{2}
{x}\).}
{}{}{}{}{}{}}

\IMGanswer{000032_02-figure0.pdf}
\IMGanswer{000032_02-figure1.pdf}
\IMGanswer{000032_02-figure2.pdf}
\IMGanswer{000032_02-figure3.pdf}\End

\Data\ID{9100009902}
\Updated{2020-07-15 03:07:10}
\Level{A}
\SubArea{Rational functions}
\LANGen{
\Question{
The Newton's second law of motion 
\[ 
 F = m\cdot a
\]
states that the the acceleration \(a\) 
of the body is directly proportional to the acting force,
\(F\).
The positive constant of this proportionality is the mass of the body,
\(m\). This
law can be also rewritten in the form of inverse proportionality between convenient
other quantities. Which of the graphs describes the Newton's second law properly,
assuming one of the quantities which appear in the Newton's law is constant?}
{}{}{}{}{}{}}
\LANGpl{
\Question{
Druga zasada dynamiki Newtona
\[ 
 F = m\cdot a
\]
zakłada, że przyspieszenie ciała \(a\) 
jest wprost proporcjonalne do wartości siły wypadkowej 
\(F\).
Dodatnią stałą tej proporcjonalności jest masa ciała
\(m\). Zasadę tą możemy zapisać w postaci odwrotnej proporcjonalności pomiędzy odpowiednimi wartościami.
Który wykres poprawnie przedstawia II zasadę dynamiki Newtona, przy założeniu, ze jedna z wielkości jest stała?}
{}{}{}{}{}{}}
\LANGcs{
\Question{
Druhý Newtonův pohybový zákon se často zapisuje ve tvaru
\[F = m\cdot a,\] kde \(F\) je síla která působí na těleso o hmotnosti \(m\)
a uděluje mu zrychlení \(a\). Vyberte graf, který odpovídá uvedenému vztahu veličin.}
{}{}{}{}{}{}}
\LANGsk{
\Question{
Druhý Newtonov zákon pohybu \[F = m\cdot a\]
uvádza, že vektor zrýchlenia \(a\) je priamo úmerný pôsobiacej sile \(F\) a nepriamo úmerný hmotnosti telesa \(m\).
Ktorý z grafov správne popisuje druhý Newtonov zákon pohybu, za predpokladu, že jedna veličina je konštantná?}
{}{}{}{}{}{}}
\LANGes{
\Question{
La segunda ley del movimiento de Newton
\[ 
 F = m\cdot a
\]
establece que la aceleración \(a\) 
del cuerpo es directamente proporcional a la fuerza
\(F\). 
La constante positiva de esta proporcionalidad es la masa del cuerpo
\(m\). 
Esta ley también se puede reescribir en forma de proporcionalidad inversa entre otras magnitudes convenientes. ¿Cuál de las gráficas describe la segunda ley de Newton correctamente, suponiendo que una de las magnitudes que aparecen en la ley de Newton es constante?}
{}{}{}{}{}{}}

\IMGanswer{9100009902a.pdf}
\IMGanswer{9100009902b.pdf}
\IMGanswer{9100009902c.pdf}
\IMGanswer{9100009902d.pdf}\End

\Data\ID{9100009903}
\Updated{2020-07-15 03:07:10}
\Level{A}
\SubArea{Rational functions}
\LANGen{
\Question{
The resistance \(R\) 
of a wire is a function of a material constant
\(\rho \), the length
\(l\) and the cross
section \(S\).
It can be computed from the formula 
\[ 
 R =\rho \cdot \frac{l} 
{S}.
\]
Assuming two of the quantities are inversely proportional when the other two
quantities are constant, identify the graph which describes this relationship properly.}
{}{}{}{}{}{}}
\LANGpl{
\Question{
Opór  \(R\) 
przewodu jest funkcją  materiału stałego
\(\rho \), długości
\(l\)  i przekroju \(S\). Może być obliczony ze wzoru
\[ 
 R =\rho \cdot \frac{l} 
{S}.
\]
Zakładając, że dwie z podanych wielkości są odwrotnie proporcjonalne, gdy dwie pozostałe wielkości są stałe, wyznacz wykres, który prawidłowo przedstawia tę zależność.}
{}{}{}{}{}{}}
\LANGcs{
\Question{
Elektrický odpor tělesa \((R)\) závisí na vlastnostech materiálu \((\rho )\),
délce \((l)\) a ploše průřezu \((S)\) 
vztahem \[R =\rho \cdot \frac{l} {S}.\]
Vyberte graf, který vyjadřuje nepřímou úměrnost mezi dvěma z veličin za
předpokladu, že zbylé veličiny jsou konstantní.}
{}{}{}{}{}{}}
\LANGsk{
\Question{
Elektrický odpor vodiča \((R)\) závisí na vlastnostiach materiálu \((\rho )\),
dĺžke \((l)\) a ploche prierezu \((S)\). Môžme ho vyjadriť
vzťahom \[R =\rho \cdot \frac{l} {S}.\]
Vyberte graf, ktorý vyjadruje nepriamu úmernosť medzi dvoma veličinami za
predpokladu, že ostatné veličiny sú konštantné.}
{}{}{}{}{}{}}
\LANGes{
\Question{
La resistencia de un cable  \(R\) 
depende de las características del material
\(\rho \), la longitud
\(l\) y la sección \(S\).
Se puede calcular a partir de la fórmula
\[ 
 R =\rho \cdot \frac{l} 
{S}.
\]
Suponiendo que dos de las cantidades son inversamente proporcionales cuando las otras dos
magnitudes son constantes, identifica la gráfica que describe esta relación correctamente.}
{}{}{}{}{}{}}

\IMGanswer{000099_03-figure0.pdf}
\IMGanswer{000099_03-figure1.pdf}
\IMGanswer{000099_03-figure2.pdf}
\IMGanswer{000099_03-figure3.pdf}\End

\Data\ID{1003118301}
\Updated{2021-04-09 07:04:55}
\Level{B}
\SubArea{Rational functions}
\LANGen{
\Question{
Find the true statement about the function \( f(x)=-1+\frac3{2x-6} \).}
{The function \( f \) is decreasing on the interval \( (3;\infty) \).}{The function \( f \) is decreasing  on the interval \( (-3;\infty) \).}{The function \( f \) is decreasing  on the interval \( (-\infty;6) \).}{The function \( f \) is decreasing  on the interval \( (-1;\infty) \).}{}{}}
\LANGpl{
\Question{
Które ze zdań dotyczących funkcji  \( f(x)=-1+\frac3{2x-6} \) jest prawdziwe?}
{Funkcja  \( f \) jest malejąca w przedziale \( (3;\infty) \).}{Funkcja  \( f \) jest malejąca  w przedziale \( (-3;\infty) \).}{Funkcja \( f \) jest malejąca w przedziale \( (-\infty;6) \).}{Funkcja  \( f \) jest malejąca w przedziale \( (-1;\infty) \).}{}{}}
\LANGcs{
\Question{
Funkce \( f \) je  dána předpisem \(f(x)=-1+\frac3{2x-6} \). Vyberte pravdivý výrok.}
{Funkce \( f \) je klesající v intervalu \( (3;\infty) \).}{Funkce \( f \) je klesající v intervalu  \( (-3;\infty) \).}{Funkce \( f \) je klesající v intervalu  \( (-\infty;6) \).}{Funkce \( f \) je klesající v intervalu  \( (-1;\infty) \).}{}{}}
\LANGsk{
\Question{
Vyberte pravdivý výrok o funkcií  \( f(x)=-1+\frac3{2x-6} \).}
{Funkcia \( f \) je klesajúca na intervale \( (3;\infty) \).}{Funkcia \( f \) je klesajúca na intervale \( (-3;\infty) \).}{Funkcia \( f \) je klesajúca na intervale \( (-\infty;6) \).}{Funkcia \( f \) je klesajúca na intervale \( (-1;\infty) \).}{}{}}
\LANGes{
\Question{
Determina la proposición verdadera sobre la función \( f(x)=-1+\frac3{2x-6} \).}
{La función \( f \) decrece en el intervalo \( (3;\infty) \).}{La función \( f \) decrece en el intervalo \( (-3;\infty) \).}{La función \( f \)  decrece en el intervalo \( (-\infty;6) \).}{La función \( f \) decrece en el intervalo \( (-1;\infty) \).}{}{}}
\End

\Data\ID{1003118302}
\Updated{2021-04-09 07:04:23}
\Level{B}
\SubArea{Rational functions}
\LANGen{
\Question{
Find the true statement about the function \( f(x)=1-\frac2{0.5x-1};\ x\in[-3;1)\cup(2;6] \).}
{The function \( f \) has no maximum.}{The function \( f \) has the maximum at \( x=6 \).}{The function \( f \) has the minimum at \( x=-3 \).}{The function \( f \) is bounded.}{}{}}
\LANGpl{
\Question{
Które ze stwierdzeń dotyczących funkcji \( f(x)=1-\frac2{0{,}5x-1};\ x\in\langle-3;1)\cup(2;6\rangle \) jest prawdziwe?}
{Funkcja \( f \) nie posiada maksimum.}{Funkcja  \( f \) osiąga maksimum w \( x=6 \).}{Funkcja \( f \) osiąga maksimum w  \( x=-3 \).}{Funkcja  \( f \) jest ograniczona.}{}{}}
\LANGcs{
\Question{
Funkce \( f \)  je dána předpisem  \( f(x)=1-\frac2{0{,}5x-1};\ x\in\langle-3;1)\cup(2;6\rangle \). Vyberte pravdivý výrok.}
{Funkce \( f \) nemá maximum.}{Funkce \( f \) má maximum v bodě \( x=6 \).}{Funkce \( f \) má minimum v bodě  \( x=-3 \).}{Funkce \( f \) je omezená.}{}{}}
\LANGsk{
\Question{
Vyberte pravdivý výrok o funkcií \( f(x)=1-\frac2{0{,}5x-1};\ x\in\langle-3;1)\cup(2;6\rangle \).}
{Funkcia \( f \) nemá maximum.}{Funkcia \( f \) má maximum v bode \( x=6 \).}{Funkcia \( f \) má minimum v bode \( x=-3 \).}{Funkcia \( f \) je ohraničená.}{}{}}
\LANGes{
\Question{
Determina la proposición verdadera sobre la función  \( f(x)=1-\frac2{0.5x-1};\ x\in[-3;1)\cup(2;6] \).}
{La función \( f \) no tiene ningún máximo.}{La función \( f \) tiene su máximo en \( x=6 \).}{La función \( f \) tiene su mínimo en \( x=-3 \).}{La función \( f \) está acotada.}{}{}}
\End

\Data\ID{1003118303}
\Updated{2020-07-15 03:07:21}
\Level{B}
\SubArea{Rational functions}
\LANGen{
\Question{
Find the false statement about the function \( f(x)=\frac{4x+1}{3-2x};\ x\in[2;\infty) \).}
{The function \( f \) has no minimum.}{The function \( f \) is increasing.}{The function \( f \) has no maximum.}{The function \( f \) is bounded.}{}{}}
\LANGpl{
\Question{
Wybierz fałszywe zdanie dotyczące funkcji  \( f(x)=\frac{4x+1}{3-2x};\ x\in\langle2;\infty) \).}
{Funkcja  \( f \) nie ma minimum.}{Funkcja  \( f \) jest rosnąca.}{Funkcja  \( f \) nie ma maksimum.}{Funkcja \( f \) jest ograniczona.}{}{}}
\LANGcs{
\Question{
Funkce \( f \) je dána předpisem \(f(x)=\frac{4x+1}{3-2x};\ x\in\langle2;\infty) \). Vyberte nepravdivý výrok.}
{Funkce \( f \) nemá minimum.}{Funkce \( f \) je rostoucí.}{Funkce \( f \) nemá maximum.}{Funkce \( f \) je omezená.}{}{}}
\LANGsk{
\Question{
Vyberte nepravdivý výrok o funkcií \( f(x)=\frac{4x+1}{3-2x};\ x\in\langle2;\infty) \).}
{Funkcia \( f \) nemá minimum.}{Funkcia \( f \) je rastúca.}{Funkcia \( f \) nemá maximum.}{Funkcia \( f \) je ohraničená.}{}{}}
\LANGes{
\Question{
Determina la proposición  falsa  sobre la función  \( f(x)=\frac{4x+1}{3-2x};\ x\in[2;\infty) \).}
{La función \( f \) no tiene ninún mínimo.}{La función  \( f \) es creciente.}{La función  \( f \) no tiene ningún máximo.}{La función \( f \) está acotada.}{}{}}
\End

\Data\ID{1003123901}
\Updated{2020-07-15 03:07:21}
\Level{B}
\SubArea{Rational functions}
\LANGen{
\Question{
Given \( f(x)=\frac{2x+3}{x-5} \), find the true statement in the following list.}
{\( f(x)=2+\frac{13}{x-5} \)}{\( f(x)=2+\frac3{x-5} \)}{\( f(x)=2x+\frac3{x-5} \)}{\( f(x)=3+\frac{2x}{x-5} \)}{}{}}
\LANGpl{
\Question{
Dana jest funkcja \( f(x)=\frac{2x+3}{x-5} \), wybierz prawdziwe stwierdzenie.}
{\( f(x)=2+\frac{13}{x-5} \)}{\( f(x)=2+\frac3{x-5} \)}{\( f(x)=2x+\frac3{x-5} \)}{\( f(x)=3+\frac{2x}{x-5} \)}{}{}}
\LANGcs{
\Question{
Funkce \( f \) je dána předpisem \( f(x)=\frac{2x+3}{x-5} \). Vyberte pravdivý výrok.}
{\( f(x)=2+\frac{13}{x-5} \)}{\( f(x)=2+\frac3{x-5} \)}{\( f(x)=2x+\frac3{x-5} \)}{\( f(x)=3+\frac{2x}{x-5} \)}{}{}}
\LANGsk{
\Question{
Funkcia \(f\) je daná predpisom \( f(x)=\frac{2x+3}{x-5} \). Vyberte pravdivý výrok.}
{\( f(x)=2+\frac{13}{x-5} \)}{\( f(x)=2+\frac3{x-5} \)}{\( f(x)=2x+\frac3{x-5} \)}{\( f(x)=3+\frac{2x}{x-5} \)}{}{}}
\LANGes{
\Question{
Dada la función \( f(x)=\frac{2x+3}{x-5} \), determina la proposición verdadera.}
{\( f(x)=2+\frac{13}{x-5} \)}{\( f(x)=2+\frac3{x-5} \)}{\( f(x)=2x+\frac3{x-5} \)}{\( f(x)=3+\frac{2x}{x-5} \)}{}{}}
\End

\Data\ID{1003123902}
\Updated{2020-07-15 03:07:21}
\Level{B}
\SubArea{Rational functions}
\LANGen{
\Question{
Given \( f(x)=\frac{3x-2}{2x-3} \), find the false statement in the following list.}
{\( f(x)=\frac32-\frac2{2x-3} \)}{\( f(x)=\frac32+\frac5{4\left(x-\frac32\right)} \)}{\( f(x)=\frac32-\frac5{6-4x} \)}{\( f(x)=\frac32+\frac{\frac52}{2x-3} \)}{}{}}
\LANGpl{
\Question{
Dana jest funkcja \( f(x)=\frac{3x-2}{2x-3} \), wybierz z poniższych stwierdzeń fałszywe.}
{\( f(x)=\frac32-\frac2{2x-3} \)}{\( f(x)=\frac32+\frac5{4\left(x-\frac32\right)} \)}{\( f(x)=\frac32-\frac5{6-4x} \)}{\( f(x)=\frac32+\frac{\frac52}{2x-3} \)}{}{}}
\LANGcs{
\Question{
Funkce \( f \) je dána předpisem \( f(x)=\frac{3x-2}{2x-3} \). Vyberte nepravdivý výrok.}
{\( f(x)=\frac32-\frac2{2x-3} \)}{\( f(x)=\frac32+\frac5{4\left(x-\frac32\right)} \)}{\( f(x)=\frac32-\frac5{6-4x} \)}{\( f(x)=\frac32+\frac{\frac52}{2x-3} \)}{}{}}
\LANGsk{
\Question{
Funkcia \(f\) je daná predpisom \( f(x)=\frac{3x-2}{2x-3} \). Vyberte nepravdivý výrok.}
{\( f(x)=\frac32-\frac2{2x-3} \)}{\( f(x)=\frac32+\frac5{4\left(x-\frac32\right)} \)}{\( f(x)=\frac32-\frac5{6-4x} \)}{\( f(x)=\frac32+\frac{\frac52}{2x-3} \)}{}{}}
\LANGes{
\Question{
Dada la función \( f(x)=\frac{3x-2}{2x-3} \), determina la proposición  falsa.}
{\( f(x)=\frac32-\frac2{2x-3} \)}{\( f(x)=\frac32+\frac5{4\left(x-\frac32\right)} \)}{\( f(x)=\frac32-\frac5{6-4x} \)}{\( f(x)=\frac32+\frac{\frac52}{2x-3} \)}{}{}}
\End

\Data\ID{1003123903}
\Updated{2020-07-15 03:07:21}
\Level{B}
\SubArea{Rational functions}
\LANGen{
\Question{
Given \( f(x)=\frac{2x-1}{3x+2} \), find the true statement in the following list.}
{\( f(x)=\frac23-\frac7{9\left(x+\frac23\right)} \)}{\( f(x)=\frac23+\frac1{3\left(3x+2\right)} \)}{\( f(x)=\frac32-\frac{\frac73}{3x+2} \)}{\( f(x)=2-\frac1{3x+2} \)}{}{}}
\LANGpl{
\Question{
Dana jest funkcja \( f(x)=\frac{2x-1}{3x+2} \). Które ze stwierdzeń z poniższej listy jest prawdziwe?}
{\( f(x)=\frac23-\frac7{9\left(x+\frac23\right)} \)}{\( f(x)=\frac23+\frac1{3\left(3x+2\right)} \)}{\( f(x)=\frac32-\frac{\frac73}{3x+2} \)}{\( f(x)=2-\frac1{3x+2} \)}{}{}}
\LANGcs{
\Question{
Funkce \( f \) je dána předpisem \( f(x)=\frac{2x-1}{3x+2} \). Vyberte pravdivý výrok.}
{\( f(x)=\frac23-\frac7{9\left(x+\frac23\right)} \)}{\( f(x)=\frac23+\frac1{3\left(3x+2\right)} \)}{\( f(x)=\frac32-\frac{\frac73}{3x+2} \)}{\( f(x)=2-\frac1{3x+2} \)}{}{}}
\LANGsk{
\Question{
Funkcia \(f\) je daná predpisom \( f(x)=\frac{2x-1}{3x+2} \). Vyberte pravdivý výrok.}
{\( f(x)=\frac23-\frac7{9\left(x+\frac23\right)} \)}{\( f(x)=\frac23+\frac1{3\left(3x+2\right)} \)}{\( f(x)=\frac32-\frac{\frac73}{3x+2} \)}{\( f(x)=2-\frac1{3x+2} \)}{}{}}
\LANGes{
\Question{
Dada la función \( f(x)=\frac{2x-1}{3x+2} \), determina la proposición verdadera.}
{\( f(x)=\frac23-\frac7{9\left(x+\frac23\right)} \)}{\( f(x)=\frac23+\frac1{3\left(3x+2\right)} \)}{\( f(x)=\frac32-\frac{\frac73}{3x+2} \)}{\( f(x)=2-\frac1{3x+2} \)}{}{}}
\End

\Data\ID{1003123904}
\Updated{2021-04-09 07:04:25}
\Level{B}
\SubArea{Rational functions}
\LANGen{
\Question{
Let
\[ \begin{aligned}
f(x)&=-2-\frac5{3(x-1)}, \\
g(x)&=\frac{2x-\frac13}{1-x}, \\
h(x)&=\frac{1-6x}{3x-3}, \\ 
j(x)&=\frac5{3x-3}-2. 
\end{aligned}
\]
From the following list choose the false statement.}
{\( h=j \)}{\( f=g \)}{\( f=h \)}{\( g=h \)}{}{}}
\LANGpl{
\Question{
Niech
\[ \begin{aligned}
f(x)&=-2-\frac5{3(x-1)}, \\
g(x)&=\frac{2x-\frac13}{1-x}, \\
h(x)&=\frac{1-6x}{3x-3}, \\ 
j(x)&=\frac5{3x-3}-2. 
\end{aligned}
\]
Wybierz fałszywe stwierdzenie.}
{\( h=j \)}{\( f=g \)}{\( f=h \)}{\( g=h \)}{}{}}
\LANGcs{
\Question{
Funkce \( f \), \( g \), \( h \), \( j \) jsou dány předpisy 
\[ \begin{aligned}
f(x)&=-2-\frac5{3(x-1)}, \\
g(x)&=\frac{2x-\frac13}{1-x}, \\
h(x)&=\frac{1-6x}{3x-3}, \\ 
j(x)&=\frac5{3x-3}-2. 
\end{aligned}
\]
Vyberte nepravdivý výrok.}
{\( h=j \)}{\( f=g \)}{\( f=h \)}{\( g=h \)}{}{}}
\LANGsk{
\Question{
Funkcie \(f\), \(g\), \(h\), \(j\) sú dané predpismi
\[ \begin{aligned}
f(x)&=-2-\frac5{3(x-1)}, \\
g(x)&=\frac{2x-\frac13}{1-x}, \\
h(x)&=\frac{1-6x}{3x-3}, \\ 
j(x)&=\frac5{3x-3}-2. 
\end{aligned}
\]
Vyberte nepravdivý výrok.}
{\( h=j \)}{\( f=g \)}{\( f=h \)}{\( g=h \)}{}{}}
\LANGes{
\Question{
Sean
\[ \begin{aligned}
f(x)&=-2-\frac5{3(x-1)}, \\
g(x)&=\frac{2x-\frac13}{1-x}, \\
h(x)&=\frac{1-6x}{3x-3}, \\ 
j(x)&=\frac5{3x-3}-2. 
\end{aligned}
\]
De la siguinte lista, determina la proposición falsa .}
{\( h=j \)}{\( f=g \)}{\( f=h \)}{\( g=h \)}{}{}}
\End

\Data\ID{1003123905}
\Updated{2021-04-09 07:04:35}
\Level{B}
\SubArea{Rational functions}
\LANGen{
\Question{
Let
\[ \begin{aligned}
f(x)&=\frac{6x-5}{5-2x}, \\
g(x)&=-3-\frac5{x-2}, \\
h(x)&=-3-\frac5{x-2.5}, \\ 
j(x)&=\frac{3x-1}{2-x}.
\end{aligned}
\] 
From the following list choose the true statement.}
{\( j=g \)}{\( f=g \)}{\( j=h \)}{\( f=j \)}{}{}}
\LANGpl{
\Question{
Niech 
\[ \begin{aligned}
f(x)&=\frac{6x-5}{5-2x}, \\
g(x)&=-3-\frac5{x-2}, \\
h(x)&=-3-\frac5{x-2{,}5}, \\ 
j(x)&=\frac{3x-1}{2-x}.
\end{aligned}
\] 
Wybierz prawdziwe stwierdzenie.}
{\( j=g \)}{\( f=g \)}{\( j=h \)}{\( f=j \)}{}{}}
\LANGcs{
\Question{
Funkce \( f \), \( g \), \( h \), \( j \) jsou dány předpisy 
\[ \begin{aligned}
f(x)&=\frac{6x-5}{5-2x}, \\
g(x)&=-3-\frac5{x-2}, \\
h(x)&=-3-\frac5{x-2{,}5}, \\ 
j(x)&=\frac{3x-1}{2-x}.
\end{aligned}
\] 
Vyberte pravdivý výrok.}
{\( j=g \)}{\( f=g \)}{\( j=h \)}{\( f=j \)}{}{}}
\LANGsk{
\Question{
Funkcie \(f\), \(g\), \(h\), \(j\) sú dané predpismi
\[ \begin{aligned}
f(x)&=\frac{6x-5}{5-2x}, \\
g(x)&=-3-\frac5{x-2}, \\
h(x)&=-3-\frac5{x-2{,}5}, \\ 
j(x)&=\frac{3x-1}{2-x}.
\end{aligned}
\] 
Vyberte pravdivý výrok.}
{\( j=g \)}{\( f=g \)}{\( j=h \)}{\( f=j \)}{}{}}
\LANGes{
\Question{
Sean \[ \begin{aligned}
f(x)&=\frac{6x-5}{5-2x}, \\
g(x)&=-3-\frac5{x-2}, \\
h(x)&=-3-\frac5{x-2.5}, \\ 
j(x)&=\frac{3x-1}{2-x}.
\end{aligned}
\] 
De la siguiente lista, determina la proposición verdadera.}
{\( j=g \)}{\( f=g \)}{\( j=h \)}{\( f=j \)}{}{}}
\End

\Data\ID{1003124601}
\Updated{2020-07-15 03:07:21}
\Level{B}
\SubArea{Rational functions}
\LANGen{
\Question{
Let \( f(x)=\frac{2x}{x^2-1} \). Find the true statement.}
{\( \forall x\in(-\infty;-1)\cup(0;1)\colon f(x) < 0 \).}{The domain of \( f \) is \( (-\infty;1)\cup(1;\infty) \).}{\( \forall x\in(-1;1)\colon f(x) \leq 0 \).}{The domain of \( f \) is \( (-\infty;-1)\cup(-1;0)\cup(0;1)\cup(1;\infty) \).}{}{}}
\LANGpl{
\Question{
Niech \( f(x)=\frac{2x}{x^2-1} \). Wybierz prawdziwe stwierdzenie.}
{\( \forall x\in(-\infty;-1)\cup(0;1)\colon f(x) < 0 \).}{Dziedziną funkcji \( f \) jest \( (-\infty;1)\cup(1;\infty) \).}{\( \forall x\in(-1;1)\colon f(x) \leq 0 \).}{Dziedziną funkcji  \( f \) jest \( (-\infty;-1)\cup(-1;0)\cup(0;1)\cup(1;\infty) \).}{}{}}
\LANGcs{
\Question{
Funkce  \( f \) je dána předpisem \( f(x)=\frac{2x}{x^2-1} \). Vyberte pravdivý výrok.}
{\( \forall x\in(-\infty;-1)\cup(0;1)\colon f(x) < 0 \).}{Definiční obor funkce \( f \) je množina \( (-\infty;1)\cup(1;\infty) \).}{\( \forall x\in(-1;1)\colon f(x) \leq 0 \).}{Definiční obor funkce \( f \) je množina \( (-\infty;-1)\cup(-1;0)\cup(0;1)\cup(1;\infty) \).}{}{}}
\LANGsk{
\Question{
Funkcia \(f\) je daná predpisom \( f(x)=\frac{2x}{x^2-1} \). Vyberte pravdivý výrok.}
{\( \forall x\in(-\infty;-1)\cup(0;1)\colon f(x) < 0 \).}{Definičný obor funkcie \( f \) je \( (-\infty;1)\cup(1;\infty) \).}{\( \forall x\in(-1;1)\colon f(x) \leq 0 \).}{Definičný obor funkcie \( f \) je \( (-\infty;-1)\cup(-1;0)\cup(0;1)\cup(1;\infty) \).}{}{}}
\LANGes{
\Question{
Sea \( f(x)=\frac{2x}{x^2-1} \). Determina la proposición verdadera.}
{\( \forall x\in(-\infty;-1)\cup(0;1)\colon f(x) < 0 \).}{El dominio de la función  \( f \) es \( (-\infty;1)\cup(1;\infty) \).}{\( \forall x\in(-1;1)\colon f(x) \leq 0 \).}{El dominio de la función \( f \) es \( (-\infty;-1)\cup(-1;0)\cup(0;1)\cup(1;\infty) \).}{}{}}
\End

\Data\ID{1103030902}
\Updated{2022-04-23 05:04:42}
\Level{B}
\SubArea{Rational functions}
\IMG{1103030902.pdf}
\LANGen{
\Question{
A part of the graph of the function \( f(x)=\frac4x \)  is shown in the picture. Identify which of the following statements is true.}
{The function \( g \) defined by \( g(x)=\left|f(x)\right| \) is bounded below.}{The function \( f \) is bounded below.}{The function $h$ defined by \( h(x)=-f(x) \)  is bounded below.}{The function $m$ defined by \( m(x)=f(x)+4 \) is bounded below.}{}{}}
\LANGpl{
\Question{
Część wykresu funkcji \( f(x)=\frac4x \)  przedstawiona jest na rysunku. Które z poniższych zdań jest prawdziwe?}
{Funkcja \( g \) określona przez \( g(x)=\left|f(x)\right| \) jest ograniczona z dołu.}{Funkcja \( f \) jest ograniczona z dołu.}{Funkcja $h$ określona przez \( h(x)=-f(x) \)  jest ograniczona z dołu.}{Funkcja $m$ określona przez \( m(x)=f(x)+4 \) jest ograniczona z dołu.}{}{}}
\LANGcs{
\Question{
Na obrázku je část grafu funkce \( f(x)=\frac4x \). Vyberte pravdivý výrok.}
{Funkce \( g(x)=\left|f(x)\right| \) je zdola omezená.}{Funkce  \( f \) je zdola omezená.}{Funkce  \( h(x)=-f(x) \)  je zdola omezená.}{Funkce \( m(x)=f(x)+4 \) je zdola omezená.}{}{}}
\LANGsk{
\Question{
Na obrázku je časť grafu funkcie \( f(x)=\frac4x \). Vyberte pravdivý výrok.}
{Funkcia \( g \) definovaná ako \( g(x)=\left|f(x)\right| \) je zdola ohraničená.}{Funkcia \( f \) je zdola ohraničená.}{Funkcia \( h \) definovaná ako \( h(x)=-f(x) \) je zdola ohraničená.}{Funkcia \( m \) definovaná ako \( m(x)=f(x)+4 \) je zdola ohraničená.}{}{}}
\LANGes{
\Question{
Una parte de la gráfica de la función \( f(x)=\frac4x \)  se muestra en la imagen. Identifica cuál de las siguientes proposiciones es verdadera.}
{La función $g$ definida como \( g(x)=\left|f(x)\right| \) está acotada inferiormente.}{La función \( f \) está acotada inferiormente.}{La función $h$ definida como \( h(x)=-f(x) \) está acotada inferiormente.}{La función  $m$ definida como \( m(x)=f(x)+4 \) está acotada inferiormente.}{}{}}
\End

\Data\ID{1103124501}
\Updated{2020-07-15 03:07:10}
\Level{B}
\SubArea{Rational functions}
\LANGen{
\Question{
Identify which of the graphs below represents the function \( f(x)=\frac{4x-3}{3x+4};\ x\in [-1;5] \).}
{}{}{}{}{}{}}
\LANGpl{
\Question{
Który z poniższych wykresów przedstawia funkcję  \( f(x)=\frac{4x-3}{3x+4};\ x\in\langle-1;5\rangle \)?}
{}{}{}{}{}{}}
\LANGcs{
\Question{
Na kterém obrázku je graf funkce \( f(x)=\frac{4x-3}{3x+4};\ x\in\langle-1;5\rangle \)?}
{}{}{}{}{}{}}
\LANGsk{
\Question{
Ktorý z nasledujúcich grafov je graf funkcie \( f(x)=\frac{4x-3}{3x+4};\ x\in\langle-1;5\rangle \)?}
{}{}{}{}{}{}}
\LANGes{
\Question{
Identifica cuál de las gráficas representa la función\( f(x)=\frac{4x-3}{3x+4};\ x\in [-1;5] \).}
{}{}{}{}{}{}}

\IMGanswer{1103124501a_0.pdf}
\IMGanswer{1103124501b_0.pdf}
\IMGanswer{1103124501c_0.pdf}
\IMGanswer{1103124501d_0.pdf}\End

\Data\ID{1103129201}
\Updated{2021-04-09 07:04:11}
\Level{B}
\SubArea{Rational functions}
\LANGen{
\Question{
The thin lens equation \( \frac1a+\frac1{a'}=\frac1f \) describes the quantitative relationship between the object distance \( a \), the image distance \( a' \), and the focal length \( f \). Let the focal length of a thin lens be \( 0{.}5\,\mathrm{m} \). Choose the picture that shows the object distance vs. image distance graph, if \( a\in [0.1\,\mathrm{m};0.4\,\mathrm{m}]\cup [0.6\,\mathrm{m};3.0\,\mathrm{m}] \).}
{}{}{}{}{}{}}
\LANGpl{
\Question{
Równanie cienkich soczewek \( \frac1a+\frac1{a'}=\frac1f \) opisuje ilościową zależność pomiędzy odległością obiektu \( a \), odległością obrazu \( a' \), a długością ogniskową \( f \). Niech długość ogniskowych cienkiej soczewki wynosi \( 0{,}5\,\mathrm{m} \). Wybierz rysunek, który przedstawia odległość przedmiotu od wykresu odległości obrazu jeśli \( a\in\langle0{,}1\,\mathrm{m};0{,}4\,\mathrm{m}\rangle\cup\langle0{,}6\,\mathrm{m};3{,}0\,\mathrm{m}\rangle \).}
{}{}{}{}{}{}}
\LANGcs{
\Question{
Polohu obrazu předmětu zobrazeného tenkou čočkou popisuje zobrazovací rovnice: \( \frac1a+\frac1{a'}=\frac1f \). Na kterém obrázku je graf závislosti obrazové vzdálenosti \( a' \) na předmětové vzdálenosti  \( a \) při zobrazení tenkou čočkou s ohniskovou vzdálenosti \( f = 0{,}5\,\mathrm{m} \) pro \( a\in\langle0{,}1\,\mathrm{m};0{,}4\,\mathrm{m}\rangle\cup\langle0{,}6\,\mathrm{m};3{,}0\,\mathrm{m}\rangle \)?}
{}{}{}{}{}{}}
\LANGsk{
\Question{
Polohu obrazu predmetu zobrazeného tenkou šošovkou popisuje zobrazovacia rovnica: \( \frac1a+\frac1{a'}=\frac1f \). Na ktorom obrázku je graf závislosti obrazovej vzdialenosti \( a' \) na predmetovej vzdialenosti  \( a \) pri zobrazení tenkou šošovkou s ohniskovou vzdialenosťou \( f = 0{,}5\,\mathrm{m} \) pre \( a\in\langle0{,}1\,\mathrm{m};0{,}4\,\mathrm{m}\rangle\cup\langle0{,}6\,\mathrm{m};3{,}0\,\mathrm{m}\rangle \)?}
{}{}{}{}{}{}}
\LANGes{
\Question{
La ecuación de lente delgada \( \frac1a+\frac1{a'}=\frac1f \) describe la relación cuantitativa entre la distancia del objeto \( a \), la distancia de la imagen \( a' \), y la distancia focal \( f \). Sea la distancia focal de una lente delgada \( 0{.}5\,\mathrm{m} \). Elige la imagen que muestre la representación gráfica de la distancia del objeto respecto a la distancia de la imagen, si \( a\in [0.1\,\mathrm{m};0.4\,\mathrm{m}]\cup [0.6\,\mathrm{m};3.0\,\mathrm{m}] \).}
{}{}{}{}{}{}}

\IMGanswer{1103129201a_1.pdf}
\IMGanswer{1103129201b_0.pdf}
\IMGanswer{1103129201c.pdf}
\IMGanswer{1103129201d.pdf}\End

\Data\ID{2000006701}
\Updated{2022-08-25 10:08:29}
\Level{B}
\SubArea{Rational functions}
\IMG{2000006701_6-figure0.pdf}
\LANGen{
\Question{
A part of the graph of the function \( f(x)=-\frac2x \)  is shown in the picture. Identify which of the following statements is true.}
{The function \( g \) defined by \( g(x)=-\left|f(x)\right| \) is an even function.}{The function \( g \) defined by \( g(x)=-\left|f(x)\right| \) is bounded below.}{The function $f$ is decreasing on \( (-\infty;0)\).}{The function $m$ defined by \( m(x)=f(x)-3 \) is bounded.}{}{}}
\LANGpl{
\Question{
Na rysunku pokazano fragment wykresu funkcji \( f(x)=-\frac2x \). Określ, które z poniższych stwierdzeń jest prawdziwe.}
{Funkcja \( g \) określona przez \( g(x)=-\left|f(x)\right| \) jest funkcją parzystą.}{Funkcja \( g \) określona przez \( g(x)=-\left|f(x)\right| \) jest ograniczona z dołu.}{Funkcja $f$ maleje w przedziale \( (-\infty;0)\).}{Funkcja $m$ określona przez \( m(x)=f(x)-3 \) jest ograniczona.}{}{}}
\LANGcs{
\Question{
Na obrázku je část grafu funkce \( f(x)=-\frac2x \). Vyberte pravdivý výrok.}
{Funkce   \( g(x)=-\left|f(x)\right| \) je sudá.}{Funkce  \( g(x)=-\left|f(x)\right| \) je omezená zdola.}{Funkce $f$ je klesající na intervalu \( (-\infty;0)\).}{Funkce  \( m(x)=f(x)-3 \) je omezená.}{}{}}
\LANGsk{
\Question{
Na obrázku je časť grafu funkcie \( f(x)=-\frac2x \). Vyberte pravdivý výrok.}
{Funkcia \( g \) definovaná \( g(x)=-\left|f(x)\right| \) je párna.}{Funkcia \( g \) definovaná \( g(x)=-\left|f(x)\right| \) je ohraničená zdola.}{Funkcia $f$ klesá na \( (-\infty;0)\).}{Funkcia $m$ definovaná na \( m(x)=f(x)-3 \) je ohraničená.}{}{}}
\LANGes{
\Question{
En la imagen se puede ver una parte de la gráfica de la función \( f(x)=-\frac2x \). Identifica cuál de las siguientes proposiciones es verdadera.}
{La función \( g \) definida por \( g(x)=-\left|f(x)\right| \) es una función par.}{La función \( g \) definida por \( g(x)=-\left|f(x)\right| \) está acotada inferiormente.}{La función $f$ está disminuyendo en \( (-\infty;0)\).}{La función $m$ definida por \( m(x)=f(x)-3 \) está acotada.}{}{}}
\End

\Data\ID{2000006704}
\Updated{2023-01-23 10:01:35}
\Level{B}
\SubArea{Rational functions}
\LANGen{
\Question{
Let \(X\) and
\(Y \) be intersections of the
graph of the function \(f(x) = \frac{3x-5}
{2+x}\) 
with \(x\)-
and \(y\)-axis,
respectively. Find these points.}
{\(X = \left[\frac{5}{3};0\right]\),
\(Y = \left[0;-\frac{5}{2}\right]\)}{\(X = \left[-\frac{5}{2};0\right]\),
\(Y = \left[0;\frac{5}{3}\right]\)}{\(X = \left[0;\frac{5}{3}\right]\),
\(Y = \left[-\frac{5}{2};0\right]\)}{\(X = \left[\frac{5}{2};0\right]\),
\(Y = \left[0;-\frac{5}{3}\right]\)}{}{}}
\LANGpl{
\Question{
\(X\) i
\(Y \)  to punkty przecięcia wykresu funkcji \(f(x) = \frac{3x-5}
{2+x}\) 
z osiami  \(x\) i \(y\). Znajdź te punkty.}
{\(X = \left[\frac{5}{3};0\right]\),
\(Y = \left[0;-\frac{5}{2}\right]\)}{\(X = \left[-\frac{5}{2};0\right]\),
\(Y = \left[0;\frac{5}{3}\right]\)}{\(X = \left[0;\frac{5}{3}\right]\),
\(Y = \left[-\frac{5}{2};0\right]\)}{\(X = \left[\frac{5}{2};0\right]\),
\(Y = \left[0;-\frac{5}{3}\right]\)}{}{}}
\LANGcs{
\Question{
Je dána funkce \(f(x) = \frac{3x-5}{2+x}\). Průsečíky grafu funkce \(f\) s osami \(x\), \(y\) označme po řadě \(X\), \(Y \). Určete tyto body.}
{\(X = \left[\frac{5}{3};0\right]\),
\(Y = \left[0;-\frac{5}{2}\right]\)}{\(X = \left[-\frac{5}{2};0\right]\),
\(Y = \left[0;\frac{5}{3}\right]\)}{\(X = \left[0;\frac{5}{3}\right]\),
\(Y = \left[-\frac{5}{2};0\right]\)}{\(X = \left[\frac{5}{2};0\right]\),
\(Y = \left[0;-\frac{5}{3}\right]\)}{}{}}
\LANGsk{
\Question{
Nech \(X\) a
\(Y\) sú priesečníky grafu funkcie \(f(x) = \frac{3x-5}
{2+x}\) 
s osou \(x\)
a osou \(y\).
Nájdite tieto body.}
{\(X = \left[\frac{5}{3};0\right]\),
\(Y = \left[0;-\frac{5}{2}\right]\)}{\(X = \left[-\frac{5}{2};0\right]\),
\(Y = \left[0;\frac{5}{3}\right]\)}{\(X = \left[0;\frac{5}{3}\right]\),
\(Y = \left[-\frac{5}{2};0\right]\)}{\(X = \left[\frac{5}{2};0\right]\),
\(Y = \left[0;-\frac{5}{3}\right]\)}{}{}}
\LANGes{
\Question{
Sean \(X\) e
\(Y\) las intersecciones de la
gráfica de la función \(f(x) = \frac{3x-5}
{2+x}\) 
con los ejes \(x\)
e \(y\). Determina estos puntos.}
{\(X = \left[\frac{5}{3};0\right]\),
\(Y = \left[0;-\frac{5}{2}\right]\)}{\(X = \left[-\frac{5}{2};0\right]\),
\(Y = \left[0;\frac{5}{3}\right]\)}{\(X = \left[0;\frac{5}{3}\right]\),
\(Y = \left[-\frac{5}{2};0\right]\)}{\(X = \left[\frac{5}{2};0\right]\),
\(Y = \left[0;-\frac{5}{3}\right]\)}{}{}}
\End

\Data\ID{2000006705}
\Updated{2023-01-23 10:01:35}
\Level{B}
\SubArea{Rational functions}
\LANGen{
\Question{
Calculate the range of the function \(f(x)= \frac{-2} 
{2x-1} + 3\).}
{\((-\infty ;3)\cup (3;\infty )\)}{\(\left(-\infty ;\frac12)\cup (\frac12;\infty \right)\)}{\((-\infty ;-3)\cup (-3;\infty )\)}{\((-\infty ;1)\cup (1;\infty )\)}{}{}}
\LANGpl{
\Question{
Znajdź zbiór wartości funkcji \(f(x)= \frac{-2} 
{2x-1} + 3\).}
{\((-\infty ;3)\cup (3;\infty )\)}{\(\left(-\infty ;\frac12)\cup (\frac12;\infty \right)\)}{\((-\infty ;-3)\cup (-3;\infty )\)}{\((-\infty ;1)\cup (1;\infty )\)}{}{}}
\LANGcs{
\Question{
Je dána funkce \(f(x)= \frac{-2}{2x-1} + 3\). Určete její obor hodnot.}
{\((-\infty ;3)\cup (3;\infty )\)}{\(\left(-\infty ;\frac12)\cup (\frac12;\infty \right)\)}{\((-\infty ;-3)\cup (-3;\infty )\)}{\((-\infty ;1)\cup (1;\infty )\)}{}{}}
\LANGsk{
\Question{
Určte obor hodnôt funkcie \(f(x)= \frac{-2} 
{2x-1} + 3\).}
{\((-\infty ;3)\cup (3;\infty )\)}{\(\left(-\infty ;\frac12)\cup (\frac12;\infty \right)\)}{\((-\infty ;-3)\cup (-3;\infty )\)}{\((-\infty ;1)\cup (1;\infty )\)}{}{}}
\LANGes{
\Question{
Determina el rango de la función \(f(x)= \frac{-2} 
{2x-1} + 3\).}
{\((-\infty ;3)\cup (3;\infty )\)}{\(\left(-\infty ;\frac12)\cup (\frac12;\infty \right)\)}{\((-\infty ;-3)\cup (-3;\infty )\)}{\((-\infty ;1)\cup (1;\infty )\)}{}{}}
\End

\Data\ID{2010009901}
\Updated{2022-08-25 10:08:18}
\Level{B}
\SubArea{Rational functions}
\LANGen{
\Question{
Find the domain \(\mathrm{Dom}(f)\) 
and range \(\mathop{\mathrm{Ran}}(f)\) of
the function \(f(x) = \frac{x-3} 
{x+1}\).}
{\begin{align*}
\mathrm{Dom}(f) &= (-\infty ;-1)\cup (-1;\infty ),\\
\mathop{\mathrm{Ran}}(f) &= (-\infty ;1)\cup (1;\infty )
\end{align*}}{\begin{align*}
\mathrm{Dom}(f) &= (-\infty ;1)\cup (1;\infty ),\\
\mathop{\mathrm{Ran}}(f) &= (-\infty ;-1)\cup (-1;\infty )
\end{align*}}{\begin{align*}
\mathrm{Dom}(f) &= (-\infty ;3)\cup (3;\infty ),\\
\mathop{\mathrm{Ran}}(f) &= (-\infty ;-1)\cup (-1;\infty )
\end{align*}}{\begin{align*}
\mathrm{Dom}(f) &= (-\infty ;-3)\cup (-3;\infty ),\\
\mathop{\mathrm{Ran}}(f) &= (-\infty ;1)\cup (1;\infty )
\end{align*}}{}{}}
\LANGpl{
\Question{
Znajdź dziedzinę \(\mathrm{Dom}(f)\) 
i zbiór wartości \(\mathop{\mathrm{Ran}}(f)\) funkcji \(f(x) = \frac{x-3} 
{x+1}\).}
{\begin{align*}
\mathrm{Dom}(f) &= (-\infty ;-1)\cup (-1;\infty ),\\
\mathop{\mathrm{Ran}}(f) &= (-\infty ;1)\cup (1;\infty )
\end{align*}}{\begin{align*}
\mathrm{Dom}(f) &= (-\infty ;1)\cup (1;\infty ),\\
\mathop{\mathrm{Ran}}(f) &= (-\infty ;-1)\cup (-1;\infty )
\end{align*}}{\begin{align*}
\mathrm{Dom}(f) &= (-\infty ;3)\cup (3;\infty ),\\
\mathop{\mathrm{Ran}}(f) &= (-\infty ;-1)\cup (-1;\infty )
\end{align*}}{\begin{align*}
\mathrm{Dom}(f) &= (-\infty ;-3)\cup (-3;\infty ),\\
\mathop{\mathrm{Ran}}(f) &= (-\infty ;1)\cup (1;\infty )
\end{align*}}{}{}}
\LANGcs{
\Question{
Určete definiční obor \(\mathrm{D}(f)\) 
a obor hodnot \(\mathop{\mathrm{H}}(f)\) funkce \(f(x) = \frac{x-3} 
{x+1}\).}
{\begin{align*}
\mathrm{D}(f) &= (-\infty ;-1)\cup (-1;\infty ),\\
\mathop{\mathrm{H}}(f) &= (-\infty ;1)\cup (1;\infty )
\end{align*}}{\begin{align*}
\mathrm{D}(f) &= (-\infty ;1)\cup (1;\infty ),\\
\mathop{\mathrm{H}}(f) &= (-\infty ;-1)\cup (-1;\infty )
\end{align*}}{\begin{align*}
\mathrm{D}(f) &= (-\infty ;3)\cup (3;\infty ),\\
\mathop{\mathrm{H}}(f) &= (-\infty ;-1)\cup (-1;\infty )
\end{align*}}{\begin{align*}
\mathrm{D}(f) &= (-\infty ;-3)\cup (-3;\infty ),\\
\mathop{\mathrm{H}}(f) &= (-\infty ;1)\cup (1;\infty )
\end{align*}}{}{}}
\LANGsk{
\Question{
Určte definičný obor \(\mathrm{D}(f)\) 
a obor hodnôt \(\mathop{\mathrm{H}}(f)\) funkce \(f(x) = \frac{x-3} 
{x+1}\).}
{\begin{align*}
\mathrm{D}(f) &= (-\infty ;-1)\cup (-1;\infty ),\\
\mathop{\mathrm{H}}(f) &= (-\infty ;1)\cup (1;\infty )
\end{align*}}{\begin{align*}
\mathrm{D}(f) &= (-\infty ;1)\cup (1;\infty ),\\
\mathop{\mathrm{H}}(f) &= (-\infty ;-1)\cup (-1;\infty )
\end{align*}}{\begin{align*}
\mathrm{D}(f) &= (-\infty ;3)\cup (3;\infty ),\\
\mathop{\mathrm{H}}(f) &= (-\infty ;-1)\cup (-1;\infty )
\end{align*}}{\begin{align*}
\mathrm{D}(f) &= (-\infty ;-3)\cup (-3;\infty ),\\
\mathop{\mathrm{H}}(f) &= (-\infty ;1)\cup (1;\infty )
\end{align*}}{}{}}
\LANGes{
\Question{
Determina el dominio \(\mathrm{Dom}(f)\) 
y el rango \(\mathop{\mathrm{Ran}}(f)\)
de la función \(f(x) = \frac{x-3} 
{x+1}\).}
{\begin{align*}
\mathrm{Dom}(f) &= (-\infty ;-1)\cup (-1;\infty ),\\
\mathop{\mathrm{Ran}}(f) &= (-\infty ;1)\cup (1;\infty )
\end{align*}}{\begin{align*}
\mathrm{Dom}(f) &= (-\infty ;1)\cup (1;\infty ),\\
\mathop{\mathrm{Ran}}(f) &= (-\infty ;-1)\cup (-1;\infty )
\end{align*}}{\begin{align*}
\mathrm{Dom}(f) &= (-\infty ;3)\cup (3;\infty ),\\
\mathop{\mathrm{Ran}}(f) &= (-\infty ;-1)\cup (-1;\infty )
\end{align*}}{\begin{align*}
\mathrm{Dom}(f) &= (-\infty ;-3)\cup (-3;\infty ),\\
\mathop{\mathrm{Ran}}(f) &= (-\infty ;1)\cup (1;\infty )
\end{align*}}{}{}}
\End

\Data\ID{2010009902}
\Updated{2022-08-25 10:08:18}
\Level{B}
\SubArea{Rational functions}
\LANGen{
\Question{
Consider the function \(f(x) = \frac{-1} 
 {x+2}-1 \).
Find all \(x\) 
such that \(f(x) > 0\).}
{\(x\in (-3;-2)\)}{\(x\in (-2;3)\)}{\(x\in \left (-\infty ;-3\right )\cup (-2;\infty )\)}{\(x\in \left (-\infty ;-2\right )\cup (3;\infty )\)}{}{}}
\LANGpl{
\Question{
Dana jest funckja \(f(x) = \frac{-1} 
 {x+2}-1 \).
Znajdź wszystkie \(x\), 
dla których \(f(x) > 0\).}
{\(x\in (-3;-2)\)}{\(x\in (-2;3)\)}{\(x\in \left (-\infty ;-3\right )\cup (-2;\infty )\)}{\(x\in \left (-\infty ;-2\right )\cup (3;\infty )\)}{}{}}
\LANGcs{
\Question{
Je dána funkce \(f(x) = \frac{-1} 
 {x+2}-1 \).
Určete všechna taková \(x\), 
pro která platí, že \(f(x) > 0\).}
{\(x\in (-3;-2)\)}{\(x\in (-2;3)\)}{\(x\in \left (-\infty ;-3\right )\cup (-2;\infty )\)}{\(x\in \left (-\infty ;-2\right )\cup (3;\infty )\)}{}{}}
\LANGsk{
\Question{
Je daná funkcia \(f(x) = \frac{-1} 
 {x+2}-1 \).
Určte všetky také \(x\), 
pre ktoré platí, že \(f(x) > 0\).}
{\(x\in (-3;-2)\)}{\(x\in (-2;3)\)}{\(x\in \left (-\infty ;-3\right )\cup (-2;\infty )\)}{\(x\in \left (-\infty ;-2\right )\cup (3;\infty )\)}{}{}}
\LANGes{
\Question{
Considera la función \(f(x) = \frac{-1} 
 {x+2}-1 \).
Determina todos los \(x\)  tales que \(f(x) > 0\).}
{\(x\in (-3;-2)\)}{\(x\in (-2;3)\)}{\(x\in \left (-\infty ;-3\right )\cup (-2;\infty )\)}{\(x\in \left (-\infty ;-2\right )\cup (3;\infty )\)}{}{}}
\End

\Data\ID{2010009903}
\Updated{2022-08-25 10:08:18}
\Level{B}
\SubArea{Rational functions}
\LANGen{
\Question{
Consider the function \(f(x) = \frac{6} 
 {x-1}-1 \).
Find all \(x\) 
such that \(f(x) < 0\).}
{\(x\in \left (-\infty ;1\right )\cup (7;\infty )\)}{\(x\in \left (-\infty ;-7\right )\cup (-1;\infty )\)}{\(x\in (7;\infty)\)}{\(x\in (-\infty;7)\)}{}{}}
\LANGpl{
\Question{
Dana jest funkcja \(f(x) = \frac{6} 
 {x-1}-1 \).
Znajdź wszystkie \(x\),
dla których \(f(x) < 0\).}
{\(x\in \left (-\infty ;1\right )\cup (7;\infty )\)}{\(x\in \left (-\infty ;-7\right )\cup (-1;\infty )\)}{\(x\in (7;\infty)\)}{\(x\in (-\infty;7)\)}{}{}}
\LANGcs{
\Question{
Je dána funkce \(f(x) = \frac{6} 
 {x-1}-1 \).
Určete všechna taková \(x\), 
pro která platí, že \(f(x) < 0\).}
{\(x\in \left (-\infty ;1\right )\cup (7;\infty )\)}{\(x\in \left (-\infty ;-7\right )\cup (-1;\infty )\)}{\(x\in (7;\infty)\)}{\(x\in (-\infty;7)\)}{}{}}
\LANGsk{
\Question{
Je daná funkcia \(f(x) = \frac{6} 
 {x-1}-1 \).
Nájdite všetky také \(x\) 
pre ktoré platí \(f(x) < 0\).}
{\(x\in \left (-\infty ;1\right )\cup (7;\infty )\)}{\(x\in \left (-\infty ;-7\right )\cup (-1;\infty )\)}{\(x\in (7;\infty)\)}{\(x\in (-\infty;7)\)}{}{}}
\LANGes{
\Question{
Considera la función \(f(x) = \frac{6} 
 {x-1}-1 \).
Determina todos los \(x\) tales que 
\(f(x) < 0\).}
{\(x\in \left (-\infty ;1\right )\cup (7;\infty )\)}{\(x\in \left (-\infty ;-7\right )\cup (-1;\infty )\)}{\(x\in (7;\infty)\)}{\(x\in (-\infty;7)\)}{}{}}
\End

\Data\ID{2010015101}
\Updated{2022-08-25 10:08:01}
\Level{B}
\SubArea{Rational functions}
\LANGen{
\Question{
Let by \(X\) and \(Y\) denote the intersection points of the graph of the function \(f(x)=\frac{2}{x+3}-1\) with \(x\) and \(y\)-axis, respectively. Find coordinates of \(X\) and \(Y\).}
{\(X = [-1;0]\),
\(Y = \left[0;-\frac13\right]\)}{\(X = [1;0]\),
\(Y = \left[0;\frac13\right]\)}{\(X = \left[-\frac13;0\right]\),
\(Y = [0;-1]\)}{\(X = [-3;0]\),
\(Y = [0;-1]\)}{}{}}
\LANGpl{
\Question{
Punkty \(X\) i \(Y\) oznaczają punkty przecięcia wykresu funkcji \(f(x)=\frac{2}{x+3}-1\) odpowiednio z osią \(x\) i \(y\). Znajdź współrzędne punktów \(X\) i \(Y\).}
{\(X = [-1;0]\),
\(Y = \left[0;-\frac13\right]\)}{\(X = [1;0]\),
\(Y = \left[0;\frac13\right]\)}{\(X = \left[-\frac13;0\right]\),
\(Y = [0;-1]\)}{\(X = [-3;0]\),
\(Y = [0;-1]\)}{}{}}
\LANGcs{
\Question{
Označme po řadě \(X\) a \(Y\) průsečíky grafu funkce \(f(x)=\frac{2}{x+3}-1\) se souřadnými osami \(x\) a \(y\). Určete souřadnice bodů \(X\) a \(Y\).}
{\(X = [-1;0]\),
\(Y = \left[0;-\frac13\right]\)}{\(X = [1;0]\),
\(Y = \left[0;\frac13\right]\)}{\(X = \left[-\frac13;0\right]\),
\(Y = [0;-1]\)}{\(X = [-3;0]\),
\(Y = [0;-1]\)}{}{}}
\LANGsk{
\Question{
Označme po poradí \(X\) a \(Y\) priesečníky grafu funkcie \(f(x)=\frac{2}{x+3}-1\) so súradnicovými osami \(x\) a \(y\). Určte súradnice bodov \(X\) a \(Y\).}
{\(X = [-1;0]\),
\(Y = \left[0;-\frac13\right]\)}{\(X = [1;0]\),
\(Y = \left[0;\frac13\right]\)}{\(X = \left[-\frac13;0\right]\),
\(Y = [0;-1]\)}{\(X = [-3;0]\),
\(Y = [0;-1]\)}{}{}}
\LANGes{
\Question{
Sean \(X\) e \(Y\) las intersecciones de la gráfica de la función \(f(x)=\frac{2}{x+3}-1\) con los ejes \(x\) e \(y\) respectivamente. Calcula dichos puntos de intersección.}
{\(X = [-1;0]\),
\(Y = \left[0;-\frac13\right]\)}{\(X = [1;0]\),
\(Y = \left[0;\frac13\right]\)}{\(X = \left[-\frac13;0\right]\),
\(Y = [0;-1]\)}{\(X = [-3;0]\),
\(Y = [0;-1]\)}{}{}}
\End

\Data\ID{2010015102}
\Updated{2022-08-25 10:08:01}
\Level{B}
\SubArea{Rational functions}
\IMG{2010015101-figure0.pdf}
\LANGen{
\Question{
Identify the function that is graphed in the picture.}
{\(f(x) = \frac{-x-4} 
{x+3}\)}{\(f(x) = \frac{-x+3} 
{x-4}\)}{\(f(x) = \frac{-x+3} 
{x+1}\)}{\(f(x) = \frac{-x-1} 
{x+3}\)}{}{}}
\LANGpl{
\Question{
Wskaż funkcję przedstawioną na wykresie.}
{\(f(x) = \frac{-x-4} 
{x+3}\)}{\(f(x) = \frac{-x+3} 
{x-4}\)}{\(f(x) = \frac{-x+3} 
{x+1}\)}{\(f(x) = \frac{-x-1} 
{x+3}\)}{}{}}
\LANGcs{
\Question{
K danému grafu funkce přiřaďte správný funkční předpis.}
{\(f(x) = \frac{-x-4} 
{x+3}\)}{\(f(x) = \frac{-x+3} 
{x-4}\)}{\(f(x) = \frac{-x+3} 
{x+1}\)}{\(f(x) = \frac{-x-1} 
{x+3}\)}{}{}}
\LANGsk{
\Question{
K danému grafu funkcie priraďte správny funkčný predpis.}
{\(f(x) = \frac{-x-4} 
{x+3}\)}{\(f(x) = \frac{-x+3} 
{x-4}\)}{\(f(x) = \frac{-x+3} 
{x+1}\)}{\(f(x) = \frac{-x-1} 
{x+3}\)}{}{}}
\LANGes{
\Question{
Identifica una posible expresión analítica para la gráfica de la función que está en la imagen.}
{\(f(x) = \frac{-x-4} 
{x+3}\)}{\(f(x) = \frac{-x+3} 
{x-4}\)}{\(f(x) = \frac{-x+3} 
{x+1}\)}{\(f(x) = \frac{-x-1} 
{x+3}\)}{}{}}
\End

\Data\ID{2010015103}
\Updated{2022-08-25 10:08:01}
\Level{B}
\SubArea{Rational functions}
\IMG{2010015101-figure1.pdf}
\LANGen{
\Question{
Identify the function that is graphed in the picture.}
{\(f(x) = \frac{1-2x} 
{x-1}\)}{\(f(x) = \frac{1+2x} 
{x+1}\)}{\(f(x) = \frac{x-2} 
{x-1}\)}{\(f(x) = \frac{x+2} 
{x-1}\)}{}{}}
\LANGpl{
\Question{
Wskaż funkcję przedstawioną na wykresie.}
{\(f(x) = \frac{1-2x} 
{x-1}\)}{\(f(x) = \frac{1+2x} 
{x+1}\)}{\(f(x) = \frac{x-2} 
{x-1}\)}{\(f(x) = \frac{x+2} 
{x-1}\)}{}{}}
\LANGcs{
\Question{
K danému grafu funkce přiřaďte správný funkční předpis.}
{\(f(x) = \frac{1-2x} 
{x-1}\)}{\(f(x) = \frac{1+2x} 
{x+1}\)}{\(f(x) = \frac{x-2} 
{x-1}\)}{\(f(x) = \frac{x+2} 
{x-1}\)}{}{}}
\LANGsk{
\Question{
K danému grafu funkcie priraďte správny funkčný predpis.}
{\(f(x) = \frac{1-2x} 
{x-1}\)}{\(f(x) = \frac{1+2x} 
{x+1}\)}{\(f(x) = \frac{x-2} 
{x-1}\)}{\(f(x) = \frac{x+2} 
{x-1}\)}{}{}}
\LANGes{
\Question{
Identifica una posible expresión analítica para la gráfica de la función que está en la imagen.}
{\(f(x) = \frac{1-2x} 
{x-1}\)}{\(f(x) = \frac{1+2x} 
{x+1}\)}{\(f(x) = \frac{x-2} 
{x-1}\)}{\(f(x) = \frac{x+2} 
{x-1}\)}{}{}}
\End

\Data\ID{2010015104}
\Updated{2022-08-25 10:08:01}
\Level{B}
\SubArea{Rational functions}
\LANGen{
\Question{
Find intersection points of the graph of the rational function \(f(x)=\frac{2x-6}{x+3}\) with \(y\)-axis.}
{\(Y = \left [0; -2\right ]\)}{\(Y = \left [3; 0\right ]\)}{\(Y = \left [0; -3\right ]\)}{\(Y = \left [3; -3\right ]\)}{}{}}
\LANGpl{
\Question{
Znajdź punkty przecięcia wykresu funkcji wymiernej \(f(x)=\frac{2x-6}{x+3}\) z osią \(y\).}
{\(Y = \left [0; -2\right ]\)}{\(Y = \left [3; 0\right ]\)}{\(Y = \left [0; -3\right ]\)}{\(Y = \left [3; -3\right ]\)}{}{}}
\LANGcs{
\Question{
Určete společné body grafu lineární lomené funkce  \(f(x)=\frac{2x-6}{x+3}\) s osou \(y\).}
{\(Y = \left [0; -2\right ]\)}{\(Y = \left [3; 0\right ]\)}{\(Y = \left [0; -3\right ]\)}{\(Y = \left [3; -3\right ]\)}{}{}}
\LANGsk{
\Question{
Určte spoločné body grafu lineárnej lomenej funkcie  \(f(x)=\frac{2x-6}{x+3}\) s osou \(y\).}
{\(Y = \left [0; -2\right ]\)}{\(Y = \left [3; 0\right ]\)}{\(Y = \left [0; -3\right ]\)}{\(Y = \left [3; -3\right ]\)}{}{}}
\LANGes{
\Question{
Determina los puntos de intersección de la gráfica de la función racional \(f(x)=\frac{2x-6}{x+3}\) con el eje \(y\).}
{\(Y = \left [0; -2\right ]\)}{\(Y = \left [3; 0\right ]\)}{\(Y = \left [0; -3\right ]\)}{\(Y = \left [3; -3\right ]\)}{}{}}
\End

\Data\ID{2010015105}
\Updated{2022-08-25 10:08:01}
\Level{B}
\SubArea{Rational functions}
\LANGen{
\Question{
Find the intersection points of the graph of the rational function \(f(x)=\frac{2x-6}{x+3}\) with \(x\)-axis.}
{\(X = \left [3; 0\right ]\)}{\(X = \left [0; -2\right ]\)}{\(X = \left [0; 3\right ]\)}{\(X = \left [3; -3\right ]\)}{}{}}
\LANGpl{
\Question{
Znajdź punkty przecięcia wykresu funkcji wymiernej \(f(x)=\frac{2x-6}{x+3}\) z osią \(x\).}
{\(X = \left [3; 0\right ]\)}{\(X = \left [0; -2\right ]\)}{\(X = \left [0; 3\right ]\)}{\(X = \left [3; -3\right ]\)}{}{}}
\LANGcs{
\Question{
Určete společné body grafu lineární lomené funkce \(f(x)=\frac{2x-6}{x+3}\) s osou \(x\).}
{\(X = \left [3; 0\right ]\)}{\(X = \left [0; -2\right ]\)}{\(X = \left [0; 3\right ]\)}{\(X = \left [3; -3\right ]\)}{}{}}
\LANGsk{
\Question{
Určte spoločné body grafu lineárnej lomenej funkcie \(f(x)=\frac{2x-6}{x+3}\) s osou \(x\).}
{\(X = \left [3; 0\right ]\)}{\(X = \left [0; -2\right ]\)}{\(X = \left [0; 3\right ]\)}{\(X = \left [3; -3\right ]\)}{}{}}
\LANGes{
\Question{
Determina los puntos de intersección de la gráfica de la función racional \(f(x)=\frac{2x-6}{x+3}\) con el eje \(x\).}
{\(X = \left [3; 0\right ]\)}{\(X = \left [0; -2\right ]\)}{\(X = \left [0; 3\right ]\)}{\(X = \left [3; -3\right ]\)}{}{}}
\End

\Data\ID{9000002901}
\Updated{2020-12-03 06:12:36}
\Level{B}
\SubArea{Rational functions}
\LANGen{
\Question{
Find the domain of the function \(f\colon y = \frac{1} 
{x-2} + 1\).}
{\(\mathbb{R}\setminus \{2\}\)}{\(\mathbb{R}\setminus \{ - 1\}\)}{\(\mathbb{R}\setminus \{0\}\)}{\(\mathbb{R}\)}{}{}}
\LANGpl{
\Question{
Znajdź dziedzinę funkcji  \(f\colon y = \frac{1} 
{x-2} + 1\).}
{\(\mathbb{R}\setminus \{2\}\)}{\(\mathbb{R}\setminus \{ - 1\}\)}{\(\mathbb{R}\setminus \{0\}\)}{\(\mathbb{R}\)}{}{}}
\LANGcs{
\Question{
Určete definiční obor funkce \(f\colon y = \frac{1} 
{x-2} + 1\).}
{\(\mathbb{R}\setminus \{2\}\)}{\(\mathbb{R}\setminus \{ - 1\}\)}{\(\mathbb{R}\setminus \{0\}\)}{\(\mathbb{R}\)}{}{}}
\LANGsk{
\Question{
Definičným oborom funkcie  \(f\colon y = \frac{1} 
{x-2} + 1\) je:}
{\(\mathbb{R}\setminus \{2\}\)}{\(\mathbb{R}\setminus \{ - 1\}\)}{\(\mathbb{R}\setminus \{0\}\)}{\(\mathbb{R}\)}{}{}}
\LANGes{
\Question{
Determina el dominio de la función \(f(x) = \frac{1} 
{x-2} + 1\).}
{\(\mathbb{R}\setminus \{2\}\)}{\(\mathbb{R}\setminus \{ - 1\}\)}{\(\mathbb{R}\setminus \{0\}\)}{\(\mathbb{R}\)}{}{}}
\End

\Data\ID{9000002902}
\Updated{2020-07-15 03:07:34}
\Level{B}
\SubArea{Rational functions}
\LANGen{
\Question{
The point \(A\) is on the
graph of the function \(f(x)= \frac{-10} 
{x+5}\).
The \(x\)-coordinate
of the point \(A\) 
is \(10\). Find the
\(y\)-coordinate
of the point \(A\).}
{\(-\frac{2} 
{3}\)}{\(\frac{2}
{3}\)}{\(-\frac{1} 
{5}\)}{\(-\frac{1} 
{3}\)}{}{}}
\LANGpl{
\Question{
Punkt  \(A\) należy do wykresu funkcji  \(f\colon y = \frac{-10} 
{x+5}\).
Współrzędna \(x\)
punktu  \(A\) 
to \(10\). Znajdź współrzędną 
\(y\)
punktu  \(A\).}
{\(-\frac{2} 
{3}\)}{\(\frac{2}
{3}\)}{\(-\frac{1} 
{5}\)}{\(-\frac{1} 
{3}\)}{}{}}
\LANGcs{
\Question{
Určete chybějící souřadnici bodu
\(A = [10;?]\) tak, aby bod
\(A\) ležel na grafu funkce
dané předpisem \(f\colon y = \frac{-10} 
{x+5}\).}
{\(-\frac{2} 
{3}\)}{\(\frac{2}
{3}\)}{\(-\frac{1} 
{5}\)}{\(-\frac{1} 
{3}\)}{}{}}
\LANGsk{
\Question{
Určte chýbajúcu súradnicu bodu
\(A = [10;?]\) tak, aby bod
\(A\) patril grafu funkcie
danej predpisom \(f\colon y = \frac{-10} 
{x+5}\).}
{\(-\frac{2} 
{3}\)}{\(\frac{2}
{3}\)}{\(-\frac{1} 
{5}\)}{\(-\frac{1} 
{3}\)}{}{}}
\LANGes{
\Question{
El punto  \(A\) está en la
gráfica de la función  \(f(x)= \frac{-10} 
{x+5}\).
La coordenada \(x\) del punto \(A\) 
es \(10\). Calcula la coordenada
\(y\) del punto \(A\).}
{\(-\frac{2} 
{3}\)}{\(\frac{2}
{3}\)}{\(-\frac{1} 
{5}\)}{\(-\frac{1} 
{3}\)}{}{}}
\End

\Data\ID{9000002903}
\Updated{2021-04-09 07:04:55}
\Level{B}
\SubArea{Rational functions}
\LANGen{
\Question{
In the following list identify a point which is on the graph of the function
\(f(x) = \frac{3} 
{x} - 5\).}
{\(A = \left [-6;-\frac{11}
 {2} \right ]\)}{\(A = \left [-1;-2\right ]\)}{\(A = \left [-3;-\frac{5}
{2}\right ]\)}{\(A = \left [\frac{1}
{2};-1\right ]\)}{}{}}
\LANGpl{
\Question{
Z poniższej listy wybierz punkt, który należy do wykresu funkcji
\(f\colon y = \frac{3} 
{x} - 5\).}
{\(A = \left [-6;-\frac{11}
 {2} \right ]\)}{\(A = \left [-1;-2\right ]\)}{\(A = \left [-3;-\frac{5}
{2}\right ]\)}{\(A = \left [\frac{1}
{2};-1\right ]\)}{}{}}
\LANGcs{
\Question{
Který z bodů je bodem grafu funkce dané předpisem
\(f\colon y = \frac{3} 
{x} - 5\)?}
{\(A = \left [-6;-\frac{11}
 {2} \right ]\)}{\(A = \left [-1;-2\right ]\)}{\(A = \left [-3;-\frac{5}
{2}\right ]\)}{\(A = \left [\frac{1}
{2};-1\right ]\)}{}{}}
\LANGsk{
\Question{
Vyberte bod, ktorý patrí grafu funkcie 
\(f\colon y = \frac{3} 
{x} - 5\).}
{\(A = \left [-6;-\frac{11}
 {2} \right ]\)}{\(A = \left [-1;-2\right ]\)}{\(A = \left [-3;-\frac{5}
{2}\right ]\)}{\(A = \left [\frac{1}
{2};-1\right ]\)}{}{}}
\LANGes{
\Question{
En la siguiente lista, identifica un punto que está en la gráfica de la función
\(f(x) = \frac{3} 
{x} - 5\).}
{\(A = \left [-6;-\frac{11}
 {2} \right ]\)}{\(A = \left [-1;-2\right ]\)}{\(A = \left [-3;-\frac{5}
{2}\right ]\)}{\(A = \left [\frac{1}
{2};-1\right ]\)}{}{}}
\End

\Data\ID{9000002904}
\Updated{2022-04-23 05:04:00}
\Level{B}
\SubArea{Rational functions}
\LANGen{
\Question{
Let by \(X\) and \(Y\) denote the intersection points of the graph of the function \(f(x) =  - \frac{1}
{x-1} + 1 \) 
 with \(x\) and \(y\)-axis, respectively. Find coordinates of \(X\) and \(Y\).}
{\(X = [2;0]\),
\(Y = [0;2]\)}{\(X = [1;0]\),
\(Y = [0;1]\)}{\(X = [0;2]\),
\(Y = [2;0]\)}{\(X = Y = [0;0]\)}{}{}}
\LANGpl{
\Question{
Niech  \(X\) i 
\(Y \) punktami przecięcia wykresu funkcji  \(f\colon y = - \frac{1}
{x-1} + 1\) 
odpowiednio z osia współrzędnych  \(x\)-
i \(y\). Wyznacz te punkty.}
{\(X = [2;0]\),
\(Y = [0;2]\)}{\(X = [1;0]\),
\(Y = [0;1]\)}{\(X = [0;2]\),
\(Y = [2;0]\)}{\(X = Y = [0;0]\)}{}{}}
\LANGcs{
\Question{
Je dána funkce \(f\colon y = - \frac{1}
{x-1} + 1\).
Průsečíky grafu funkce \(f\) 
s osami \(x\),
\(y\) označme
po řadě \(X\),
\(Y \). Určete
souřadnice bodů \(X\) 
a \(Y \).}
{\(X = [2;0]\),
\(Y = [0;2]\)}{\(X = [1;0]\),
\(Y = [0;1]\)}{\(X = [0;2]\),
\(Y = [2;0]\)}{\(X = Y = [0;0]\)}{}{}}
\LANGsk{
\Question{
Daná je funkcia \(f\colon y = - \frac{1}
{x-1} + 1\).
Priesečníky grafu funkcie \(f\) 
s osami \(x\),
\(y\) označme
v tomto poradí \(X\),
\(Y \). Určte
súradnice bodov \(X\) 
a \(Y \).}
{\(X = [2;0]\),
\(Y = [0;2]\)}{\(X = [1;0]\),
\(Y = [0;1]\)}{\(X = [0;2]\),
\(Y = [2;0]\)}{\(X = Y = [0;0]\)}{}{}}
\LANGes{
\Question{
Sean \(X\) e
\(Y \) las intersecciones de la
gráfica de la función \(f(x) = - \frac{1}
{x-1} + 1\) 
con los ejes \(x\)
e \(y\)
respectivamente. Calcula dichos puntos de intersección.}
{\(X = [2;0]\),
\(Y = [0;2]\)}{\(X = [1;0]\),
\(Y = [0;1]\)}{\(X = [0;2]\),
\(Y = [2;0]\)}{\(X = Y = [0;0]\)}{}{}}
\End

\Data\ID{9000002905}
\Updated{2022-05-29 12:05:06}
\Level{B}
\SubArea{Rational functions}
\LANGen{
\Question{
Find the range of the function \(f(x)= \frac{1} 
{x-2} + 1\).}
{\((-\infty ;1)\cup (1;\infty )\)}{\(\mathbb{R}\)}{\((-\infty ;2)\cup (2;\infty )\)}{\((-\infty ;-1)\cup (-1;\infty )\)}{}{}}
\LANGpl{
\Question{
Określ zbiór wartości funkcji  \(f\colon y = \frac{1} 
{x-2} + 1\).}
{\((-\infty ;1)\cup (1;\infty )\)}{\(\mathbb{R}\)}{\((-\infty ;2)\cup (2;\infty )\)}{\((-\infty ;-1)\cup (-1;\infty )\)}{}{}}
\LANGcs{
\Question{
Je dána funkce \(f\colon y = \frac{1} 
{x-2} + 1\).
Určete její obor hodnot.}
{\((-\infty ;1)\cup (1;\infty )\)}{\(\mathbb{R}\)}{\((-\infty ;2)\cup (2;\infty )\)}{\((-\infty ;-1)\cup (-1;\infty )\)}{}{}}
\LANGsk{
\Question{
Určte obor hodnôt danej funkcie \(f\colon y = \frac{1} 
{x-2} + 1\).}
{\((-\infty ;1)\cup (1;\infty )\)}{\(\mathbb{R}\)}{\((-\infty ;2)\cup (2;\infty )\)}{\((-\infty ;-1)\cup (-1;\infty )\)}{}{}}
\LANGes{
\Question{
Determina el rango de la función \(f(x)= \frac{1} 
{x-2} + 1\).}
{\((-\infty ;1)\cup (1;\infty )\)}{\(\mathbb{R}\)}{\((-\infty ;2)\cup (2;\infty )\)}{\((-\infty ;-1)\cup (-1;\infty )\)}{}{}}
\End

\Data\ID{9000002906}
\Updated{2021-04-09 07:04:12}
\Level{B}
\SubArea{Rational functions}
\LANGen{
\Question{
Find the domain of the function
\(f(x) = - \frac{3}
{x-1} - 2\) if we have to ensure
that the range of \(f\) 
is \((-1;1] \).}
{\((-2;0] \)}{\([ - 2;0)\)}{\((0;2] \)}{\((0;4)\)}{}{}}
\LANGpl{
\Question{
Wyznacz dziedzinę funkcji
\(f\colon y = - \frac{3}
{x-1} - 2\) jeśli pewne jest, że jej zakres to \(f\) 
is \((-1;1] \).}
{\((-2;0] \)}{\([ - 2;0)\)}{\((0;2] \)}{\((0;4)\)}{}{}}
\LANGcs{
\Question{
Určete definiční obor funkce
\(f\colon y = - \frac{3}
{x-1} - 2\), jejímž oborem
hodnot je interval \((-1;1\rangle \).}
{\((-2;0\rangle \)}{\(\langle - 2;0)\)}{\((0;2\rangle \)}{\((0;4)\)}{}{}}
\LANGsk{
\Question{
Určte definičný obor funkcie
\(f\colon y = - \frac{3}
{x-1} - 2\), ktorej obor hodnôt je interval \((-1;1\rangle \).}
{\((-2;0\rangle \)}{\(\langle - 2;0)\)}{\((0;2\rangle \)}{\((0;4)\)}{}{}}
\LANGes{
\Question{
Determina el dominio de la función
\(f(x) = - \frac{3}
{x-1} - 2\)  si tenemos que asegurarnos
de que el rango de la función \(f\) 
es \((-1;1] \).}
{\((-2;0] \)}{\([ - 2;0)\)}{\((0;2] \)}{\((0;4)\)}{}{}}
\End

\Data\ID{9000002907}
\Updated{2020-07-15 03:07:34}
\Level{B}
\SubArea{Rational functions}
\IMG{000029_07-figure0.pdf}
\LANGen{
\Question{
Identify a possible analytic expression for the function graphed in the
picture.}
{\(y = -2 + \frac{1} 
{x+1}\)}{\(y = - \frac{1}
{x+1} - 2\)}{\(y = \frac{1} 
{x+2} - 1\)}{\(y = 2 + \frac{1} 
{x+1}\)}{}{}}
\LANGpl{
\Question{
Określ prawdopodobny wzór funkcji, której wykres przedstawiony jest na rysunku poniżej.}
{\(y = -2 + \frac{1} 
{x+1}\)}{\(y = - \frac{1}
{x+1} - 2\)}{\(y = \frac{1} 
{x+2} - 1\)}{\(y = 2 + \frac{1} 
{x+1}\)}{}{}}
\LANGcs{
\Question{
Určete předpis lineární lomené funkce, jejíž graf je na obrázku.}
{\(y = -2 + \frac{1} 
{x+1}\)}{\(y = - \frac{1}
{x+1} - 2\)}{\(y = \frac{1} 
{x+2} - 1\)}{\(y = 2 + \frac{1} 
{x+1}\)}{}{}}
\LANGsk{
\Question{
Určte predpis lineárnej lomenej funkcie, ktorej graf je znázornený na obrázku.}
{\(y = -2 + \frac{1} 
{x+1}\)}{\(y = - \frac{1}
{x+1} - 2\)}{\(y = \frac{1} 
{x+2} - 1\)}{\(y = 2 + \frac{1} 
{x+1}\)}{}{}}
\LANGes{
\Question{
Identifica una posible expresión analítica para la función graficada en
la imagen.}
{\(y = -2 + \frac{1} 
{x+1}\)}{\(y = - \frac{1}
{x+1} - 2\)}{\(y = \frac{1} 
{x+2} - 1\)}{\(y = 2 + \frac{1} 
{x+1}\)}{}{}}
\End

\Data\ID{9000003103}
\Updated{2021-04-09 07:04:59}
\Level{B}
\SubArea{Rational functions}
\IMG{000031_03-figure0.pdf}
\LANGen{
\Question{
Identify a possible analytic expression for the function graphed in the picture.}
{\(y = 1 -\frac{2} 
{x}\)}{\(y = -1 + \frac{2} 
{x}\)}{\(y = 1 + \frac{2} 
{x}\)}{\(y = -1 -\frac{2} 
{x}\)}{}{}}
\LANGpl{
\Question{
Określ prawdopodobny wzór funkcji, której wykres przedstawiony jest na rysunku poniżej.}
{\(y = 1 -\frac{2} 
{x}\)}{\(y = -1 + \frac{2} 
{x}\)}{\(y = 1 + \frac{2} 
{x}\)}{\(y = -1 -\frac{2} 
{x}\)}{}{}}
\LANGcs{
\Question{
K danému grafu funkce přiřaďte správný funkční předpis.}
{\(y = 1 -\frac{2} 
{x}\)}{\(y = -1 + \frac{2} 
{x}\)}{\(y = 1 + \frac{2} 
{x}\)}{\(y = -1 -\frac{2} 
{x}\)}{}{}}
\LANGsk{
\Question{
Predpis funkcie, ktorej graf vidíte na obrázku, je:}
{\(y = 1 -\frac{2} 
{x}\)}{\(y = -1 + \frac{2} 
{x}\)}{\(y = 1 + \frac{2} 
{x}\)}{\(y = -1 -\frac{2} 
{x}\)}{}{}}
\LANGes{
\Question{
Identifica una posible expresión analítica para la gráfica de la función de la imagen.}
{\(y = 1 -\frac{2} 
{x}\)}{\(y = -1 + \frac{2} 
{x}\)}{\(y = 1 + \frac{2} 
{x}\)}{\(y = -1 -\frac{2} 
{x}\)}{}{}}
\End

\Data\ID{9000003104}
\Updated{2021-04-09 07:04:47}
\Level{B}
\SubArea{Rational functions}
\IMG{000031_04-figure0.pdf}
\LANGen{
\Question{
Identify a possible analytic expression for the function graphed in the picture.}
{\(y = -2 -\frac{1} 
{x}\)}{\(y = 2 + \frac{1} 
{x}\)}{\(y = -2 + \frac{1} 
{x}\)}{\(y = 2 -\frac{1} 
{x}\)}{}{}}
\LANGpl{
\Question{
Określ prawdopodobny wzór funkcji, której wykres przedstawiony jest na rysunku poniżej.}
{\(y = -2 -\frac{1} 
{x}\)}{\(y = 2 + \frac{1} 
{x}\)}{\(y = -2 + \frac{1} 
{x}\)}{\(y = 2 -\frac{1} 
{x}\)}{}{}}
\LANGcs{
\Question{
K danému grafu funkce přiřaďte správný funkční předpis.}
{\(y = -2 -\frac{1} 
{x}\)}{\(y = 2 + \frac{1} 
{x}\)}{\(y = -2 + \frac{1} 
{x}\)}{\(y = 2 -\frac{1} 
{x}\)}{}{}}
\LANGsk{
\Question{
Predpis funkcie, ktorej graf vidíte na obrázku, je:}
{\(y = -2 -\frac{1} 
{x}\)}{\(y = 2 + \frac{1} 
{x}\)}{\(y = -2 + \frac{1} 
{x}\)}{\(y = 2 -\frac{1} 
{x}\)}{}{}}
\LANGes{
\Question{
Identifica una posible expresión analítica para la gráfica de la función de la imagen.}
{\(y = -2 -\frac{1} 
{x}\)}{\(y = 2 + \frac{1} 
{x}\)}{\(y = -2 + \frac{1} 
{x}\)}{\(y = 2 -\frac{1} 
{x}\)}{}{}}
\End

\Data\ID{9000003105}
\Updated{2021-04-09 07:04:06}
\Level{B}
\SubArea{Rational functions}
\IMG{000031_05-figure0.pdf}
\LANGen{
\Question{
Identify a possible analytic expression for the function graphed in the picture.}
{\(y = \frac{1} 
{x-2}\)}{\(y = - \frac{1}
{x-2}\)}{\(y = - \frac{1}
{x+2}\)}{\(y = \frac{1} 
{x+2}\)}{}{}}
\LANGpl{
\Question{
Określ prawdopodobny wzór funkcji, której wykres przedstawiony jest na rysunku poniżej.}
{\(y = \frac{1} 
{x-2}\)}{\(y = - \frac{1}
{x-2}\)}{\(y = - \frac{1}
{x+2}\)}{\(y = \frac{1} 
{x+2}\)}{}{}}
\LANGcs{
\Question{
K danému grafu funkce přiřaďte správný funkční předpis.}
{\(y = \frac{1} 
{x-2}\)}{\(y = - \frac{1}
{x-2}\)}{\(y = - \frac{1}
{x+2}\)}{\(y = \frac{1} 
{x+2}\)}{}{}}
\LANGsk{
\Question{
Predpis funkcie, ktorej graf vidíte na obrázku, je:}
{\(y = \frac{1} 
{x-2}\)}{\(y = - \frac{1}
{x-2}\)}{\(y = - \frac{1}
{x+2}\)}{\(y = \frac{1} 
{x+2}\)}{}{}}
\LANGes{
\Question{
Identifica una posible expresión analítica para la gráfica de la función de la imagen.}
{\(y = \frac{1} 
{x-2}\)}{\(y = - \frac{1}
{x-2}\)}{\(y = - \frac{1}
{x+2}\)}{\(y = \frac{1} 
{x+2}\)}{}{}}
\End

\Data\ID{9000003106}
\Updated{2020-07-15 03:07:34}
\Level{B}
\SubArea{Rational functions}
\IMG{000031_06-figure0.pdf}
\LANGen{
\Question{
Identify a possible analytic expression for the function graphed in the picture.}
{\(y = \frac{2} 
{x+1}\)}{\(y = \frac{1} 
{x+2}\)}{\(y = - \frac{1}
{x+2}\)}{\(y = \frac{1} 
{x-1}\)}{}{}}
\LANGpl{
\Question{
Określ prawdopodobny wzór funkcji, której wykres przedstawiony jest na rysunku poniżej.}
{\(y = \frac{2} 
{x+1}\)}{\(y = \frac{1} 
{x+2}\)}{\(y = - \frac{1}
{x+2}\)}{\(y = \frac{1} 
{x-1}\)}{}{}}
\LANGcs{
\Question{
K danému grafu funkce přiřaďte správný funkční předpis.}
{\(y = \frac{2} 
{x+1}\)}{\(y = \frac{1} 
{x+2}\)}{\(y = - \frac{1}
{x+2}\)}{\(y = \frac{1} 
{x-1}\)}{}{}}
\LANGsk{
\Question{
Predpis funkcie, ktorej graf vidíte na obrázku, je:}
{\(y = \frac{2} 
{x+1}\)}{\(y = \frac{1} 
{x+2}\)}{\(y = - \frac{1}
{x+2}\)}{\(y = \frac{1} 
{x-1}\)}{}{}}
\LANGes{
\Question{
Identifica una posible expresión analítica para la gráfica de la función que está en la imagen.}
{\(y = \frac{2} 
{x+1}\)}{\(y = \frac{1} 
{x+2}\)}{\(y = - \frac{1}
{x+2}\)}{\(y = \frac{1} 
{x-1}\)}{}{}}
\End

\Data\ID{9000003107}
\Updated{2020-07-15 03:07:34}
\Level{B}
\SubArea{Rational functions}
\IMG{000031_07-figure0.pdf}
\LANGen{
\Question{
Identify a possible analytic expression for the function graphed in the picture.}
{\(y = -2 + \frac{1} 
{x+1}\)}{\(y = 2 + \frac{1} 
{x+1}\)}{\(y = 2 + \frac{1} 
{x-1}\)}{\(y = -2 + \frac{1} 
{x-1}\)}{}{}}
\LANGpl{
\Question{
Określ prawdopodobny wzór funkcji, której wykres przedstawiony jest na rysunku poniżej.}
{\(y = -2 + \frac{1} 
{x+1}\)}{\(y = 2 + \frac{1} 
{x+1}\)}{\(y = 2 + \frac{1} 
{x-1}\)}{\(y = -2 + \frac{1} 
{x-1}\)}{}{}}
\LANGcs{
\Question{
K danému grafu funkce přiřaďte správný funkční předpis.}
{\(y = -2 + \frac{1} 
{x+1}\)}{\(y = 2 + \frac{1} 
{x+1}\)}{\(y = 2 + \frac{1} 
{x-1}\)}{\(y = -2 + \frac{1} 
{x-1}\)}{}{}}
\LANGsk{
\Question{
Predpis funkcie, ktorej graf vidíte na obrázku, je:}
{\(y = -2 + \frac{1} 
{x+1}\)}{\(y = 2 + \frac{1} 
{x+1}\)}{\(y = 2 + \frac{1} 
{x-1}\)}{\(y = -2 + \frac{1} 
{x-1}\)}{}{}}
\LANGes{
\Question{
Identifica una posible expresión analítica para la gráfica de la función que está en la imagen.}
{\(y = -2 + \frac{1} 
{x+1}\)}{\(y = 2 + \frac{1} 
{x+1}\)}{\(y = 2 + \frac{1} 
{x-1}\)}{\(y = -2 + \frac{1} 
{x-1}\)}{}{}}
\End

\Data\ID{9000003108}
\Updated{2020-07-15 03:07:34}
\Level{B}
\SubArea{Rational functions}
\IMG{000031_08-figure0.pdf}
\LANGen{
\Question{
Identify a possible analytic expression for the function graphed in the picture.}
{\(y = -2 - \frac{1} 
{x-1}\)}{\(y = -1 - \frac{1} 
{x-2}\)}{\(y = -2 + \frac{1} 
{x-1}\)}{\(y = 1 - \frac{1} 
{x-2}\)}{}{}}
\LANGpl{
\Question{
Określ prawdopodobny wzór funkcji, której wykres przedstawiony jest na rysunku poniżej.}
{\(y = -2 - \frac{1} 
{x-1}\)}{\(y = -1 - \frac{1} 
{x-2}\)}{\(y = -2 + \frac{1} 
{x-1}\)}{\(y = 1 - \frac{1} 
{x-2}\)}{}{}}
\LANGcs{
\Question{
K danému grafu funkce přiřaďte správný funkční předpis.}
{\(y = -2 - \frac{1} 
{x-1}\)}{\(y = -1 - \frac{1} 
{x-2}\)}{\(y = -2 + \frac{1} 
{x-1}\)}{\(y = 1 - \frac{1} 
{x-2}\)}{}{}}
\LANGsk{
\Question{
Predpis funkcie, ktorej graf vidíte na obrázku, je:}
{\(y = -2 - \frac{1} 
{x-1}\)}{\(y = -1 - \frac{1} 
{x-2}\)}{\(y = -2 + \frac{1} 
{x-1}\)}{\(y = 1 - \frac{1} 
{x-2}\)}{}{}}
\LANGes{
\Question{
Identifica una posible expresión analítica para la gráfica de la función que está en la imagen.}
{\(y = -2 - \frac{1} 
{x-1}\)}{\(y = -1 - \frac{1} 
{x-2}\)}{\(y = -2 + \frac{1} 
{x-1}\)}{\(y = 1 - \frac{1} 
{x-2}\)}{}{}}
\End

\Data\ID{9000003109}
\Updated{2022-04-23 05:04:05}
\Level{B}
\SubArea{Rational functions}
\IMG{000031_09-figure0.pdf}
\LANGen{
\Question{
Identify the function that is graphed in the picture.}
{\(f(x) = \frac{x+3} 
{x+2}\)}{\(f(x) = \frac{x+2} 
{x+1}\)}{\(f(x) = \frac{x-2} 
{x+1}\)}{\(f(x) = -\frac{x+3}
{x+2}\)}{}{}}
\LANGpl{
\Question{
Określ prawdopodobny wzór funkcji, której wykres przedstawiony jest na rysunku poniżej.}
{\(y= \frac{x+3} 
{x+2}\)}{\(y = \frac{x+2} 
{x+1}\)}{\(y = \frac{x-2} 
{x+1}\)}{\(y = -\frac{x+3}
{x+2}\)}{}{}}
\LANGcs{
\Question{
K danému grafu funkce přiřaďte správný funkční předpis.}
{\(y = \frac{x+3} 
{x+2}\)}{\(y = \frac{x+2} 
{x+1}\)}{\(y = \frac{x-2} 
{x+1}\)}{\(y = -\frac{x+3}
{x+2}\)}{}{}}
\LANGsk{
\Question{
Predpis funkcie, ktorej graf vidíte na obrázku, je:}
{\(y = \frac{x+3} 
{x+2}\)}{\(y = \frac{x+2} 
{x+1}\)}{\(y = \frac{x-2} 
{x+1}\)}{\(y = -\frac{x+3}
{x+2}\)}{}{}}
\LANGes{
\Question{
Identifica una posible expresión analítica para la gráfica de la función que está en la imagen.}
{\(y = \frac{x+3} 
{x+2}\)}{\(y = \frac{x+2} 
{x+1}\)}{\(y = \frac{x-2} 
{x+1}\)}{\(y = -\frac{x+3}
{x+2}\)}{}{}}
\End

\Data\ID{9000003110}
\Updated{2022-04-23 05:04:05}
\Level{B}
\SubArea{Rational functions}
\IMG{000031_10-figure0.pdf}
\LANGen{
\Question{
Identify the function that is graphed in the picture.}
{\(f(x) = \frac{2-x} 
{1-x}\)}{\(f(x) = \frac{x-2} 
{x+1}\)}{\(f(x) = -\frac{2-x}
{1-x}\)}{\(f(x) = \frac{x-1} 
{x+1}\)}{}{}}
\LANGpl{
\Question{
Określ prawdopodobny wzór funkcji, której wykres przedstawiony jest na rysunku poniżej.}
{\(y = \frac{2-x} 
{1-x}\)}{\(y = \frac{x-2} 
{x+1}\)}{\(y = -\frac{2-x}
{1-x}\)}{\(y = \frac{x-1} 
{x+1}\)}{}{}}
\LANGcs{
\Question{
K danému grafu funkce přiřaďte správný funkční předpis.}
{\(y = \frac{2-x} 
{1-x}\)}{\(y = \frac{x-2} 
{x+1}\)}{\(y = -\frac{2-x}
{1-x}\)}{\(y = \frac{x-1} 
{x+1}\)}{}{}}
\LANGsk{
\Question{
Predpis funkcie, ktorej graf vidíte na obrázku, je:}
{\(y = \frac{2-x} 
{1-x}\)}{\(y = \frac{x-2} 
{x+1}\)}{\(y = -\frac{2-x}
{1-x}\)}{\(y = \frac{x-1} 
{x+1}\)}{}{}}
\LANGes{
\Question{
Identifica una posible expresión analítica para la gráfica de la función que está en la imagen.}
{\(y = \frac{2-x} 
{1-x}\)}{\(y = \frac{x-2} 
{x+1}\)}{\(y = -\frac{2-x}
{1-x}\)}{\(y = \frac{x-1} 
{x+1}\)}{}{}}
\End

\Data\ID{9000007501}
\Updated{2021-04-09 07:04:18}
\Level{B}
\SubArea{Rational functions}
\LANGen{
\Question{
The graph of the function 
\[ 
 f(x) = 1 + \frac{3} 
{x + 2}
\]
 is a hyperbola. Find the center of this hyperbola.}
{\(S = [-2;1]\)}{\(S = [3;1]\)}{\(S = [1;3]\)}{\(S = [1;-2]\)}{\(S = [-2;3]\)}{}}
\LANGpl{
\Question{
Wykresem funkcji
\[ 
 f(x) = 1 + \frac{3} 
{x + 2}
\]
 jest hiperbola. Wyznacz współrzędne środka symetrii tej hiperboli.}
{\(S = [-2;1]\)}{\(S = [3;1]\)}{\(S = [1;3]\)}{\(S = [1;-2]\)}{\(S = [-2;3]\)}{}}
\LANGcs{
\Question{
Určete střed hyperboly dané předpisem
\(f(x) = 1 + \frac{3} 
{x+2}\).}
{\(S = [-2;1]\)}{\(S = [3;1]\)}{\(S = [1;3]\)}{\(S = [1;-2]\)}{\(S = [-2;3]\)}{}}
\LANGsk{
\Question{
Určte stred hyperboly danej predpisom
\(f(x) = 1 + \frac{3} 
{x+2}\).}
{\(S = [-2;1]\)}{\(S = [3;1]\)}{\(S = [1;3]\)}{\(S = [1;-2]\)}{\(S = [-2;3]\)}{}}
\LANGes{
\Question{
La gráfica de la función 
\[ 
 f(x) = 1 + \frac{3} 
{x + 2}
\]
es una hipérbola. Determina el centro de la hipérbola.}
{\(S = [-2;1]\)}{\(S = [3;1]\)}{\(S = [1;3]\)}{\(S = [1;-2]\)}{\(S = [-2;3]\)}{}}
\End

\Data\ID{9000007502}
\Updated{2021-04-09 07:04:42}
\Level{B}
\SubArea{Rational functions}
\LANGen{
\Question{
The graph of the function 
\[ 
 f(x) = 2 - \frac{3} 
{x - 2}
\]
 is a hyperbola. Find the center of this hyperbola.}
{\(S = [2;2]\)}{\(S = [-2;2]\)}{\(S = [2;3]\)}{\(S = [-2;3]\)}{\(S = [2;0]\)}{}}
\LANGpl{
\Question{
Wykresem funkcji 
\[ 
 f(x) = 2 - \frac{3} 
{x - 2}
\]
jest hiperbola. Wyznacz środek tej hiperboli.}
{\(S = [2;2]\)}{\(S = [-2;2]\)}{\(S = [2;3]\)}{\(S = [-2;3]\)}{\(S = [2;0]\)}{}}
\LANGcs{
\Question{
Určete střed hyperboly dané předpisem
\(f(x) = 2 - \frac{3} 
{x-2}\).}
{\(S = [2;2]\)}{\(S = [-2;2]\)}{\(S = [2;3]\)}{\(S = [-2;3]\)}{\(S = [2;0]\)}{}}
\LANGsk{
\Question{
Určte stred hyperboly danej predpisom 
\(f(x) = 2 - \frac{3} 
{x-2}\).}
{\(S = [2;2]\)}{\(S = [-2;2]\)}{\(S = [2;3]\)}{\(S = [-2;3]\)}{\(S = [2;0]\)}{}}
\LANGes{
\Question{
La gráfica de la función
\[ 
 f(x) = 2 - \frac{3} 
{x - 2}
\]
es una hipérbola. Determina el centro de la hipérbola.}
{\(S = [2;2]\)}{\(S = [-2;2]\)}{\(S = [2;3]\)}{\(S = [-2;3]\)}{\(S = [2;0]\)}{}}
\End

\Data\ID{9000007503}
\Updated{2021-04-09 07:04:04}
\Level{B}
\SubArea{Rational functions}
\LANGen{
\Question{
The graph of the function 
\[ 
 f(x) = 1 + \frac{1} 
{2(x - 2)}
\]
 is a hyperbola. Find the center of this hyperbola.}
{\(S = [2;1]\)}{\(S = [1;1]\)}{\(S = [1;2]\)}{\(S = [-1;1]\)}{\(S = [2;2]\)}{}}
\LANGpl{
\Question{
Wykresem funkcji 
\[ 
 f(x) = 1 + \frac{1} 
{2(x - 2)}
\]
jest hiperbola. Wyznacz środek tej hiperboli.}
{\(S = [2;1]\)}{\(S = [1;1]\)}{\(S = [1;2]\)}{\(S = [-1;1]\)}{\(S = [2;2]\)}{}}
\LANGcs{
\Question{
Určete střed hyperboly dané předpisem
\(f(x) = 1 + \frac{1} 
{2(x-2)}\).}
{\(S = [2;1]\)}{\(S = [1;1]\)}{\(S = [1;2]\)}{\(S = [-1;1]\)}{\(S = [2;2]\)}{}}
\LANGsk{
\Question{
Určte stred hyperboly danej predpisom 
 \(f\colon y = 1 + \frac{1} 
{2(x-2)}\).}
{\(S = [2;1]\)}{\(S = [1;1]\)}{\(S = [1;2]\)}{\(S = [-1;1]\)}{\(S = [2;2]\)}{}}
\LANGes{
\Question{
La gráfica de la función 
\[ 
 f(x) = 1 + \frac{1} 
{2(x - 2)}
\]
 es una hipérbola. Determina el centro de la hipérbola.}
{\(S = [2;1]\)}{\(S = [1;1]\)}{\(S = [1;2]\)}{\(S = [-1;1]\)}{\(S = [2;2]\)}{}}
\End

\Data\ID{9000007504}
\Updated{2021-04-10 09:04:07}
\Level{B}
\SubArea{Rational functions}
\LANGen{
\Question{
The graph of the function 
\[ 
 f(x) = \frac{3} 
{-2(x + 3)} - 1
\]
 is a hyperbola. Find the center of this hyperbola.}
{\(S = [-3;-1]\)}{\(S = [3;-1]\)}{\(S = [3;1]\)}{\(S = \left [\frac{3}
{2};-1\right ]\)}{\(S = \left [-\frac{3}
{2};-1\right ]\)}{}}
\LANGpl{
\Question{
Wykresem funkcji 
\[ 
 f(x) = \frac{3} 
{-2(x + 3)} - 1
\]
jest hiperbola. Wyznacz środek tej hiperboli.}
{\(S = [-3;-1]\)}{\(S = [3;-1]\)}{\(S = [3;1]\)}{\(S = \left [\frac{3}
{2};-1\right ]\)}{\(S = \left [-\frac{3}
{2};-1\right ]\)}{}}
\LANGcs{
\Question{
Určete střed hyperboly dané předpisem
\(f(x) = \frac{3} 
{-2(x+3)} - 1\).}
{\(S = [-3;-1]\)}{\(S = [3;-1]\)}{\(S = [3;1]\)}{\(S = \left [\frac{3}
{2};-1\right ]\)}{\(S = \left [-\frac{3}
{2};-1\right ]\)}{}}
\LANGsk{
\Question{
Určte stred hyperboly danej predpisom 
 \(f\colon y = \frac{3} 
{-2(x+3)} - 1\).}
{\(S = [-3;-1]\)}{\(S = [3;-1]\)}{\(S = [3;1]\)}{\(S = \left [\frac{3}
{2};-1\right ]\)}{\(S = \left [-\frac{3}
{2};-1\right ]\)}{}}
\LANGes{
\Question{
La gráfica de la función
\[ 
 f(x) = \frac{3} 
{-2(x + 3)} - 1
\]
  es una hipérbola. Determina su centro.}
{\(S = [-3;-1]\)}{\(S = [3;-1]\)}{\(S = [3;1]\)}{\(S = \left [\frac{3}
{2};-1\right ]\)}{\(S = \left [-\frac{3}
{2};-1\right ]\)}{}}
\End

\Data\ID{9000007505}
\Updated{2021-04-10 09:04:25}
\Level{B}
\SubArea{Rational functions}
\LANGen{
\Question{
The graph of the function 
\[ 
 f(x) = \frac{1} 
{-x + 3} + 2
\]
 is a hyperbola. Find the center of this hyperbola.}
{\(S = [3;2]\)}{\(S = [-3;2]\)}{\(S = [1;2]\)}{\(S = [2;3]\)}{\(S = [3;1]\)}{}}
\LANGpl{
\Question{
Wykresem funkcji 
\[ 
 f(x) = \frac{1} 
{-x + 3} + 2
\]
jest hiperbola. Wyznacz środek tej hiperboli.}
{\(S = [3;2]\)}{\(S = [-3;2]\)}{\(S = [1;2]\)}{\(S = [2;3]\)}{\(S = [3;1]\)}{}}
\LANGcs{
\Question{
Určete střed hyperboly dané předpisem
\(f(x) = \frac{1} 
{-x+3} + 2\).}
{\(S = [3;2]\)}{\(S = [-3;2]\)}{\(S = [1;2]\)}{\(S = [2;3]\)}{\(S = [3;1]\)}{}}
\LANGsk{
\Question{
Určte stred hyperboly danej predpisom 
 \(f\colon y = \frac{1} 
{-x+3} + 2\).}
{\(S = [3;2]\)}{\(S = [-3;2]\)}{\(S = [1;2]\)}{\(S = [2;3]\)}{\(S = [3;1]\)}{}}
\LANGes{
\Question{
La gráfica de la función
\[ 
 f(x) = \frac{1} 
{-x + 3} + 2
\]
es una hipérbola.  Determina su centro.}
{\(S = [3;2]\)}{\(S = [-3;2]\)}{\(S = [1;2]\)}{\(S = [2;3]\)}{\(S = [3;1]\)}{}}
\End

\Data\ID{9000007506}
\Updated{2021-04-10 09:04:40}
\Level{B}
\SubArea{Rational functions}
\LANGen{
\Question{
The graph of the function 
\[ 
 f(x) = \frac{3x - 4} 
 {x + 2}
\]
 is a hyperbola. Find the center of this hyperbola.}
{\(S = [-2;3]\)}{\(S = [3;2]\)}{\(S = [0;-4]\)}{\(S = [0;4]\)}{\(S = [4;-2]\)}{}}
\LANGpl{
\Question{
Wykresem funkcji 
\[ 
 f(x) = \frac{3x - 4} 
 {x + 2}
\]
jest hiperbola. Wyznacz środek tej hiperboli.}
{\(S = [-2;3]\)}{\(S = [3;2]\)}{\(S = [0;-4]\)}{\(S = [0;4]\)}{\(S = [4;-2]\)}{}}
\LANGcs{
\Question{
Určete střed hyperboly dané předpisem
\(f(x) = \frac{3x-4} 
 {x+2} \).}
{\(S = [-2;3]\)}{\(S = [3;2]\)}{\(S = [0;-4]\)}{\(S = [0;4]\)}{\(S = [4;-2]\)}{}}
\LANGsk{
\Question{
Určte stred hyperboly danej predpisom 
 \(f\colon y = \frac{3x-4} 
 {x+2} \).}
{\(S = [-2;3]\)}{\(S = [3;2]\)}{\(S = [0;-4]\)}{\(S = [0;4]\)}{\(S = [4;-2]\)}{}}
\LANGes{
\Question{
La gráfica de la función
\[ 
 f(x) = \frac{3x - 4} 
 {x + 2}
\]
es una hipérbola.  Determina su centro.}
{\(S = [-2;3]\)}{\(S = [3;2]\)}{\(S = [0;-4]\)}{\(S = [0;4]\)}{\(S = [4;-2]\)}{}}
\End

\Data\ID{9000007507}
\Updated{2021-04-10 09:04:55}
\Level{B}
\SubArea{Rational functions}
\LANGen{
\Question{
The graph of the function 
\[ 
 f(x) = \frac{2x - 4} 
{3x + 2}
\]
 is a hyperbola. Find the center of this hyperbola.}
{\(S = \left [-\frac{2}
{3}; \frac{2} 
{3}\right ]\)}{\(S = \left [-\frac{3}
{2}; \frac{2} 
{3}\right ]\)}{\(S = \left [\frac{2}
{3};-\frac{3}
{2}\right ]\)}{\(S = \left [-\frac{2}
{3};-\frac{2}
{3}\right ]\)}{\(S = \left [\frac{3}
{2}; \frac{3} 
{2}\right ]\)}{}}
\LANGpl{
\Question{
Wykresem funkcji 
\[ 
 f(x) = \frac{2x - 4} 
{3x + 2}
\]
jest hiperbola. Wyznacz środek tej hiperboli.}
{\(S = \left [-\frac{2}
{3}; \frac{2} 
{3}\right ]\)}{\(S = \left [-\frac{3}
{2}; \frac{2} 
{3}\right ]\)}{\(S = \left [\frac{2}
{3};-\frac{3}
{2}\right ]\)}{\(S = \left [-\frac{2}
{3};-\frac{2}
{3}\right ]\)}{\(S = \left [\frac{3}
{2}; \frac{3} 
{2}\right ]\)}{}}
\LANGcs{
\Question{
Určete střed hyperboly dané předpisem
\(f(x) = \frac{2x-4} 
{3x+2}\).}
{\(S = \left [-\frac{2}
{3}; \frac{2} 
{3}\right ]\)}{\(S = \left [-\frac{3}
{2}; \frac{2} 
{3}\right ]\)}{\(S = \left [\frac{2}
{3};-\frac{3}
{2}\right ]\)}{\(S = \left [-\frac{2}
{3};-\frac{2}
{3}\right ]\)}{\(S = \left [\frac{3}
{2}; \frac{3} 
{2}\right ]\)}{}}
\LANGsk{
\Question{
Určte stred hyperboly danej predpisom 
 \(f\colon y = \frac{2x-4} 
{3x+2}\).}
{\(S = \left [-\frac{2}
{3}; \frac{2} 
{3}\right ]\)}{\(S = \left [-\frac{3}
{2}; \frac{2} 
{3}\right ]\)}{\(S = \left [\frac{2}
{3};-\frac{3}
{2}\right ]\)}{\(S = \left [-\frac{2}
{3};-\frac{2}
{3}\right ]\)}{\(S = \left [\frac{3}
{2}; \frac{3} 
{2}\right ]\)}{}}
\LANGes{
\Question{
La gráfica de la función
\[ 
 f(x) = \frac{2x - 4} 
{3x + 2}
\]
es una hipérbola.  Determina su centro.}
{\(S = \left [-\frac{2}
{3}; \frac{2} 
{3}\right ]\)}{\(S = \left [-\frac{3}
{2}; \frac{2} 
{3}\right ]\)}{\(S = \left [\frac{2}
{3};-\frac{3}
{2}\right ]\)}{\(S = \left [-\frac{2}
{3};-\frac{2}
{3}\right ]\)}{\(S = \left [\frac{3}
{2}; \frac{3} 
{2}\right ]\)}{}}
\End

\Data\ID{9000007508}
\Updated{2021-04-10 09:04:09}
\Level{B}
\SubArea{Rational functions}
\LANGen{
\Question{
The graph of the function 
\[ 
 f(x) = \frac{2x + 1} 
 {x + 2}
\]
 is a hyperbola. Find the center of this hyperbola.}
{\(S = [-2;2]\)}{\(S = [2;-2]\)}{\(S = [2;2]\)}{\(S = [-2;-2]\)}{\(S = [-2;3]\)}{}}
\LANGpl{
\Question{
Wykresem funkcji 
\[ 
 f(x) = \frac{2x + 1} 
 {x + 2}
\]
jest hiperbola. Wyznacz środek tej hiperboli.}
{\(S = [-2;2]\)}{\(S = [2;-2]\)}{\(S = [2;2]\)}{\(S = [-2;-2]\)}{\(S = [-2;3]\)}{}}
\LANGcs{
\Question{
Určete střed hyperboly dané předpisem
\(f(x) = \frac{2x+1} 
 {x+2} \).}
{\(S = [-2;2]\)}{\(S = [2;-2]\)}{\(S = [2;2]\)}{\(S = [-2;-2]\)}{\(S = [-2;3]\)}{}}
\LANGsk{
\Question{
Grafom funkcie
\[ 
 f(x) = \frac{2x + 1} 
 {x + 2}
\]
 je hyperbola. Ktorý z následujúcich bodov je stredom tejto hyperboly?}
{\(S = [-2;2]\)}{\(S = [2;-2]\)}{\(S = [2;2]\)}{\(S = [-2;-2]\)}{\(S = [-2;3]\)}{}}
\LANGes{
\Question{
La gráfica de la función
\[ 
 f(x) = \frac{2x + 1} 
 {x + 2}
\]
es una hipérbola.  Determina su centro.}
{\(S = [-2;2]\)}{\(S = [2;-2]\)}{\(S = [2;2]\)}{\(S = [-2;-2]\)}{\(S = [-2;3]\)}{}}
\End

\Data\ID{9000007509}
\Updated{2021-04-10 09:04:26}
\Level{B}
\SubArea{Rational functions}
\LANGen{
\Question{
The graph of the function 
\[ 
 f(x) = \frac{2x + 3} 
 {2 - x}
\]
 is a hyperbola. Find the center of this hyperbola.}
{\(S = [2;-2]\)}{\(S = [-2;2]\)}{\(S = [2;2]\)}{\(S ={\Bigl [ 2; \frac{3} 
{2}\Bigr ]}\)}{\(S ={\Bigl [ 2;-\frac{3}
{2}\Bigr ]}\)}{}}
\LANGpl{
\Question{
Wykresem funkcji
\[ 
 f(x) = \frac{2x + 3} 
 {2 - x}
\]
jest hiperbola. Wyznacz środek tej hiperboli.}
{\(S = [2;-2]\)}{\(S = [-2;2]\)}{\(S = [2;2]\)}{\(S ={\Bigl [ 2; \frac{3} 
{2}\Bigr ]}\)}{\(S ={\Bigl [ 2;-\frac{3}
{2}\Bigr ]}\)}{}}
\LANGcs{
\Question{
Určete střed hyperboly dané předpisem
\(f(x) = \frac{2x+3} 
 {2-x} \).}
{\(S = [2;-2]\)}{\(S = [-2;2]\)}{\(S = [2;2]\)}{\(S ={\Bigl [ 2; \frac{3} 
{2}\Bigr ]}\)}{\(S ={\Bigl [ 2;-\frac{3}
{2}\Bigr ]}\)}{}}
\LANGsk{
\Question{
Grafom funkcie
\[ 
 f(x) = \frac{2x + 3} 
 {2 - x}
\]
 je hyperbola. Ktorý z následujúcich bodov je stredom tejto hyperboly?}
{\(S = [2;-2]\)}{\(S = [-2;2]\)}{\(S = [2;2]\)}{\(S ={\Bigl [ 2; \frac{3} 
{2}\Bigr ]}\)}{\(S ={\Bigl [ 2;-\frac{3}
{2}\Bigr ]}\)}{}}
\LANGes{
\Question{
La gráfica de la función
\[ 
 f(x) = \frac{2x + 3} 
 {2 - x}
\]
es una hipérbola.  Determina su centro.}
{\(S = [2;-2]\)}{\(S = [-2;2]\)}{\(S = [2;2]\)}{\(S ={\Bigl [ 2; \frac{3} 
{2}\Bigr ]}\)}{\(S ={\Bigl [ 2;-\frac{3}
{2}\Bigr ]}\)}{}}
\End

\Data\ID{9000007510}
\Updated{2021-04-10 09:04:40}
\Level{B}
\SubArea{Rational functions}
\LANGen{
\Question{
The graph of the function 
\[ 
 f(x) = \frac{-x + 1} 
{1 + 3x}
\]
 is a hyperbola. Find the center of this hyperbola.}
{\(S = \left [-\frac{1}
{3};-\frac{1}
{3}\right ]\)}{\(S = \left [\frac{1}
{3}; \frac{1} 
{3}\right ]\)}{\(S = \left [1;-\frac{1}
{3}\right ]\)}{\(S = \left [-1;-\frac{1}
{3}\right ]\)}{\(S = \left [-\frac{1}
{3}; \frac{1} 
{3}\right ]\)}{}}
\LANGpl{
\Question{
Wykresem funkcji 
\[ 
 f(x) = \frac{-x + 1} 
{1 + 3x}
\]
jest hiperbola. Wyznacz środek tej hiperboli.}
{\(S = \left [-\frac{1}
{3};-\frac{1}
{3}\right ]\)}{\(S = \left [\frac{1}
{3}; \frac{1} 
{3}\right ]\)}{\(S = \left [1;-\frac{1}
{3}\right ]\)}{\(S = \left [-1;-\frac{1}
{3}\right ]\)}{\(S = \left [-\frac{1}
{3}; \frac{1} 
{3}\right ]\)}{}}
\LANGcs{
\Question{
Určete střed hyperboly dané předpisem
\(f(x) = \frac{-x+1} 
{1+3x} \).}
{\(S = \left [-\frac{1}
{3};-\frac{1}
{3}\right ]\)}{\(S = \left [\frac{1}
{3}; \frac{1} 
{3}\right ]\)}{\(S = \left [1;-\frac{1}
{3}\right ]\)}{\(S = \left [-1;-\frac{1}
{3}\right ]\)}{\(S = \left [-\frac{1}
{3}; \frac{1} 
{3}\right ]\)}{}}
\LANGsk{
\Question{
Grafom funkcie
\[ 
 f(x) = \frac{-x + 1} 
{1 + 3x}
\]
 je hyperbola. Ktorý z následujúcich bodov je stredom tejto hyperboly?}
{\(S = \left [-\frac{1}
{3};-\frac{1}
{3}\right ]\)}{\(S = \left [\frac{1}
{3}; \frac{1} 
{3}\right ]\)}{\(S = \left [1;-\frac{1}
{3}\right ]\)}{\(S = \left [-1;-\frac{1}
{3}\right ]\)}{\(S = \left [-\frac{1}
{3}; \frac{1} 
{3}\right ]\)}{}}
\LANGes{
\Question{
La gráfica de la función
\[ 
 f(x) = \frac{-x + 1} 
{1 + 3x}
\]
  es una hipérbola.  Determina su centro.}
{\(S = \left [-\frac{1}
{3};-\frac{1}
{3}\right ]\)}{\(S = \left [\frac{1}
{3}; \frac{1} 
{3}\right ]\)}{\(S = \left [1;-\frac{1}
{3}\right ]\)}{\(S = \left [-1;-\frac{1}
{3}\right ]\)}{\(S = \left [-\frac{1}
{3}; \frac{1} 
{3}\right ]\)}{}}
\End

\Data\ID{9000007601}
\Updated{2020-07-15 03:07:34}
\Level{B}
\SubArea{Rational functions}
\LANGen{
\Question{
Find the domain of the function \(f(x) = 1 + \frac{3} 
{x+2}\).}
{\(\mathbb{R}\setminus \{ - 2\}\)}{\(\mathbb{R}\setminus \{2\}\)}{\(\mathbb{R}\setminus \{1;2\}\)}{\(\mathbb{R}\setminus \{ - 2;1\}\)}{\(\mathbb{R}\)}{}}
\LANGpl{
\Question{
Wyznacz dziedzinę funkcji  \(f(x) = 1 + \frac{3} 
{x+2}\).}
{\(\mathbb{R}\setminus \{ - 2\}\)}{\(\mathbb{R}\setminus \{2\}\)}{\(\mathbb{R}\setminus \{1;2\}\)}{\(\mathbb{R}\setminus \{ - 2;1\}\)}{\(\mathbb{R}\)}{}}
\LANGcs{
\Question{
Určete definiční obor funkce \(f(x) = 1 + \frac{3} 
{x+2}\).}
{\(\mathbb{R}\setminus \{ - 2\}\)}{\(\mathbb{R}\setminus \{2\}\)}{\(\mathbb{R}\setminus \{1;2\}\)}{\(\mathbb{R}\setminus \{ - 2;1\}\)}{\(\mathbb{R}\)}{}}
\LANGsk{
\Question{
Určte definičný obor funkcie \(f(x) = 1 + \frac{3} 
{x+2}\).}
{\(\mathbb{R}\setminus \{ - 2\}\)}{\(\mathbb{R}\setminus \{2\}\)}{\(\mathbb{R}\setminus \{1;2\}\)}{\(\mathbb{R}\setminus \{ - 2;1\}\)}{\(\mathbb{R}\)}{}}
\LANGes{
\Question{
Determina el dominio de la función \(f(x) = 1 + \frac{3} 
{x+2}\).}
{\(\mathbb{R}\setminus \{ - 2\}\)}{\(\mathbb{R}\setminus \{2\}\)}{\(\mathbb{R}\setminus \{1;2\}\)}{\(\mathbb{R}\setminus \{ - 2;1\}\)}{\(\mathbb{R}\)}{}}
\End

\Data\ID{9000007602}
\Updated{2021-04-09 07:04:39}
\Level{B}
\SubArea{Rational functions}
\LANGen{
\Question{
Find the domain of the function \(f(x) = 2 - \frac{3} 
{x-2}\).}
{\(\mathbb{R}\setminus \{2\}\)}{\(\mathbb{R}\setminus \{ - 2\}\)}{\(\mathbb{R}\setminus \{ - 2;2\}\)}{\(\mathbb{R}\setminus \{ - 3\}\)}{\(\mathbb{R}\)}{}}
\LANGpl{
\Question{
Wyznacz dziedzinę funkcji \(f(x) = 2 - \frac{3} 
{x-2}\).}
{\(\mathbb{R}\setminus \{2\}\)}{\(\mathbb{R}\setminus \{ - 2\}\)}{\(\mathbb{R}\setminus \{ - 2;2\}\)}{\(\mathbb{R}\setminus \{ - 3\}\)}{\(\mathbb{R}\)}{}}
\LANGcs{
\Question{
Určete definiční obor funkce \(f(x) = 2 - \frac{3} 
{x-2}\).}
{\(\mathbb{R}\setminus \{2\}\)}{\(\mathbb{R}\setminus \{ - 2\}\)}{\(\mathbb{R}\setminus \{ - 2;2\}\)}{\(\mathbb{R}\setminus \{ - 3\}\)}{\(\mathbb{R}\)}{}}
\LANGsk{
\Question{
Určte definičný obor funkcie \(f(x) = 2 - \frac{3} 
{x-2}\).}
{\(\mathbb{R}\setminus \{2\}\)}{\(\mathbb{R}\setminus \{ - 2\}\)}{\(\mathbb{R}\setminus \{ - 2;2\}\)}{\(\mathbb{R}\setminus \{ - 3\}\)}{\(\mathbb{R}\)}{}}
\LANGes{
\Question{
Determina el dominio de la función \(f(x) = 2 - \frac{3} 
{x-2}\).}
{\(\mathbb{R}\setminus \{2\}\)}{\(\mathbb{R}\setminus \{ - 2\}\)}{\(\mathbb{R}\setminus \{ - 2;2\}\)}{\(\mathbb{R}\setminus \{ - 3\}\)}{\(\mathbb{R}\)}{}}
\End

\Data\ID{9000007606}
\Updated{2020-07-15 03:07:34}
\Level{B}
\SubArea{Rational functions}
\LANGen{
\Question{
Find the range of the function \(f(x) = 1 + \frac{3} 
{x+2}\).}
{\(\mathbb{R}\setminus \{1\}\)}{\(\mathbb{R}\setminus \{ - 2\}\)}{\(\mathbb{R}\setminus \{ - 2;1\}\)}{\([ 0;\infty )\)}{\(\mathbb{R}\)}{}}
\LANGpl{
\Question{
Wyznacz dziedzinę funkcji  \(f(x) = 1 + \frac{3} 
{x+2}\).}
{\(\mathbb{R}\setminus \{1\}\)}{\(\mathbb{R}\setminus \{ - 2\}\)}{\(\mathbb{R}\setminus \{ - 2;1\}\)}{\([ 0;\infty )\)}{\(\mathbb{R}\)}{}}
\LANGcs{
\Question{
Určete obor hodnot funkce \(f(x) = 1 + \frac{3} 
{x+2}\).}
{\(\mathbb{R}\setminus \{1\}\)}{\(\mathbb{R}\setminus \{ - 2\}\)}{\(\mathbb{R}\setminus \{ - 2;1\}\)}{\(\langle 0;\infty )\)}{\(\mathbb{R}\)}{}}
\LANGsk{
\Question{
Určte obor hodnôt funkcie \(f(x) = 1 + \frac{3} 
{x+2}\).}
{\(\mathbb{R}\setminus \{1\}\)}{\(\mathbb{R}\setminus \{ - 2\}\)}{\(\mathbb{R}\setminus \{ - 2;1\}\)}{\(\langle 0;\infty )\)}{\(\mathbb{R}\)}{}}
\LANGes{
\Question{
Determina el rango de la función  \(f(x) = 1 + \frac{3} 
{x+2}\).}
{\(\mathbb{R}\setminus \{1\}\)}{\(\mathbb{R}\setminus \{ - 2\}\)}{\(\mathbb{R}\setminus \{ - 2;1\}\)}{\([ 0;\infty )\)}{\(\mathbb{R}\)}{}}
\End

\Data\ID{9000007607}
\Updated{2020-07-15 03:07:34}
\Level{B}
\SubArea{Rational functions}
\LANGen{
\Question{
Find the range of the function \(f(x) = 2 - \frac{3} 
{x-2}\).}
{\(\mathbb{R}\setminus \{2\}\)}{\(\mathbb{R}\setminus \{ - 2\}\)}{\(\mathbb{R}\setminus \{ - 2;3\}\)}{\((0;\infty )\)}{\(\mathbb{R}\)}{}}
\LANGpl{
\Question{
Wyznacz dziedzinę funkcji \(f(x) = 2 - \frac{3} 
{x-2}\).}
{\(\mathbb{R}\setminus \{2\}\)}{\(\mathbb{R}\setminus \{ - 2\}\)}{\(\mathbb{R}\setminus \{ - 2;3\}\)}{\((0;\infty )\)}{\(\mathbb{R}\)}{}}
\LANGcs{
\Question{
Určete obor hodnot funkce \(f(x) = 2 - \frac{3} 
{x-2}\).}
{\(\mathbb{R}\setminus \{2\}\)}{\(\mathbb{R}\setminus \{ - 2\}\)}{\(\mathbb{R}\setminus \{ - 2;3\}\)}{\((0;\infty )\)}{\(\mathbb{R}\)}{}}
\LANGsk{
\Question{
Určte obor hodnôt funkcie \(f(x) = 2 - \frac{3} 
{x-2}\).}
{\(\mathbb{R}\setminus \{2\}\)}{\(\mathbb{R}\setminus \{ - 2\}\)}{\(\mathbb{R}\setminus \{ - 2;3\}\)}{\((0;\infty )\)}{\(\mathbb{R}\)}{}}
\LANGes{
\Question{
Determina el rango de la función  \(f(x) = 2 - \frac{3} 
{x-2}\).}
{\(\mathbb{R}\setminus \{2\}\)}{\(\mathbb{R}\setminus \{ - 2\}\)}{\(\mathbb{R}\setminus \{ - 2;3\}\)}{\((0;\infty )\)}{\(\mathbb{R}\)}{}}
\End

\Data\ID{9000007702}
\Updated{2021-04-09 07:04:33}
\Level{B}
\SubArea{Rational functions}
\LANGen{
\Question{
Identify a correct statement which concerns the function
\(f(x) = \frac{1} 
{-x+2}\).}
{None of the statements above is true.}{The function \(f\) 
is an increasing function.}{The function \(f\) 
is bounded below.}{The function \(f\) 
has a maximum at \(x = 2\).}{The function \(f\) 
is decreasing on \((2;\infty )\).}{}}
\LANGpl{
\Question{
Które ze stwierdzeń dotyczących funkcji
\(f(x) = \frac{1} 
{-x+2}\) jest prawdziwe?}
{Żadne ze stwierdzeń nie jest prawdziwe.}{Funkcja \(f\) 
jest funkcją rosnącą.}{Funkcja  \(f\) 
jest ograniczona z dołu.}{Funkcja  \(f\) 
osiąga maksimum w punkcie \(x = 2\).}{Funkcja  \(f\) jest malejąca w przedziale
 \((2;\infty )\).}{}}
\LANGcs{
\Question{
Je dána funkce \(f(x) = \frac{1} 
{-x+2}\).
Označte pravdivé tvrzení o funkci
\(f\).}
{Funkce nemá žádnou z uvedených vlastností.}{Funkce je rostoucí.}{Funkce je zdola omezená.}{Funkce má maximum v bodě \(2\).}{Funkce je klesající na intervalu
\((2;\infty )\).}{}}
\LANGsk{
\Question{
Daná je funkcia \(f(x) = \frac{1} 
{-x+2}\). Vyberte pravdivé tvrdenie.}
{Funkcia nemá žiadnu z uvedených vlastností.}{Funkcia je rastúca.}{Funkcia je zdola ohraničená.}{Funkcia má maximum v bode \(2\).}{Funkcia je klesajúca na intervale
\((2;\infty )\).}{}}
\LANGes{
\Question{
Identifica una proposición verdadera que se refiere a la función
\(f(x) = \frac{1} 
{-x+2}\).}
{Ninguna de las proposiciones anteriores es cierta.}{La función \(f\) 
es creciente.}{La función \(f\) 
está acotada inferiormente.}{La función \(f\) 
tiene su máximo en \(x = 2\).}{La función \(f\) 
es decreciente en \((2;\infty )\).}{}}
\End

\Data\ID{9000007707}
\Updated{2020-07-15 03:07:34}
\Level{B}
\SubArea{Rational functions}
\LANGen{
\Question{
Identify a correct statement which concerns the function
\(f(x) = 2 -\frac{1} 
{x}\).}
{None of the statements above is true.}{The function \(f\) 
is bounded above.}{The function \(f\) 
is an even function.}{The function \(f\) 
is a bounded function.}{The function \(f\) 
is an odd function.}{}}
\LANGpl{
\Question{
Wybierz poprawne stwierdzenie dotyczące funkcji
\(f(x) = 2 -\frac{1} 
{x}\).}
{Żadne ze stwierdzeń nie jest prawdziwe.}{Funkcja  \(f\) 
jest ograniczona z dołu.}{Funkcja \(f\) 
jest funkcja parzystą.}{Funkcja \(f\) 
jest funkcją ograniczoną.}{Funkcja \(f\) 
jest funkcja nieparzystą.}{}}
\LANGcs{
\Question{
Je dána funkce \(f(x) = 2 -\frac{1} 
{x}\).
Označte pravdivé tvrzení o funkci
\(f\).}
{Funkce nemá žádnou z uvedených vlastností.}{Funkce je shora omezená.}{Funkce je sudá.}{Funkce je omezená.}{Funkce je lichá.}{}}
\LANGsk{
\Question{
Daná je funkcia \(f(x) = 2 -\frac{1} 
{x}\), vyberte pravdivé tvrdenie.}
{Funkcia nemá žiadnu z uvedených vlastností.}{Funkcia je zhora ohraničená.}{Funkcia je párna.}{Funkcia je ohraničená.}{Funkcia je nepárna.}{}}
\LANGes{
\Question{
Identifica una proposición verdadera que se refiere a la función
\(f(x) = 2 -\frac{1} 
{x}\).}
{Ninguna de las proposiciones anteriores es cierta.}{La función \(f\) 
está acotada superiormente.}{La función \(f\) 
es una función par.}{La función \(f\) 
está acotada.}{La función \(f\) 
es una función impar.}{}}
\End

\Data\ID{9000007709}
\Updated{2021-04-09 07:04:18}
\Level{B}
\SubArea{Rational functions}
\LANGen{
\Question{
Identify a correct statement which concerns the function
\(f(x) = -\frac{5}
{x} - 3\).}
{None of the statements above is true.}{The function \(f\) 
is bounded above.}{The function \(f\) 
is an even function.}{The function \(f\) is a
decreasing function on \((0;\infty )\).}{The function \(f\) 
is an odd function.}{}}
\LANGpl{
\Question{
Wybierz prawdziwe stwierdzenie dotyczące funkcji
\(f(x) = -\frac{5}
{x} - 3\).}
{Żadne ze stwierdzeń nie jest prawdziwe.}{Funkcja \(f\) 
jest ograniczona z dołu.}{Funkcja  \(f\) 
jest parzysta.}{Funkcja  \(f\) jest funkcją rosnącą w przedziale \((0;\infty )\).}{Funkcja \(f\) 
jest funkcją nieparzystą.}{}}
\LANGcs{
\Question{
Je dána funkce \(f(x) = -\frac{5}
{x} - 3\).
Označte pravdivé tvrzení o funkci
\(f\).}
{Funkce nemá žádnou z uvedených vlastností.}{Funkce je shora omezená.}{Funkce je sudá.}{Funkce je klesající na intervalu
\((0;\infty )\).}{Funkce je lichá.}{}}
\LANGsk{
\Question{
Daná je funkcia \(f(x) = -\frac{5}
{x} - 3\), vyberte pravdivé tvrdenie.}
{Funkcia nemá žiadnu z uvedených vlastností.}{Funkcia je zhora ohraničená.}{Funkcia je párna.}{Funkcia je klesajúca na intervale
\((0;\infty )\).}{Funkcia je nepárna.}{}}
\LANGes{
\Question{
Identifica una proposición verdadera que se refiere a la función 
\(f(x) = -\frac{5}
{x} - 3\).}
{Ninguna de las proposiciones anteriores es cierta.}{La función \(f\) 
está acotada inferiormente.}{La función \(f\) 
es una función par.}{La función \(f\) es decreciente en  \((0;\infty )\).}{La función \(f\) 
es una función impar.}{}}
\End

\Data\ID{9000014201}
\Updated{2022-04-23 05:04:05}
\Level{B}
\SubArea{Rational functions}
\LANGen{
\Question{
Find intersection points of the graph of the rational function 
\(
f(x) = \frac{2x - 3} 
 {x - 2}
\)
with \(y\)-axis.}
{\(Y = \left [0; \frac{3} 
{2}\right ]\)}{\(Y = \left [\frac{3}
{2};0\right ]\)}{\(Y _{1} = \left [0; \frac{3} 
{2}\right ]\text{ and }Y _{2} = \left [\frac{3}
{2};0\right ]\)}{\(Y = \left [2;2\right ]\)}{}{}}
\LANGpl{
\Question{
Wyznacz punkty przecięcia wykresu funkcji wymiernej
\[ 
 f\colon y = \frac{2x - 3} 
 {x - 2}
\]
z osia współrzędnych \(y\).}
{\(Y = \left [0; \frac{3} 
{2}\right ]\)}{\(Y = \left [\frac{3}
{2};0\right ]\)}{\(Y _{1} = \left [0; \frac{3} 
{2}\right ]\text{ i }Y _{2} = \left [\frac{3}
{2};0\right ]\)}{\(Y = \left [2;2\right ]\)}{}{}}
\LANGcs{
\Question{
Určete společné body osy \(y\) a
grafu lineární lomené funkce \(f\colon y = \frac{2x-3} 
 {x-2} \).}
{\(Y = \left [0; \frac{3} 
{2}\right ]\)}{\(Y = \left [\frac{3}
{2};0\right ]\)}{\(Y _{1} = \left [0; \frac{3} 
{2}\right ] \wedge Y _{2} = \left [\frac{3}
{2};0\right ]\)}{\(Y = \left [2;2\right ]\)}{}{}}
\LANGsk{
\Question{
Určte priesečníky osy \(y\) a
grafu lineárnej lomenej funkcie \(f\colon y = \frac{2x-3} 
 {x-2} \).}
{\(Y = \left [0; \frac{3} 
{2}\right ]\)}{\(Y = \left [\frac{3}
{2};0\right ]\)}{\(Y _{1} = \left [0; \frac{3} 
{2}\right ] \wedge Y _{2} = \left [\frac{3}
{2};0\right ]\)}{\(Y = \left [2;2\right ]\)}{}{}}
\LANGes{
\Question{
Determina los puntos de intersección de la gráfica de la función racional
\[ 
 f\colon y = \frac{2x - 3} 
 {x - 2}
\]
con el eje \(y\).}
{\(Y = \left [0; \frac{3} 
{2}\right ]\)}{\(Y = \left [\frac{3}
{2};0\right ]\)}{\(Y _{1} = \left [0; \frac{3} 
{2}\right ]\text{ and }Y _{2} = \left [\frac{3}
{2};0\right ]\)}{\(Y = \left [2;2\right ]\)}{}{}}
\End

\Data\ID{9000014202}
\Updated{2022-04-23 05:04:05}
\Level{B}
\SubArea{Rational functions}
\LANGen{
\Question{
Find the intersections of the graph of the rational function
\(f(x) = \frac{x+2} 
{2-x}\) with
\(x\)-axis.}
{\(X = \left [-2;0\right ]\)}{\(X = \left [0;-2\right ]\)}{\(X_{1} = \left [0;-2\right ]\text{ and }X_{2} = \left [-2;0\right ]\)}{\(X = \left [2;0\right ]\)}{}{}}
\LANGpl{
\Question{
Wyznacz punty przecięcia wykresu funkcji wymiernej
\(f\colon y = \frac{x+2} 
{2-x}\) z osią
\(x\).}
{\(X = \left [-2;0\right ]\)}{\(X = \left [0;-2\right ]\)}{\(X_{1} = \left [0;-2\right ]\text{ i }X_{2} = \left [-2;0\right ]\)}{\(X = \left [2;0\right ]\)}{}{}}
\LANGcs{
\Question{
Určete společné body osy \(x\) a
grafu lineární lomené funkce \(f\colon y = \frac{x+2} 
{2-x}\).}
{\(X = \left [-2;0\right ]\)}{\(X = \left [0;-2\right ]\)}{\(X_{1} = \left [0;-2\right ] \wedge X_{2} = \left [-2;0\right ]\)}{\(X = \left [2;0\right ]\)}{}{}}
\LANGsk{
\Question{
Určte priesečníky osy \(x\) a
grafu lineárnej lomenej funkcie \(f\colon y = \frac{x+2} {2-x}\).}
{\(X = \left [-2;0\right ]\)}{\(X = \left [0;-2\right ]\)}{\(X_{1} = \left [0;-2\right ] \wedge X_{2} = \left [-2;0\right ]\)}{\(X = \left [2;0\right ]\)}{}{}}
\LANGes{
\Question{
Determina los puntos de intersección de la gráfica de la función racional
\(f\colon y = \frac{x+2} 
{2-x}\) con el eje 
\(x\).}
{\(X = \left [-2;0\right ]\)}{\(X = \left [0;-2\right ]\)}{\(X_{1} = \left [0;-2\right ]\text{ y }X_{2} = \left [-2;0\right ]\)}{\(X = \left [2;0\right ]\)}{}{}}
\End

\Data\ID{9000014203}
\Updated{2021-04-10 09:04:58}
\Level{B}
\SubArea{Rational functions}
\LANGen{
\Question{
Which of the statements from the following list is true for the function
\(f(x) = -\frac{2}
{x} + 1\)?}
{The function \(f\) 
is a one-to-one function.}{The function \(f\) 
is an odd function.}{The function \(f\) 
is an increasing function.}{The graph of the function \(f\) 
is a hyperbola with branches in the second and fourth quadrant.}{}{}}
\LANGpl{
\Question{
Które z poniższych stwierdzeń jest prawdziwe dla funkcji
\(f\colon y = -\frac{2}
{x} + 1\)?}
{Funkcja \(f\) 
jest funkcją jeden do jednego.}{Funkcja \(f\) 
jest funkcją nieparzystą.}{Funkcja   \(f\) 
jest funkcją rosnącą.}{Wykresem funkcji  \(f\) 
jest hiperbola, której gałęzie znajdują się w drugiej i czwartej ćwiartce układu współrzędnych.}{}{}}
\LANGcs{
\Question{
Který z následujících výroků o funkci
\(f\colon y = -\frac{2}
{x} + 1\) je
pravdivý?}
{Funkce \(f\) 
je prostá.}{Funkce \(f\) 
je lichá.}{Funkce \(f\) je
rostoucí v celém \(D(f)\).}{Grafem funkce \(f\) 
je hyperbola, jejíž větve leží ve II. a IV. kvadrantu souřadnicového
systému.}{}{}}
\LANGsk{
\Question{
Daná je funkcia \(f\colon y = -\frac{2}
{x} + 1\). Vyberte pravdivé tvrdenie o funkcii \(f\).}
{Funkcia \(f\) 
je prostá.}{Funkcia \(f\) 
je nepárna.}{Funkcia \(f\) je
rastúca na celom \(D(f)\).}{Grafom funkcie \(f\) 
je hyperbola, ktorej časti ležia v II. a IV. kvadrante súradnicového
systému.}{}{}}
\LANGes{
\Question{
¿Cuál de las proposiciones es verdadera para la función
\(f(x) = -\frac{2}
{x} + 1\)?}
{La función \(f\) 
es una función uno a uno (inyectiva).}{La función \(f\) 
es impar.}{La función \(f\) 
es creciente.}{La gráfica de la función \(f\) 
es una hipérbola cuyas ramas están en el segundo y cuarto cuadrante.}{}{}}
\End

\Data\ID{9000014206}
\Updated{2022-04-23 05:04:42}
\Level{B}
\SubArea{Rational functions}
\LANGen{
\Question{
Find the domain \(\mathrm{Dom}(f)\) 
and range \(\mathop{\mathrm{Ran}}(f)\) of
the function \(f(x) = \frac{2+x} 
{x+4}\).}
{\begin{align*}
\mathrm{Dom}(f) &= (-\infty ;-4)\cup (-4;\infty ),\\
\mathop{\mathrm{Ran}}(f) &= (-\infty ;1)\cup (1;\infty )
\end{align*}}{\begin{align*}
\mathrm{Dom}(f) &= (-\infty ;4)\cup (4;\infty ), \\
\mathop{\mathrm{Ran}}(f) &= (-\infty ;1)\cup (1;\infty )
\end{align*}}{\begin{align*}
\mathrm{Dom}(f) &= (-\infty ;2)\cup (2;\infty ),\\
\mathop{\mathrm{Ran}}(f) &= (-\infty ;4)\cup (4;\infty )
\end{align*}}{\begin{align*}
\mathrm{Dom}(f) &= (-\infty ;4)\cup (4;\infty ), \\ 
\mathop{\mathrm{Ran}}(f) &= (-\infty ;2)\cup (2;\infty )
\end{align*}}{}{}}
\LANGpl{
\Question{
Wyznacz dziedzinę  \(D(f)\) 
i zakres \(H(f)\)
funkcji \(f\colon y = \frac{2+x} 
{x+4}\).}
{\begin{align*}
D(f) &= (-\infty ;-4)\cup (-4;\infty ),\\
H(f) &= (-\infty ;1)\cup (1;\infty )
\end{align*}}{\begin{align*}
D(f) &= (-\infty ;4)\cup (4;\infty ), \\
H(f) &= (-\infty ;1)\cup (1;\infty )
\end{align*}}{\begin{align*}
D(f) &= (-\infty ;2)\cup (2;\infty ), \\
H(f) &= (-\infty ;4)\cup (4;\infty )
\end{align*}}{\begin{align*}
D(f) &= (-\infty ;4)\cup (4;\infty ), \\ 
H(f) &= (-\infty ;2)\cup (2;\infty )
\end{align*}}{}{}}
\LANGcs{
\Question{
Pro funkci \(f\colon y = \frac{2+x} 
{x+4}\) vyber
správnou kombinaci \(D(f)\) 
a \(H(f)\).}
{\begin{align*}
D(f) &= (-\infty ;-4)\cup (-4;\infty ), \\
H(f) &= (-\infty ;1)\cup (1;\infty )
\end{align*}}{\begin{align*}
D(f) &= (-\infty ;4)\cup (4;\infty ),\\
H(f) &= (-\infty ;1)\cup (1;\infty )\end{align*}}{\begin{align*}
D(f) &= (-\infty ;2)\cup (2;\infty ),\\
H(f) &= (-\infty ;4)\cup (4;\infty )\end{align*}}{\begin{align*}
D(f) &= (-\infty ;4)\cup (4;\infty ), \\
H(f) &= (-\infty ;2)\cup (2;\infty )\end{align*}}{}{}}
\LANGsk{
\Question{
Určte \(D(f)\) a \(H(f)\) funkcie \(f\colon y = \frac{2+x} 
{x+4}\).}
{\begin{align*}
D(f) &= (-\infty ;-4)\cup (-4;\infty ), \\
H(f) &= (-\infty ;1)\cup (1;\infty )\end{align*}}{\begin{align*}
D(f) &= (-\infty ;4)\cup (4;\infty ), \\
H(f) &= (-\infty ;1)\cup (1;\infty )\end{align*}}{\begin{align*}
D(f) &= (-\infty ;2)\cup (2;\infty ),\\
H(f) &= (-\infty ;4)\cup (4;\infty )\end{align*}}{\begin{align*}
D(f) &= (-\infty ;4)\cup (4;\infty ), \\
H(f) &= (-\infty ;2)\cup (2;\infty )\end{align*}}{}{}}
\LANGes{
\Question{
Determina el dominio \(\mathrm{Dom}(f)\) 
y el rango  \(\mathop{\mathrm{Ran}}(f)\) de la función \(f(x) = \frac{2+x} 
{x+4}\).}
{\begin{align*}
\mathrm{Dom}(f) &= (-\infty ;-4)\cup (-4;\infty ),\\
\mathop{\mathrm{Ran}}(f) &= (-\infty ;1)\cup (1;\infty )
\end{align*}}{\begin{align*}
\mathrm{Dom}(f) &= (-\infty ;4)\cup (4;\infty ), \\
\mathop{\mathrm{Ran}}(f) &= (-\infty ;1)\cup (1;\infty )
\end{align*}}{\begin{align*}
\mathrm{Dom}(f) &= (-\infty ;2)\cup (2;\infty ),\\
\mathop{\mathrm{Ran}}(f) &= (-\infty ;4)\cup (4;\infty )
\end{align*}}{\begin{align*}
\mathrm{Dom}(f) &= (-\infty ;4)\cup (4;\infty ), \\ 
\mathop{\mathrm{Ran}}(f) &= (-\infty ;2)\cup (2;\infty )
\end{align*}}{}{}}
\End

\Data\ID{9000014207}
\Updated{2020-07-15 03:07:34}
\Level{B}
\SubArea{Rational functions}
\IMG{000142_07-figure0.pdf}
\LANGen{
\Question{
Identify a possible analytic expression for the function
\(f\) 
graphed in the picture.}
{\(f(x)= \frac{x+1} 
{x+2}\)}{\(f(x)= \frac{1} 
{x+2} - 1\)}{\(f(x) = \frac{1} 
{x+2} + 1\)}{\(f(x) = \frac{x+1} 
{x-2}\)}{}{}}
\LANGpl{
\Question{
Wyznacz wzór funkcji 
\(f\), której wykres przedstawiony jest  na rysunku poniżej.}
{\(f\colon y = \frac{x+1} 
{x+2}\)}{\(f\colon y = \frac{1} 
{x+2} - 1\)}{\(f\colon y = \frac{1} 
{x+2} + 1\)}{\(f\colon y = \frac{x+1} 
{x-2}\)}{}{}}
\LANGcs{
\Question{
Vyber správný předpis funkce \(f\),
jejíž graf je znázorněn na obrázku:}
{\(f\colon y = \frac{x+1} 
{x+2}\)}{\(f\colon y = \frac{1} 
{x+2} - 1\)}{\(f\colon y = \frac{1} 
{x+2} + 1\)}{\(f\colon y = \frac{x+1} 
{x-2}\)}{}{}}
\LANGsk{
\Question{
Určte správny predpis funkcie \(f\),
ktorej graf je znázornený na obrázku.}
{\(f\colon y = \frac{x+1} 
{x+2}\)}{\(f\colon y = \frac{1} 
{x+2} - 1\)}{\(f\colon y = \frac{1} 
{x+2} + 1\)}{\(f\colon y = \frac{x+1} 
{x-2}\)}{}{}}
\LANGes{
\Question{
Identifica una posible expresión analítica para la función \(f\) graficada en la imagen.}
{\(f(x)= \frac{x+1} 
{x+2}\)}{\(f(x)= \frac{1} 
{x+2} - 1\)}{\(f(x) = \frac{1} 
{x+2} + 1\)}{\(f(x) = \frac{x+1} 
{x-2}\)}{}{}}
\End

\Data\ID{9000014208}
\Updated{2020-07-15 03:07:34}
\Level{B}
\SubArea{Rational functions}
\IMG{000142_08-figure0.pdf}
\LANGen{
\Question{
Identify a possible analytic expression for the function
\(f\) 
graphed in the picture.}
{\(f(x) = \frac{2x+1} 
 {x-1} \)}{\(f(x) = \frac{3} 
{x-1} - 2\)}{\(f(x)= \frac{2x-1} 
 {x+1} \)}{\(f(x) = \frac{2x+2} 
 {x-1} \)}{}{}}
\LANGpl{
\Question{
Wyznacz wzór funkcji 
\(f\), której wykres przedstawiony jest  na rysunku poniżej.}
{\(f\colon y = \frac{2x+1} 
 {x-1} \)}{\(f\colon y = \frac{3} 
{x-1} - 2\)}{\(f\colon y = \frac{2x-1} 
 {x+1} \)}{\(f\colon y = \frac{2x+2} 
 {x-1} \)}{}{}}
\LANGcs{
\Question{
Vyber správný předpis funkce \(f\),
jejíž graf je znázorněn na obrázku:}
{\(f\colon y = \frac{2x+1} 
 {x-1} \)}{\(f\colon y = \frac{3} 
{x-1} - 2\)}{\(f\colon y = \frac{2x-1} 
 {x+1} \)}{\(f\colon y = \frac{2x+2} 
 {x-1} \)}{}{}}
\LANGsk{
\Question{
Určte správny predpis funkcie \(f\),
ktorej graf je znázornený na obrázku.}
{\(f\colon y = \frac{2x+1} 
 {x-1} \)}{\(f\colon y = \frac{3} 
{x-1} - 2\)}{\(f\colon y = \frac{2x-1} 
 {x+1} \)}{\(f\colon y = \frac{2x+2} 
 {x-1} \)}{}{}}
\LANGes{
\Question{
Identifica una posible expresión analítica para la función \(f\)  graficada en la imagen.}
{\(f(x) = \frac{2x+1} 
 {x-1} \)}{\(f(x) = \frac{3} 
{x-1} - 2\)}{\(f(x)= \frac{2x-1} 
 {x+1} \)}{\(f(x) = \frac{2x+2} 
 {x-1} \)}{}{}}
\End

\Data\ID{9000014209}
\Updated{2022-04-23 05:04:42}
\Level{B}
\SubArea{Rational functions}
\LANGen{
\Question{
Consider the function \(f(x) = \frac{3x+1} 
 {x-2} \).
Find all \(x\) 
such that \(f(x) > 0\).}
{\(x\in \left (-\infty ;-\frac{1}
{3}\right )\cup (2;\infty )\)}{\(x\in \left (-\frac{1}
{3};\infty \right )\)}{\(x\in (2;3)\)}{\(x\in (-\infty ;-3)\cup (2;\infty )\)}{}{}}
\LANGpl{
\Question{
Rozważ funkcję  \(f\colon y = \frac{3x+1} 
 {x-2} \).
Znajdź wszystkie wartości \(x\), 
dla których  \(f(x) > 0\).}
{\(x\in \left (-\infty ;-\frac{1}
{3}\right )\cup (2;\infty )\)}{\(x\in \left (-\frac{1}
{3};\infty \right )\)}{\(x\in (2;3)\)}{\(x\in (-\infty ;-3)\cup (2;\infty )\)}{}{}}
\LANGcs{
\Question{
Je dána funkce \(f\colon y = \frac{3x+1} 
 {x-2} \).
Pro které \(x\) 
platí \(f(x) > 0\)?}
{\(x\in \left (-\infty ;-\frac{1}
{3}\right )\cup (2;\infty )\)}{\(x\in \left (-\frac{1}
{3};\infty \right )\)}{\(x\in (2;3)\)}{\(x\in (-\infty ;-3)\cup (2;\infty )\)}{}{}}
\LANGsk{
\Question{
Daná je funkcia \(f\colon y = \frac{3x+1} 
 {x-2} \).
Určte všetky \(x\), pre ktoré 
platí \(f(x) > 0\).}
{\(x\in \left (-\infty ;-\frac{1}
{3}\right )\cup (2;\infty )\)}{\(x\in \left (-\frac{1}
{3};\infty \right )\)}{\(x\in (2;3)\)}{\(x\in (-\infty ;-3)\cup (2;\infty )\)}{}{}}
\LANGes{
\Question{
Considera la función \(f(x) = \frac{3x+1} 
 {x-2} \).
Determina todos los \(x\) 
tales que \(f(x) > 0\).}
{\(x\in \left (-\infty ;-\frac{1}
{3}\right )\cup (2;\infty )\)}{\(x\in \left (-\frac{1}
{3};\infty \right )\)}{\(x\in (2;3)\)}{\(x\in (-\infty ;-3)\cup (2;\infty )\)}{}{}}
\End

\Data\ID{9000014210}
\Updated{2022-04-23 05:04:42}
\Level{B}
\SubArea{Rational functions}
\LANGen{
\Question{
Consider the function \(f(x) = \frac{2x+1} 
 {x+3} \).
Find all \(x\) 
such that \(f(x) < 0\).}
{\(x\in \left (-3;-\frac{1}
{2}\right )\)}{\(x\in (-\infty ;-3)\cup \left (\frac{1}
{2};\infty \right )\)}{\(x\in (-3;\infty )\)}{\(x\in \left (-\infty ;-\frac{1}
{2}\right )\)}{}{}}
\LANGpl{
\Question{
Rozważ funkcję  \(f\colon y = \frac{2x+1} 
 {x+3} \).
Znajdź wszystkie wartości  \(x\),
dla których \(f(x) < 0\).}
{\(x\in \left (-3;-\frac{1}
{2}\right )\)}{\(x\in (-\infty ;-3)\cup \left (\frac{1}
{2};\infty \right )\)}{\(x\in (-3;\infty )\)}{\(x\in \left (-\infty ;-\frac{1}
{2}\right )\)}{}{}}
\LANGcs{
\Question{
Je dána funkce \(f\colon y = \frac{2x+1} 
 {x+3} \).
Pro které \(x\) 
platí \(f(x) < 0\)?}
{\(x\in \left (-3;-\frac{1}
{2}\right )\)}{\(x\in (-\infty ;-3)\cup (\frac{1}
{2};\infty )\)}{\(x\in (-3;\infty )\)}{\(x\in \left (-\infty ;-\frac{1}
{2}\right )\)}{}{}}
\LANGsk{
\Question{
Daná je funkcia \(f\colon y = \frac{2x+1} 
 {x+3} \).
Určte všetky \(x\), pre ktoré 
platí \(f(x) < 0\).}
{\(x\in \left (-3;-\frac{1}
{2}\right )\)}{\(x\in (-\infty ;-3)\cup (\frac{1}
{2};\infty )\)}{\(x\in (-3;\infty )\)}{\(x\in \left (-\infty ;-\frac{1}
{2}\right )\)}{}{}}
\LANGes{
\Question{
Considera la función \(f(x) = \frac{2x+1} 
 {x+3} \).
 Determina todos los \(x\) 
tales que  \(f(x) < 0\).}
{\(x\in \left (-3;-\frac{1}
{2}\right )\)}{\(x\in (-\infty ;-3)\cup \left (\frac{1}
{2};\infty \right )\)}{\(x\in (-3;\infty )\)}{\(x\in \left (-\infty ;-\frac{1}
{2}\right )\)}{}{}}
\End

\Data\ID{9100002909}
\Updated{2020-07-15 03:07:10}
\Level{B}
\SubArea{Rational functions}
\LANGen{
\Question{
Identify a possible graph of the function
\(f(x) = \frac{1} 
{x+2}\).}
{}{}{}{}{}{}}
\LANGpl{
\Question{
Który z przedstawionych poniżej wykresów jest wykresem funkcji 
\(f\colon y = \frac{1} 
{x+2}\)?}
{}{}{}{}{}{}}
\LANGcs{
\Question{
Přiřaďte graf k funkci dané předpisem
\(f\colon y = \frac{1} 
{x+2}\).}
{}{}{}{}{}{}}
\LANGsk{
\Question{
Určte, ktorý z nasledujúcich grafov je graf funkcie
\(f\colon y = \frac{1} 
{x+2}\).}
{}{}{}{}{}{}}
\LANGes{
\Question{
Identifica una posible gráfica de la función
\(f(x) = \frac{1} 
{x+2}\).}
{}{}{}{}{}{}}

\IMGanswer{000029_09-figure0.pdf}
\IMGanswer{000029_09-figure1.pdf}
\IMGanswer{000029_09-figure2.pdf}
\IMGanswer{000029_09-figure3.pdf}\End

\Data\ID{9100003203}
\Updated{2020-07-15 03:07:10}
\Level{B}
\SubArea{Rational functions}
\LANGen{
\Question{
Identify a possible graph of the function
\(f(x) = 1 -\frac{2} 
{x}\).}
{}{}{}{}{}{}}
\LANGpl{
\Question{
Który z przedstawionych poniżej wykresów jest wykresem funkcji 
\(f(x) = 1 -\frac{2} 
{x}\)?}
{}{}{}{}{}{}}
\LANGcs{
\Question{
K funkci určené předpisem \(f(x) = 1 -\frac{2} 
{x}\) 
najděte její graf.}
{}{}{}{}{}{}}
\LANGsk{
\Question{
Určte, ktorý z nasledujúcich grafov je graf funkcie
\(f(x) = 1 -\frac{2} 
{x}\).}
{}{}{}{}{}{}}
\LANGes{
\Question{
Identifica una posible gráfica de la función
\(f(x) = 1 -\frac{2} 
{x}\).}
{}{}{}{}{}{}}

\IMGanswer{000032_03-figure0.pdf}
\IMGanswer{000032_03-figure1.pdf}
\IMGanswer{000032_03-figure2.pdf}
\IMGanswer{000032_03-figure3.pdf}\End

\Data\ID{9100003204}
\Updated{2020-07-15 03:07:10}
\Level{B}
\SubArea{Rational functions}
\LANGen{
\Question{
Identify a possible graph of the function
\(f(x) = -2 -\frac{1} 
{x}\).}
{}{}{}{}{}{}}
\LANGpl{
\Question{
Który z przedstawionych poniżej wykresów jest wykresem funkcji 
\(f(x)= -2 -\frac{1} 
{x}\)?}
{}{}{}{}{}{}}
\LANGcs{
\Question{
K funkci určené předpisem \(f(x) = -2 -\frac{1} 
{x}\) 
najděte její graf.}
{}{}{}{}{}{}}
\LANGsk{
\Question{
Určte, ktorý z nasledujúcich grafov je graf funkcie
\(f(x) = -2 -\frac{1} 
{x}\).}
{}{}{}{}{}{}}
\LANGes{
\Question{
Identifica una posible gráfica de la función
\(f(x) = -2 -\frac{1} 
{x}\).}
{}{}{}{}{}{}}

\IMGanswer{000032_04-figure0.pdf}
\IMGanswer{000032_04-figure1.pdf}
\IMGanswer{000032_04-figure2.pdf}
\IMGanswer{000032_04-figure3.pdf}\End

\Data\ID{9100003205}
\Updated{2020-07-15 03:07:10}
\Level{B}
\SubArea{Rational functions}
\LANGen{
\Question{
Identify a possible graph of the function
\(f(x) = \frac{1} 
{x-2}\).}
{}{}{}{}{}{}}
\LANGpl{
\Question{
Który z przedstawionych poniżej wykresów jest wykresem funkcji 
\(f(x) = \frac{1} 
{x-2}\)?}
{}{}{}{}{}{}}
\LANGcs{
\Question{
K funkci určené předpisem \(f(x) = \frac{1} 
{x-2}\) 
najděte její graf.}
{}{}{}{}{}{}}
\LANGsk{
\Question{
Určte, ktorý z nasledujúcich grafov je graf funkcie
\(f(x) = \frac{1} 
{x-2}\).}
{}{}{}{}{}{}}
\LANGes{
\Question{
Identifica una posible gráfica de la función
\(f(x) = \frac{1} 
{x-2}\).}
{}{}{}{}{}{}}

\IMGanswer{000032_05-figure0.pdf}
\IMGanswer{000032_05-figure1.pdf}
\IMGanswer{000032_05-figure2.pdf}
\IMGanswer{000032_05-figure3.pdf}\End

\Data\ID{9100003206}
\Updated{2020-07-15 03:07:10}
\Level{B}
\SubArea{Rational functions}
\LANGen{
\Question{
Identify a possible graph of the function
\(f(x) = \frac{2} 
{x+1}\).}
{}{}{}{}{}{}}
\LANGpl{
\Question{
Który z przedstawionych poniżej wykresów jest wykresem funkcji 
\(f(x) = \frac{2} 
{x+1}\)?}
{}{}{}{}{}{}}
\LANGcs{
\Question{
K funkci určené předpisem \(f(x) = \frac{2} 
{x+1}\) 
najděte její graf.}
{}{}{}{}{}{}}
\LANGsk{
\Question{
Určte, ktorý z nasledujúcich grafov je graf funkcie
\(f(x) = \frac{2} 
{x+1}\).}
{}{}{}{}{}{}}
\LANGes{
\Question{
Identifica una posible gráfica de la función
\(f(x) = \frac{2} 
{x+1}\).}
{}{}{}{}{}{}}

\IMGanswer{9100003206a.pdf}
\IMGanswer{9100003206b.pdf}
\IMGanswer{9100003206c.pdf}
\IMGanswer{9100003206d.pdf}\End

\Data\ID{9100003207}
\Updated{2020-07-15 03:07:10}
\Level{B}
\SubArea{Rational functions}
\LANGen{
\Question{
Identify a possible graph of the function
\(f(x)= 2 + \frac{1} 
{x+1}\).}
{}{}{}{}{}{}}
\LANGpl{
\Question{
Który z przedstawionych poniżej wykresów jest wykresem funkcji 
\(f(x) = 2 + \frac{1} 
{x+1}\)?}
{}{}{}{}{}{}}
\LANGcs{
\Question{
K funkci určené předpisem \(f(x) = 2 + \frac{1} 
{x+1}\) 
najděte její graf.}
{}{}{}{}{}{}}
\LANGsk{
\Question{
Určte, ktorý z nasledujúcich grafov je graf funkcie
\(f(x) = 2 + \frac{1} 
{x+1}\).}
{}{}{}{}{}{}}
\LANGes{
\Question{
Identifica una posible gráfica de la función
\(f(x)= 2 + \frac{1} 
{x+1}\).}
{}{}{}{}{}{}}

\IMGanswer{000032_07-figure0.pdf}
\IMGanswer{000032_07-figure1.pdf}
\IMGanswer{000032_07-figure2.pdf}
\IMGanswer{000032_07-figure3.pdf}\End

\Data\ID{9100003208}
\Updated{2020-07-15 03:07:10}
\Level{B}
\SubArea{Rational functions}
\LANGen{
\Question{
Identify a possible graph of the function
\(f(x) = 2 - \frac{1} 
{x-1}\).}
{}{}{}{}{}{}}
\LANGpl{
\Question{
Który z przedstawionych poniżej wykresów jest wykresem funkcji 
\(f(x) = 2 - \frac{1} 
{x-1}\)?}
{}{}{}{}{}{}}
\LANGcs{
\Question{
K funkci určené předpisem \(f(x) = 2 - \frac{1} 
{x-1}\) 
najděte její graf.}
{}{}{}{}{}{}}
\LANGsk{
\Question{
Určte, ktorý z nasledujúcich grafov je graf funkcie
\(f(x) = 2 - \frac{1} 
{x-1}\).}
{}{}{}{}{}{}}
\LANGes{
\Question{
Identifica una posible gráfica de la función
\(f(x) = 2 - \frac{1} 
{x-1}\).}
{}{}{}{}{}{}}

\IMGanswer{000032_08-figure0.pdf}
\IMGanswer{000032_08-figure1.pdf}
\IMGanswer{000032_08-figure2.pdf}
\IMGanswer{000032_08-figure3.pdf}\End

\Data\ID{9100003209}
\Updated{2020-07-15 03:07:10}
\Level{B}
\SubArea{Rational functions}
\LANGen{
\Question{
Identify a possible graph of the function
\(f(x) = -2 + \frac{1} 
{x+1}\).}
{}{}{}{}{}{}}
\LANGpl{
\Question{
Który z przedstawionych poniżej wykresów jest wykresem funkcji 
\(f(x) = -2 + \frac{1} 
{x+1}\)?}
{}{}{}{}{}{}}
\LANGcs{
\Question{
K funkci určené předpisem \(f(x)= -2 + \frac{1} 
{x+1}\) 
najděte její graf.}
{}{}{}{}{}{}}
\LANGsk{
\Question{
Určte, ktorý z nasledujúcich grafov je graf funkcie
\(f(x) = -2 + \frac{1} 
{x+1}\).}
{}{}{}{}{}{}}
\LANGes{
\Question{
Identifica una posible gráfica de la función
\(f(x) = -2 + \frac{1} 
{x+1}\).}
{}{}{}{}{}{}}

\IMGanswer{000032_09-figure0.pdf}
\IMGanswer{000032_09-figure1.pdf}
\IMGanswer{000032_09-figure2.pdf}
\IMGanswer{000032_09-figure3.pdf}\End

\Data\ID{9100003210}
\Updated{2020-07-15 03:07:10}
\Level{B}
\SubArea{Rational functions}
\LANGen{
\Question{
Identify a possible graph of the function
\(f(x) = -2 - \frac{1} 
{x-1}\).}
{}{}{}{}{}{}}
\LANGpl{
\Question{
Który z przedstawionych poniżej wykresów jest wykresem funkcji 
\(f(x) = -2 - \frac{1} 
{x-1}\)?}
{}{}{}{}{}{}}
\LANGcs{
\Question{
K funkci určené předpisem \(f(x) = -2 - \frac{1} 
{x-1}\) 
najděte její graf.}
{}{}{}{}{}{}}
\LANGsk{
\Question{
Určte, ktorý z nasledujúcich grafov je graf funkcie
\(f(x)= -2 - \frac{1} 
{x-1}\).}
{}{}{}{}{}{}}
\LANGes{
\Question{
Identifica una posible gráfica de la función
\(f(x) = -2 - \frac{1} 
{x-1}\).}
{}{}{}{}{}{}}

\IMGanswer{000032_10-figure0.pdf}
\IMGanswer{000032_10-figure1.pdf}
\IMGanswer{000032_10-figure2.pdf}
\IMGanswer{000032_10-figure3.pdf}\End

\Data\ID{9100014204}
\Updated{2020-07-15 04:07:14}
\Level{B}
\SubArea{Rational functions}
\LANGen{
\Question{
Identify a possible graph for the function
\(f(x) = \frac{2} 
{x+1} - 1\).}
{}{}{}{}{}{}}
\LANGpl{
\Question{
Który z przedstawionych poniżej wykresów jest wykresem funkcji 
\(f\colon y = \frac{2} 
{x+1} - 1\)?}
{}{}{}{}{}{}}
\LANGcs{
\Question{
Přiřaďte k dané funkci \(f\colon y = \frac{2} 
{x+1} - 1\) 
odpovídající graf:}
{}{}{}{}{}{}}
\LANGsk{
\Question{
Určte, ktorý z nasledujúcich grafov je graf funkcie
\(f\colon y = \frac{2} 
{x+1} - 1\).}
{}{}{}{}{}{}}
\LANGes{
\Question{
Identifica una posible gráfica de la función
\(f(x) = \frac{2} 
{x+1} - 1\).}
{}{}{}{}{}{}}

\IMGanswer{000142_04-figure0.pdf}
\IMGanswer{000142_04-figure1.pdf}
\IMGanswer{000142_04-figure2.pdf}
\IMGanswer{000142_04-figure3.pdf}\End

\Data\ID{9100014205}
\Updated{2020-07-15 04:07:14}
\Level{B}
\SubArea{Rational functions}
\LANGen{
\Question{
Identify a possible graph for the function
\(f(x) = 2 - \frac{1} 
{x-3}\).}
{}{}{}{}{}{}}
\LANGpl{
\Question{
Który z przedstawionych poniżej wykresów jest wykresem funkcji
\(f\colon y = 2 - \frac{1} 
{x-3}\)?}
{}{}{}{}{}{}}
\LANGcs{
\Question{
Přiřaďte k dané funkci \(f\colon y = 2 - \frac{1} 
{x-3}\) 
odpovídající graf:}
{}{}{}{}{}{}}
\LANGsk{
\Question{
Určte, ktorý z nasledujúcich grafov je graf funkcie
\(f\colon y = 2 - \frac{1} 
{x-3}\).}
{}{}{}{}{}{}}
\LANGes{
\Question{
Identifica una posible gráfica de la función
\(f(x) = 2 - \frac{1} 
{x-3}\).}
{}{}{}{}{}{}}

\IMGanswer{000142_05-figure0.pdf}
\IMGanswer{000142_05-figure1.pdf}
\IMGanswer{000142_05-figure2.pdf}
\IMGanswer{000142_05-figure3.pdf}\End

\Data\ID{1003019401}
\Updated{2022-08-25 10:08:58}
\Level{C}
\SubArea{Rational functions}
\LANGen{
\Question{
Identify which of the given functions is even.}
{\( m(x)=\frac{x^4-3}{x^2} \)}{\( h(x)=\frac{x^3+2}{x^2} \)}{\( g(x)=\frac{x^2}{x+1} \)}{\( f(x)=\frac{x^3}{x^2-5} \)}{}{}}
\LANGpl{
\Question{
Która z poniższych funkcji jest parzysta?}
{\( m(x)=\frac{x^4-3}{x^2} \)}{\( h(x)=\frac{x^3+2}{x^2} \)}{\( g(x)=\frac{x^2}{x+1} \)}{\( f(x)=\frac{x^3}{x^2-5} \)}{}{}}
\LANGcs{
\Question{
Která z následujících funkcí je sudá?}
{\( m(x)=\frac{x^4-3}{x^2} \)}{\( h(x)=\frac{x^3+2}{x^2} \)}{\( g(x)=\frac{x^2}{x+1} \)}{\( f(x)=\frac{x^3}{x^2-5} \)}{}{}}
\LANGsk{
\Question{
Ktorá z následujúcich funkcií je párna?}
{\( m(x)=\frac{x^4-3}{x^2} \)}{\( h(x)=\frac{x^3+2}{x^2} \)}{\( g(x)=\frac{x^2}{x+1} \)}{\( f(x)=\frac{x^3}{x^2-5} \)}{}{}}
\LANGes{
\Question{
Identifica cuál de las funciones dadas es par.}
{\( m(x)=\frac{x^4-3}{x^2} \)}{\( h(x)=\frac{x^3+2}{x^2} \)}{\( g(x)=\frac{x^2}{x+1} \)}{\( f(x)=\frac{x^3}{x^2-5} \)}{}{}}
\End

\Data\ID{1003019402}
\Updated{2022-08-25 10:08:58}
\Level{C}
\SubArea{Rational functions}
\LANGen{
\Question{
Identify which of the given functions  is neither even nor odd.}
{\( h(x) = \frac{x^3-7}{2x^2} \)}{\( f(x) = \frac{x^2+3}{3x^4} \)}{\( g(x) = \frac{5x^2-6}{x} \)}{\( m(x) = \frac{x^3}{2x^2+5} \)}{}{}}
\LANGpl{
\Question{
Która z podanych funkcji nie jest ani parzysta ani nieparzysta?}
{\( h(x) = \frac{x^3-7}{2x^2} \)}{\( f(x) = \frac{x^2+3}{3x^4} \)}{\( g(x) = \frac{5x^2-6}{x} \)}{\( m(x) = \frac{x^3}{2x^2+5} \)}{}{}}
\LANGcs{
\Question{
Která z následujících funkcí není ani sudá ani lichá?}
{\( h(x) = \frac{x^3-7}{2x^2} \)}{\( f(x) = \frac{x^2+3}{3x^4} \)}{\( g(x) = \frac{5x^2-6}{x} \)}{\( m(x) = \frac{x^3}{2x^2+5} \)}{}{}}
\LANGsk{
\Question{
Ktorá z následujúcich funkcií nie je párna ani nepárna?}
{\( h(x) = \frac{x^3-7}{2x^2} \)}{\( f(x) = \frac{x^2+3}{3x^4} \)}{\( g(x) = \frac{5x^2-6}{x} \)}{\( m(x) = \frac{x^3}{2x^2+5} \)}{}{}}
\LANGes{
\Question{
Identifica cuál de las funciones dadas no es par ni impar.}
{\( h(x) = \frac{x^3-7}{2x^2} \)}{\( f(x) = \frac{x^2+3}{3x^4} \)}{\( g(x) = \frac{5x^2-6}{x} \)}{\( m(x) = \frac{x^3}{2x^2+5} \)}{}{}}
\End

\Data\ID{1003028401}
\Updated{2022-08-25 10:08:04}
\Level{C}
\SubArea{Rational functions}
\LANGen{
\Question{
Let \( f(x)=\frac{3x-9}{x^2-3} \). Which of the statements about the domain and the range of the function \( f \) is true?}
{\( 3\in D(f) \wedge 3\in H(f) \)}{\( 3\in D(f) \wedge 3\notin H(f) \)}{\( 3\notin D(f) \wedge 3\in H(f) \)}{\( 3\notin D(f) \wedge 3\notin H(f) \)}{}{}}
\LANGpl{
\Question{
Niech \( f(x)=\frac{3x-9}{x^2-3} \). Które ze stwierdzeń dotyczących dziedziny i zakresu funkcji  \( f \) jest prawdziwe?}
{\( 3\in D(f) \wedge 3\in H(f) \)}{\( 3\in D(f) \wedge 3\notin H(f) \)}{\( 3\notin D(f) \wedge 3\in H(f) \)}{\( 3\notin D(f) \wedge 3\notin H(f) \)}{}{}}
\LANGcs{
\Question{
Funkce \(f\) je dána předpisem \( f(x)=\frac{3x-9}{x^2-3} \). Vyberte pravdivý výrok o definičním oboru \(D(f)\) a oboru hodnot \(H(f)\) funkce \(f\).}
{\( 3\in D(f) \wedge 3\in H(f) \)}{\( 3\in D(f) \wedge 3\notin H(f) \)}{\( 3\notin D(f) \wedge 3\in H(f) \)}{\( 3\notin D(f) \wedge 3\notin H(f) \)}{}{}}
\LANGsk{
\Question{
Daná je funkcia \( f(x)=\frac{3x-9}{x^2-3} \). Ktoré tvrdenie o definičnom obore a obore hodnôt funkcie \( f \) je pravdivé?}
{\( 3\in D(f) \wedge 3\in H(f) \)}{\( 3\in D(f) \wedge 3\notin H(f) \)}{\( 3\notin D(f) \wedge 3\in H(f) \)}{\( 3\notin D(f) \wedge 3\notin H(f) \)}{}{}}
\LANGes{
\Question{
Sea \( f(x)=\frac{3x-9}{x^2-3} \). ¿Cuál de las proposiciones sobre el dominio y el rango de la función \( f \)  es verdadera?}
{\( 3\in D(f) \wedge 3\in H(f) \)}{\( 3\in D(f) \wedge 3\notin H(f) \)}{\( 3\notin D(f) \wedge 3\in H(f) \)}{\( 3\notin D(f) \wedge 3\notin H(f) \)}{}{}}
\End

\Data\ID{1003028402}
\Updated{2022-08-25 10:08:04}
\Level{C}
\SubArea{Rational functions}
\LANGen{
\Question{
Let \( f(x)=\frac{2x-4}{x^2-4} \). Which of the statements about the domain and the range of the function \( f \) is true?}
{\( -2\notin D(f) \wedge -2\in H(f) \)}{\( -2\in D(f) \wedge -2\notin H(f) \)}{\( -2\in D(f) \wedge -2\in H(f) \)}{\( -2\notin D(f) \wedge -2\notin H(f) \)}{}{}}
\LANGpl{
\Question{
Niech \( f(x)=\frac{2x-4}{x^2-4} \). Które ze stwierdzeń dotyczących dziedziny i zakresu funkcji  \( f \) jest prawdziwe?}
{\( -2\notin D(f) \wedge -2\in H(f) \)}{\( -2\in D(f) \wedge -2\notin H(f) \)}{\( -2\in D(f) \wedge -2\in H(f) \)}{\( -2\notin D(f) \wedge -2\notin H(f) \)}{}{}}
\LANGcs{
\Question{
Funkce \(f\) je dána předpisem \( f(x)=\frac{2x-4}{x^2-4} \). Vyberte pravdivý výrok o definičním oboru \(D(f)\) a oboru hodnot \(H(f)\) funkce \(f\).}
{\( -2\notin D(f) \wedge -2\in H(f) \)}{\( -2\in D(f) \wedge -2\notin H(f) \)}{\( -2\in D(f) \wedge -2\in H(f) \)}{\( -2\notin D(f) \wedge -2\notin H(f) \)}{}{}}
\LANGsk{
\Question{
Daná je funkcia \( f(x)=\frac{2x-4}{x^2-4} \). Ktoré tvrdenie o definičnom obore a obore hodnôt funkcie \( f \) je pravdivé?}
{\( -2\notin D(f) \wedge -2\in H(f) \)}{\( -2\in D(f) \wedge -2\notin H(f) \)}{\( -2\in D(f) \wedge -2\in H(f) \)}{\( -2\notin D(f) \wedge -2\notin H(f) \)}{}{}}
\LANGes{
\Question{
Sea \( f(x)=\frac{2x-4}{x^2-4} \). ¿Cuál de las proposiciones sobre el dominio y el rango de la función \( f \) es verdadera?}
{\( -2\notin D(f) \wedge -2\in H(f) \)}{\( -2\in D(f) \wedge -2\notin H(f) \)}{\( -2\in D(f) \wedge -2\in H(f) \)}{\( -2\notin D(f) \wedge -2\notin H(f) \)}{}{}}
\End

\Data\ID{1003030903}
\Updated{2022-08-25 10:08:38}
\Level{C}
\SubArea{Rational functions}
\LANGen{
\Question{
Identify which of the following functions is bounded below.}
{\( h(x)=\frac{2x-4}{x-2} \)}{\( m(x)=\frac{x^2-4}{x-2} \)}{\( g(x)=\frac{4-x^2}{x-2} \)}{\( f(x)= -\left|\frac{x^2-4}{x-2}\right| \)}{}{}}
\LANGpl{
\Question{
Określ, która z podanych funkcji jest ograniczona z dołu.}
{\( h(x)=\frac{2x-4}{x-2} \)}{\( m(x)=\frac{x^2-4}{x-2} \)}{\( g(x)=\frac{4-x^2}{x-2} \)}{\( f(x)= -\left|\frac{x^2-4}{x-2}\right| \)}{}{}}
\LANGcs{
\Question{
Která z následujících funkcí je zdola omezená?}
{\( h(x)=\frac{2x-4}{x-2} \)}{\( m(x)=\frac{x^2-4}{x-2} \)}{\( g(x)=\frac{4-x^2}{x-2} \)}{\( f(x)= -\left|\frac{x^2-4}{x-2}\right| \)}{}{}}
\LANGsk{
\Question{
Ktorá z nasledujúcich funkcií je zdola ohraničená?}
{\( h(x)=\frac{2x-4}{x-2} \)}{\( m(x)=\frac{x^2-4}{x-2} \)}{\( g(x)=\frac{4-x^2}{x-2} \)}{\( f(x)= -\left|\frac{x^2-4}{x-2}\right| \)}{}{}}
\LANGes{
\Question{
Identifica cuál de las siguientes funciones está acotada inferiormente.}
{\( h(x)=\frac{2x-4}{x-2} \)}{\( m(x)=\frac{x^2-4}{x-2} \)}{\( g(x)=\frac{4-x^2}{x-2} \)}{\( f(x)= -\left|\frac{x^2-4}{x-2}\right| \)}{}{}}
\End

\Data\ID{1003118304}
\Updated{2020-07-15 03:07:21}
\Level{C}
\SubArea{Rational functions}
\LANGen{
\Question{
Identify which of the following functions is bounded.}
{\( h(x)=\frac{3x-6}{2x-4} \)}{\( f(x)=\frac{3x-6}{2x} \)}{\( g(x)=3-\frac6{2x} \)}{\( m(x)=\left|\frac{4x-3}{2x-6}\right| \)}{}{}}
\LANGpl{
\Question{
Która z poniższych funkcji jest ograniczona?}
{\( h(x)=\frac{3x-6}{2x-4} \)}{\( f(x)=\frac{3x-6}{2x} \)}{\( g(x)=3-\frac6{2x} \)}{\( m(x)=\left|\frac{4x-3}{2x-6}\right| \)}{}{}}
\LANGcs{
\Question{
Která z následujících funkcí je omezená?}
{\( h(x)=\frac{3x-6}{2x-4} \)}{\( f(x)=\frac{3x-6}{2x} \)}{\( g(x)=3-\frac6{2x} \)}{\( m(x)=\left|\frac{4x-3}{2x-6}\right| \)}{}{}}
\LANGsk{
\Question{
Určte, ktorá nasledujúcich funkcií je ohraničená.}
{\( h(x)=\frac{3x-6}{2x-4} \)}{\( f(x)=\frac{3x-6}{2x} \)}{\( g(x)=3-\frac6{2x} \)}{\( m(x)=\left|\frac{4x-3}{2x-6}\right| \)}{}{}}
\LANGes{
\Question{
Identifica cuál de las siguintes funciones está acotada.}
{\( h(x)=\frac{3x-6}{2x-4} \)}{\( f(x)=\frac{3x-6}{2x} \)}{\( g(x)=3-\frac6{2x} \)}{\( m(x)=\left|\frac{4x-3}{2x-6}\right| \)}{}{}}
\End

\Data\ID{1003118305}
\Updated{2020-07-15 03:07:21}
\Level{C}
\SubArea{Rational functions}
\LANGen{
\Question{
Find the false statement about the function \( f(x)=\left|\frac1{2-3x}-3\right| \).}
{The domain of the function \( f \) is the set \( \left(-\infty;\frac32\right)\cup\left(\frac32;\infty\right) \).}{The range of the function \( f \) is the interval \( \left[0;\infty\right) \).}{The function \( f \) has the minimum at \( x=\frac59 \).}{The function \( f \) is bounded below.}{}{}}
\LANGpl{
\Question{
Wybierz fałszywe stwierdzenie dotyczące funkcji \( f(x)=\left|\frac1{2-3x}-3\right| \).}
{Dziedziną funkcji \( f \) jest zbiór \( \left(-\infty;\frac32\right)\cup\left(\frac32;\infty\right) \).}{Zakres funkcji  \( f \) mieści się w przedziale \( \left\langle0;\infty\right) \).}{Funkcja  \( f \) osiąga minimum w \( x=\frac59 \).}{Funkcja \( f \) jest ograniczona z dołu.}{}{}}
\LANGcs{
\Question{
Funkce \( f \) je dána předpisem \( f(x)=\left|\frac1{2-3x}-3\right| \). Vyberte nepravdivý výrok.}
{Definiční obor funkce \( f \) je množina \( \left(-\infty;\frac32\right)\cup\left(\frac32;\infty\right) \).}{Obor hodnot funkce  \( f \) je interval \( \left\langle0;\infty\right) \).}{Funkce \( f \) má minimum v bodě \( x=\frac59 \).}{Funkce \( f \) je zdola omezená.}{}{}}
\LANGsk{
\Question{
Vyberte nepravdivý výrok o funkcií \( f(x)=\left|\frac1{2-3x}-3\right| \).}
{Definičným oborom funkcie \( f \) je množina \( \left(-\infty;\frac32\right)\cup\left(\frac32;\infty\right) \).}{Oborom hodnôt funkcie \( f \) je interval \( \left\langle0;\infty\right) \).}{Funkcia \( f \) má minimum v bode \( x=\frac59 \).}{Funkcia \( f \) je zdola ohraničená.}{}{}}
\LANGes{
\Question{
Determina la proposición  falsa sobre la función\( f(x)=\left|\frac1{2-3x}-3\right| \).}
{El dominio de la función  \( f \) es el conjunto\( \left(-\infty;\frac32\right)\cup\left(\frac32;\infty\right) \).}{El rango de la función \( f \) es el intervalo \( \left[0;\infty\right) \).}{La función \( f \) tiene su mínimo en \( x=\frac59 \).}{La función \( f \) está acotada inferiormente.}{}{}}
\End

\Data\ID{1003118306}
\Updated{2021-04-09 07:04:41}
\Level{C}
\SubArea{Rational functions}
\LANGen{
\Question{
Find the true statement about the function \( f(x)=\left|\frac{4x-4}{2x-1}\right| \).}
{The domain of the function \( f \) is the set \( \left(-\infty;\frac12\right)\cup\left(\frac12;\infty\right) \).}{The range of the function \( f \) is the set \( [0;2)\cup(2;\infty) \).}{The function \( f \) has the minimum at \( x=4 \).}{The function \( f \) is an injective (one-to-one) function.}{}{}}
\LANGpl{
\Question{
Wybierz prawdziwe stwierdzenie dotyczące funkcji \( f(x)=\left|\frac{4x-4}{2x-1}\right| \).}
{Dziedziną funkcji \( f \) jest zbiór  \( \left(-\infty;\frac12\right)\cup\left(\frac12;\infty\right) \).}{Zakresem funkcji  \( f \) jest zbiór \( \langle0;2)\cup(2;\infty) \).}{Funkcja \( f \) osiąga minimum w  \( x=4 \).}{Funkcja \( f \) jest funkcją iniekcyjną (jeden do jednego).}{}{}}
\LANGcs{
\Question{
Funkce \( f \) je dána předpisem  \( f(x)=\left|\frac{4x-4}{2x-1}\right| \). Vyberte pravdivý výrok.}
{Definiční obor funkce \( f \) je množina \( \left(-\infty;\frac12\right)\cup\left(\frac12;\infty\right) \).}{Obor hodnot funkce  \( f \) je množina \( \langle0;2)\cup(2;\infty) \).}{Funkce \( f \) má minimum v bodě \( x=4 \).}{Funkce \( f \) je prostá.}{}{}}
\LANGsk{
\Question{
Vyberte pravdivý výrok o funkcií \( f(x)=\left|\frac{4x-4}{2x-1}\right| \).}
{Definičným oborom funkcie \( f \) je množina \( \left(-\infty;\frac12\right)\cup\left(\frac12;\infty\right) \).}{Oborom hodnôt funkcie \( f \) je množina \( \langle0;2)\cup(2;\infty) \).}{Funkcia \( f \) má minimum v bode \( x=4 \).}{Funkcia \( f \) je prostá.}{}{}}
\LANGes{
\Question{
Determina la proposición verdadera sobre la función \( f(x)=\left|\frac{4x-4}{2x-1}\right| \).}
{El dominio de la función \( f \) es el conjunto \( \left(-\infty;\frac12\right)\cup\left(\frac12;\infty\right) \).}{El rango de la función \( f \) es el conjunto \( [0;2)\cup(2;\infty) \).}{La función \( f \) tiene su mínimo en \( x=4 \).}{La función \( f \) es una función inyectiva (uno a uno).}{}{}}
\End

\Data\ID{1003118307}
\Updated{2020-07-15 03:07:21}
\Level{C}
\SubArea{Rational functions}
\LANGen{
\Question{
Identify which of the following functions has the maximum at \( x=-\frac12 \).}
{\( m(x)=-\left|\frac{4x+2}{x-2}\right| \)}{\( g(x)=\left|-\frac{5x+10}{2x-1}\right| \)}{\( f(x)=-\left|\frac{2x+1}{4x+2}\right| \)}{\( h(x)=-\left|\frac{x+1}{2x-2}\right| \)}{}{}}
\LANGpl{
\Question{
Która z poniższych funkcji osiąga maksimum w \( x=-\frac12 \)?}
{\( m(x)=-\left|\frac{4x+2}{x-2}\right| \)}{\( g(x)=\left|-\frac{5x+10}{2x-1}\right| \)}{\( f(x)=-\left|\frac{2x+1}{4x+2}\right| \)}{\( h(x)=-\left|\frac{x+1}{2x-2}\right| \)}{}{}}
\LANGcs{
\Question{
Která z následujících funkcí má maximum v bodě \( x=-\frac12 \)?}
{\( m(x)=-\left|\frac{4x+2}{x-2}\right| \)}{\( g(x)=\left|-\frac{5x+10}{2x-1}\right| \)}{\( f(x)=-\left|\frac{2x+1}{4x+2}\right| \)}{\( h(x)=-\left|\frac{x+1}{2x-2}\right| \)}{}{}}
\LANGsk{
\Question{
Určte, ktorá z daných funkcií má maximum v bode \( x=-\frac12 \).}
{\( m(x)=-\left|\frac{4x+2}{x-2}\right| \)}{\( g(x)=\left|-\frac{5x+10}{2x-1}\right| \)}{\( f(x)=-\left|\frac{2x+1}{4x+2}\right| \)}{\( h(x)=-\left|\frac{x+1}{2x-2}\right| \)}{}{}}
\LANGes{
\Question{
Identifica cuál de las siguintes funciones tiene su máximo en \( x=-\frac12 \).}
{\( m(x)=-\left|\frac{4x+2}{x-2}\right| \)}{\( g(x)=\left|-\frac{5x+10}{2x-1}\right| \)}{\( f(x)=-\left|\frac{2x+1}{4x+2}\right| \)}{\( h(x)=-\left|\frac{x+1}{2x-2}\right| \)}{}{}}
\End

\Data\ID{1103082701}
\Updated{2021-04-09 07:04:07}
\Level{C}
\SubArea{Rational functions}
\IMG{1103082701.pdf}
\LANGen{
\Question{
Function \( f \)  is given completely by the graph. Identify which of the following statements is false.}
{\( f(x)=\frac1x;\ x\in[-2;-0.5] \)}{\( f(x)=\left|-\frac1x\right|;\ x\in[-2;-0.5] \)}{\( f(x)=\frac1{|x|} ;\ x\in[-2;-0.5]  \)}{\( f(x)=-\frac1x;\ x\in[-2;-0.5]  \)}{}{}}
\LANGpl{
\Question{
Funkcję \( f \)  przedstawiono na wykresie poniżej. Które z poniższych zdań jest fałszywe?}
{\( f(x)=\frac1x;\ x\in\langle-2;-0{,}5\rangle \)}{\( f(x)=\left|-\frac1x\right|;\ x\in\langle-2;-0{,}5\rangle \)}{\( f(x)=\frac1{|x|} ;\ x\in\langle-2;-0{,}5\rangle  \)}{\( f(x)=-\frac1x;\ x\in\langle-2;-0{,}5\rangle  \)}{}{}}
\LANGcs{
\Question{
Funkce \( f \)  je dána grafem. Určete, které z následujících tvrzení je nepravdivé.}
{\( f(x)=\frac1x;\ x\in\langle-2;-0{,}5\rangle \)}{\( f(x)=\left|-\frac1x\right|;\ x\in\langle-2;-0{,}5\rangle \)}{\( f(x)=\frac1{|x|} ;\ x\in\langle-2;-0{,}5\rangle  \)}{\( f(x)=-\frac1x;\ x\in\langle-2;-0{,}5\rangle  \)}{}{}}
\LANGsk{
\Question{
Funkcia \( f \) je daná grafom. Vyberte nepravdivé tvrdenie o funkcii \( f \).}
{\( f(x)=\frac1x;\ x\in\langle-2;-0{,}5\rangle \)}{\( f(x)=\left|-\frac1x\right|;\ x\in\langle-2;-0{,}5\rangle \)}{\( f(x)=\frac1{|x|} ;\ x\in\langle-2;-0{,}5\rangle  \)}{\( f(x)=-\frac1x;\ x\in\langle-2;-0{,}5\rangle  \)}{}{}}
\LANGes{
\Question{
La función  \( f \)   está definida completamente por la gráfica. ¿Cuál de las siguientes proposiciones es falsa?}
{\( f(x)=\frac1x;\ x\in[-2;-0.5] \)}{\( f(x)=\left|-\frac1x\right|;\ x\in[-2;-0.5] \)}{\( f(x)=\frac1{|x|} ;\ x\in[-2;-0.5]  \)}{\( f(x)=-\frac1x;\ x\in[-2;-0.5]  \)}{}{}}
\End

\Data\ID{1103102304}
\Updated{2021-04-09 07:04:32}
\Level{C}
\SubArea{Rational functions}
\IMG{1103102304.pdf}
\LANGen{
\Question{
Function \( f \) is given completely by the graph. Which of the following statements is false?}
{\(  f(x)=\frac{|x|}x,\ x\in[-5;0)\cup(0;5] \)}{\(  f(x)=\left|\frac{|x|}x\right|,\ x\in[-5;0)\cup(0;5] \)}{\(  f(x)=1,\ x\in[-5;0)\cup(0;5] \)}{\(  f(x)=\frac{x}x,\ x\in[-5;0)\cup(0;5] \)}{}{}}
\LANGpl{
\Question{
Funkcję \( f \) przedstawiono za pomocą wykresu. Które z poniższych stwierdzeń jest fałszywe?}
{\(  f(x)=\frac{|x|}x,\ x\in\langle-5;0)\cup(0;5\rangle \)}{\(  f(x)=\left|\frac{|x|}x\right|,\ x\in\langle-5;0)\cup(0;5\rangle \)}{\(  f(x)=1,\ x\in\langle-5;0)\cup(0;5\rangle \)}{\(  f(x)=\frac{x}x,\ x\in\langle-5;0)\cup(0;5\rangle \)}{}{}}
\LANGcs{
\Question{
Funkce \( f \) je dána uvedeným grafem. Vyberte nepravdivý výrok.}
{\(  f(x)=\frac{|x|}x,\ x\in\langle-5;0)\cup(0;5\rangle \)}{\(  f(x)=\left|\frac{|x|}x\right|,\ x\in\langle-5;0)\cup(0;5\rangle \)}{\(  f(x)=1,\ x\in\langle-5;0)\cup(0;5\rangle \)}{\(  f(x)=\frac{x}x,\ x\in\langle-5;0)\cup(0;5\rangle \)}{}{}}
\LANGsk{
\Question{
Funkcia \( f \) je daná grafom. Vyberte nepravdivé tvrdenie o funkcii \( f \).}
{\(  f(x)=\frac{|x|}x,\ x\in\langle-5;0)\cup(0;5\rangle \)}{\(  f(x)=\left|\frac{|x|}x\right|,\ x\in\langle-5;0)\cup(0;5\rangle \)}{\(  f(x)=1,\ x\in\langle-5;0)\cup(0;5\rangle \)}{\(  f(x)=\frac{x}x,\ x\in\langle-5;0)\cup(0;5\rangle \)}{}{}}
\LANGes{
\Question{
La función \( f \) está definida completamente por la gráfica. ¿Cuál de las siguientes proposiciones es  falsa?}
{\(  f(x)=\frac{|x|}x,\ x\in[-5;0)\cup(0;5] \)}{\(  f(x)=\left|\frac{|x|}x\right|,\ x\in[-5;0)\cup(0;5] \)}{\(  f(x)=1,\ x\in[-5;0)\cup(0;5] \)}{\(  f(x)=\frac{x}x,\ x\in[-5;0)\cup(0;5] \)}{}{}}
\End

\Data\ID{1103124502}
\Updated{2020-07-15 03:07:10}
\Level{C}
\SubArea{Rational functions}
\LANGen{
\Question{
Identify which of the graphs below represents the function \( f(x)=\left|\frac{1-2x}{x-4}\right|;\ x\in [-\frac52;\frac52] \).}
{}{}{}{}{}{}}
\LANGpl{
\Question{
Który z poniższych wykresów przedstawia funkcję \( f(x)=\left|\frac{1-2x}{x-4}\right|;\ x\in\langle-\frac52;\frac52\rangle \)?}
{}{}{}{}{}{}}
\LANGcs{
\Question{
Na kterém obrázku je graf funkce \( f(x)=\left|\frac{1-2x}{x-4}\right|;\ x\in\langle-\frac52;\frac52\rangle \)?}
{}{}{}{}{}{}}
\LANGsk{
\Question{
Ktorý z nasledujúcich grafov je graf funkcie \( f(x)=\left|\frac{1-2x}{x-4}\right|;\ x\in\langle-\frac52;\frac52\rangle \)?}
{}{}{}{}{}{}}
\LANGes{
\Question{
Identifica cuál de las gráficas representa la función  \( f(x)=\left|\frac{1-2x}{x-4}\right|;\ x\in [-\frac52;\frac52] \).}
{}{}{}{}{}{}}

\IMGanswer{1103124502a_1.pdf}
\IMGanswer{1103124502b_0.pdf}
\IMGanswer{1103124502c_0.pdf}
\IMGanswer{1103124502d_0.pdf}\End

\Data\ID{1103124602}
\Updated{2022-05-29 12:05:45}
\Level{C}
\SubArea{Rational functions}
\LANGen{
\Question{
Let \( f(x)=\frac{x^2-x-6}{x^2-9} \). One of the following pictures  shows a part of the graph of \( f \).  Choose the picture.}
{}{}{}{}{}{}}
\LANGpl{
\Question{
Niech  \( f(x)=\frac{x^2-x-6}{x^2-9} \). Jeden z poniższych rysunków przedstawia część wykresu funkcji  \( f \).  Wybierz ten rysunek.}
{}{}{}{}{}{}}
\LANGcs{
\Question{
Na kterém obrázku je část grafu funkce \( f(x)=\frac{x^2-x-6}{x^2-9} \)?}
{}{}{}{}{}{}}
\LANGsk{
\Question{
Daná je funkcia \( f(x)=\frac{x^2-x-6}{x^2-9} \). Ktorý z nasledujúcich obrázkov predstavuje časť grafu funkcie \( f \)?}
{}{}{}{}{}{}}
\LANGes{
\Question{
Sea \( f(x)=\frac{x^2-x-6}{x^2-9} \). Una de las siguientes imágenes muestra una parte de la gráfica de la función \( f \).  Elige la imágen.}
{}{}{}{}{}{}}

\IMGanswer{1103124602a.pdf}
\IMGanswer{1103124602b.pdf}
\IMGanswer{1103124602c.pdf}
\IMGanswer{1103124602d.pdf}\End

\Data\ID{2010009904}
\Updated{2022-08-25 10:08:18}
\Level{C}
\SubArea{Rational functions}
\IMG{2010009901-figure0_0.pdf}
\LANGen{
\Question{
A part of the graph of the function \( f(x)=\frac{-3}x \)  is shown in the picture. Identify which of the following statements is true.}
{The function \( g \) defined by \( g(x)=-\left|f(x)\right| \) is bounded above.}{The function \( m \) defined by \( m(x)=\left|f(x)\right| \) is bounded above.}{The function \( h \) defined by \( h(x)=-f(x)\) is bounded below.}{The function \( f \) is bounded below.}{}{}}
\LANGpl{
\Question{
Na rysunku przedstawiono część funkcji \( f(x)=\frac{-3}x \). Wskaż prawdziwe stwierdzenie.}
{Funkcja \( g \) określona jako \( g(x)=-\left|f(x)\right| \) jest ograniczona z góry.}{Funkcja \( m \) określona jako \( m(x)=\left|f(x)\right| \) jest ograniczona z dołu.}{Funkcja \( h \) określona jako \( h(x)=-f(x)\) jest ograniczona z dołu.}{Funkcja \( f \) jest ograniczona z dołu.}{}{}}
\LANGcs{
\Question{
Na obrázku je část grafu funkce \( f(x)=\frac{-3}x \). Vyberte pravdivý výrok.}
{Funkce \( g \) definovaná vztahem \( g(x)=-\left|f(x)\right| \) je shora omezená.}{Funkce \( m \) definovaná vztahem \( m(x)=\left|f(x)\right| \) je shora omezená.}{Funkce \( h \) defined by \( h(x)=-f(x)\) je zdola omezená.}{Funkce \( f \) je zdola omezená.}{}{}}
\LANGsk{
\Question{
Na obrázku je časť grafu funkcie \( f(x)=\frac{-3}x \). Vyberte pravdivý výrok.}
{Funkcia \( g \) definovaná vzťahom \( g(x)=-\left|f(x)\right| \) je zhora ohraničená.}{Funkcia \( m \) definovaná vzťahom \( m(x)=\left|f(x)\right| \) je zhora ohraničená.}{Funkcia \( h \) definovaná vzťahom \( h(x)=-f(x)\) je zdola ohraničená.}{Funkcia \( f \) je zdola ohraničená.}{}{}}
\LANGes{
\Question{
Una parte de la gráfica de la función \( f(x)=\frac{-3}x \) se muestra en la imagen. Identifica cuál de las siguientes proposiciones es verdadera.}
{La función \( g \) definida como  \( g(x)=-\left|f(x)\right| \) está acotada superiormente.}{La función \( m \) definida como  \( m(x)=\left|f(x)\right| \) está acotada superiormente.}{La función \( h \) definida como  \( h(x)=-f(x)\) está acotada inferiormente.}{La función \( f \) está acotada inferiormente.}{}{}}
\End

\Data\ID{2010017301}
\Updated{2022-08-25 10:08:47}
\Level{C}
\SubArea{Rational functions}
\LANGen{
\Question{
Identify which of the given functions is odd.}
{\( m(x)=\frac{x^2-4}{x^3} \)}{\( h(x)=\frac{|x|-1}{x^2} \)}{\( g(x)=2x^3-16 \)}{\( f(x)=\frac{x^2}{x^3-1} \)}{}{}}
\LANGpl{
\Question{
Wskaż funkcję nieparzystą.}
{\( m(x)=\frac{x^2-4}{x^3} \)}{\( h(x)=\frac{|x|-1}{x^2} \)}{\( g(x)=2x^3-16 \)}{\( f(x)=\frac{x^2}{x^3-1} \)}{}{}}
\LANGcs{
\Question{
Určete, která z následujících funkcí je lichá.}
{\( m(x)=\frac{x^2-4}{x^3} \)}{\( h(x)=\frac{|x|-1}{x^2} \)}{\( g(x)=2x^3-16 \)}{\( f(x)=\frac{x^2}{x^3-1} \)}{}{}}
\LANGsk{
\Question{
Určte, ktorá z nasledujúcich funkcií je nepárna.}
{\( m(x)=\frac{x^2-4}{x^3} \)}{\( h(x)=\frac{|x|-1}{x^2} \)}{\( g(x)=2x^3-16 \)}{\( f(x)=\frac{x^2}{x^3-1} \)}{}{}}
\LANGes{
\Question{
Identifica cuál de las funciones dadas es impar.}
{\( m(x)=\frac{x^2-4}{x^3} \)}{\( h(x)=\frac{|x|-1}{x^2} \)}{\( g(x)=2x^3-16 \)}{\( f(x)=\frac{x^2}{x^3-1} \)}{}{}}
\End

\Data\ID{2010017302}
\Updated{2022-08-25 10:08:47}
\Level{C}
\SubArea{Rational functions}
\IMG{2010017301-figure0.pdf}
\LANGen{
\Question{
Find the interval where the function \(f(x) = -\left |2+\frac{1} 
{x}\right |\) is a decreasing
function. The function \(f\) 
is graphed in the picture.}
{\(\left[ -\frac12; 0\right)\)}{\((-\infty ;0)\)}{\(\left[ -\frac12; \infty\right)\)}{\(\left(-\infty ; -\frac12\right)\)}{}{}}
\LANGpl{
\Question{
Znajdź przedział, w którym funkcja \(f(x) = -\left |2+\frac{1} 
{x}\right |\) jest funkcją malejącą. Funkcja \(f\) 
jest przedstawiona na rysunku.}
{\(\left\langle -\frac12; 0\right)\)}{\((-\infty ;0)\)}{\(\left\langle -\frac12; \infty\right)\)}{\(\left(-\infty ; -\frac12\right)\)}{}{}}
\LANGcs{
\Question{
Určete interval, v němž je funkce  \(f(x) = -\left |2+\frac{1} 
{x}\right |\) klesající.
Graf funkce \(f\) 
je znázorněn na obrázku.}
{\(\left\langle -\frac12; 0\right)\)}{\((-\infty ;0)\)}{\(\left\langle -\frac12; \infty\right)\)}{\(\left(-\infty ; -\frac12\right)\)}{}{}}
\LANGsk{
\Question{
Určte interval, v ktorom je funkcia  \(f(x) = -\left |2+\frac{1} 
{x}\right |\) klesajúca.
Graf funkcie \(f\) 
je znázornený na obrázku.}
{\(\left\langle -\frac12; 0\right)\)}{\((-\infty ;0)\)}{\(\left\langle -\frac12; \infty\right)\)}{\(\left(-\infty ; -\frac12\right)\)}{}{}}
\LANGes{
\Question{
Determina el intervalo donde la función \(f(x) = -\left |2+\frac{1} 
{x}\right |\) es decreciente. La gráfica de la función \(f\) 
está representada en la imagen.}
{\(\left[ -\frac12; 0\right)\)}{\((-\infty ;0)\)}{\(\left[ -\frac12; \infty\right)\)}{\(\left(-\infty ; -\frac12\right)\)}{}{}}
\End

\Data\ID{2010017304}
\Updated{2022-08-25 10:08:47}
\Level{C}
\SubArea{Rational functions}
\LANGen{
\Question{
Consider the functions 
\[ 
 \text{$f(x)= -\frac{2} 
{3x}$ and $g(x) = \frac{k} 
{x}$.}
\]
Identify the value of the coefficient \(k\) 
which ensures that the graphs of both functions are symmetric about
\(y\)-axis.}
{\( k=\frac23\)}{\( k=\frac32\)}{\( k=-\frac23\)}{\( k=-\frac32\)}{}{}}
\LANGpl{
\Question{
Dane są dwie funkcje 
\[ 
 \text{$f(x)= -\frac{2} 
{3x}$ i $g(x) = \frac{k} 
{x}$.}
\]
Określ wartość współczynnika \(k\), dla którego wykresy obu funkcji są symetryczne względem osi
\(y\).}
{\( k=\frac23\)}{\( k=\frac32\)}{\( k=-\frac23\)}{\( k=-\frac32\)}{}{}}
\LANGcs{
\Question{
Jsou dány funkce  
\[ 
 \text{$f(x)= -\frac{2} 
{3x}$   a   $g(x) = \frac{k} 
{x}$.}
\]
Určete hodnotu koeficientu  \(k\) 
tak, aby grafy obou funkcí byly symetrické podle osy
\(y\).}
{\( k=\frac23\)}{\( k=\frac32\)}{\( k=-\frac23\)}{\( k=-\frac32\)}{}{}}
\LANGsk{
\Question{
Sú dané funkcie  
\[ 
 \text{$f(x)= -\frac{2} 
{3x}$   a   $g(x) = \frac{k} 
{x}$.}
\]
Určte hodnotu koeficientu  \(k\) 
tak, aby grafy oboch funkcií boli symetrické podľa osi
\(y\).}
{\( k=\frac23\)}{\( k=\frac32\)}{\( k=-\frac23\)}{\( k=-\frac32\)}{}{}}
\LANGes{
\Question{
Considera las funciones
\[ 
 \text{$f(x)= -\frac{2} 
{3x}$ y $g(x) = \frac{k} 
{x}$.}
\]
Identifica el valor del coeficiente \(k\) 
que asegura que las gráficas de ambas funciones sean simétricas respecto al eje
\(y\)}
{\( k=\frac23\)}{\( k=\frac32\)}{\( k=-\frac23\)}{\( k=-\frac32\)}{}{}}
\End

\Data\ID{2010017305}
\Updated{2022-08-25 10:08:47}
\Level{C}
\SubArea{Rational functions}
\IMG{2010017301-figure5_0.pdf}
\LANGen{
\Question{
The picture shows parts of the graphs of the functions 
\[ 
 \text{$f(x)= \frac{k_{1}} 
 {x} $ and $g(x) = \frac{k_{2}} 
 {x} $.}
\]
Find the relationship between \(k_{1}\) 
and \(k_{2}\)?}
{\( k_1 < k_2\)}{\( k_1 \geq k_2\)}{\( k_1 = k_2\)}{The relationship between \(k_1\) and \(k_2\) cannot be determined from the picture.}{}{}}
\LANGpl{
\Question{
Rysunek przedstawia fragmenty wykresów funkcji  
\[ 
 \text{$f(x)= \frac{k_{1}} 
 {x} $ i $g(x) = \frac{k_{2}} 
 {x} $.}
\]
Jaka jest zależność pomiędzy \(k_{1}\) 
a \(k_{2}\)?}
{\( k_1 < k_2\)}{\( k_1 \geq k_2\)}{\( k_1 = k_2\)}{Zależności pomiędzy \(k_1\) a \(k_2\) nie można określić na podstawie rysunku.}{}{}}
\LANGcs{
\Question{
Na obrázku jsou části grafů funkcí  
\[ 
 \text{$f(x)= \frac{k_{1}} 
 {x} $   a   $g(x) = \frac{k_{2}} 
 {x} $.}
\]
Určete vztah mezi koeficienty \(k_{1}\) 
a \(k_{2}\).}
{\( k_1 < k_2\)}{\( k_1 \geq k_2\)}{\( k_1 = k_2\)}{Vztah mezi  \(k_1\) a \(k_2\) není možné z obrázku určit.}{}{}}
\LANGsk{
\Question{
Na obrázku sú časti grafov funkcií  
\[ 
 \text{$f(x)= \frac{k_{1}} 
 {x} $   a   $g(x) = \frac{k_{2}} 
 {x} $.}
\]
Určte vzťah medzi koeficientmi \(k_{1}\) 
a \(k_{2}\).}
{\( k_1 < k_2\)}{\( k_1 \geq k_2\)}{\( k_1 = k_2\)}{Vzťah medzi \(k_1\) a \(k_2\) sa nedá z obrázka určiť.}{}{}}
\LANGes{
\Question{
La imagen muestra partes de las gráficas de las funciones.
\[ 
 \text{$f(x)= \frac{k_{1}} 
 {x} $ y $g(x) = \frac{k_{2}} 
 {x} $.}
\]
Determina la relación entre\(k_{1}\) 
y \(k_{2}\).}
{\( k_1 < k_2\)}{\( k_1 \geq k_2\)}{\( k_1 = k_2\)}{La relación entre \(k_1\) y \(k_2\) no se puede determinar a base a la imagen.}{}{}}
\End

\Data\ID{2110017303}
\Updated{2022-08-25 10:08:02}
\Level{C}
\SubArea{Rational functions}
\LANGen{
\Question{
Let \( f(x)=\frac{x^2-1}{x^2-3x+2} \). One of the following pictures  shows a part of the graph of \( f \).  Choose the picture.}
{}{}{}{}{}{}}
\LANGpl{
\Question{
Dana jest funkcja \( f(x)=\frac{x^2-1}{x^2-3x+2} \). Wskaż, który z poniższych rysunków przedstawia fragment wykresu funkcji \( f \).}
{}{}{}{}{}{}}
\LANGcs{
\Question{
Je dána funkce  \( f(x)=\frac{x^2-1}{x^2-3x+2} \). Určete, který z obrázků znázorňuje část grafu funkce  \( f \).}
{}{}{}{}{}{}}
\LANGsk{
\Question{
Je daná funkcia  \( f(x)=\frac{x^2-1}{x^2-3x+2} \). Určte, ktorý z obrázkov znázorňuje časť grafu funkcie  \( f \).}
{}{}{}{}{}{}}
\LANGes{
\Question{
Sea \( f(x)=\frac{x^2-1}{x^2-3x+2}\). Una de las siguientes imágenes muestra una parte de la gráfica de la función \( f \).  Elige la imágen.}
{}{}{}{}{}{}}

\IMGanswer{2010017301-figure1.pdf}
\IMGanswer{2010017301-figure2.pdf}
\IMGanswer{2010017301-figure3.pdf}
\IMGanswer{2010017301-figure4.pdf}\End

\Data\ID{9000002908}
\Updated{2022-04-23 05:04:17}
\Level{C}
\SubArea{Rational functions}
\IMG{000029_08-figure0.pdf}
\LANGen{
\Question{
Find the intervals where the function
\(f(x) = \left |1 + \frac{1} 
{x}\right |\) is an increasing
function. The function \(f\) 
is graphed in the picture.}
{\([ - 1;0)\)}{\((-\infty ;1] \)}{\((-\infty ;0)\)}{\((0;\infty )\)}{}{}}
\LANGpl{
\Question{
Określ przedziały, w których  funkcja
\(f\colon y = \left |1 + \frac{1} 
{x}\right |\) jest funkcją rosnącą.  Wykres funkcji \(f\) 
przedstawiono na rysunku poniżej.}
{\([ - 1;0)\)}{\((-\infty ;1] \)}{\((-\infty ;0)\)}{\((0;\infty )\)}{}{}}
\LANGcs{
\Question{
Z grafu na obrázku určete interval, v němž je funkce
\(f\colon y = \left |1 + \frac{1} 
{x}\right |\) 
rostoucí.}
{\(\langle - 1;0)\)}{\((-\infty ;1\rangle \)}{\((-\infty ;0)\)}{\((0;\infty )\)}{}{}}
\LANGsk{
\Question{
Z grafu na obrázku určte interval, na ktorom daná funkcia
\(f\colon y = \left |1 + \frac{1} 
{x}\right |\) rastie.}
{\(\langle - 1;0)\)}{\((-\infty ;1\rangle \)}{\((-\infty ;0)\)}{\((0;\infty )\)}{}{}}
\LANGes{
\Question{
Determina los intervalos donde la función
\(f(x) = \left |1 + \frac{1} 
{x}\right |\) es creciente. La gráfica de la función  \(f\) 
está representada en la imagen.}
{\([ - 1;0)\)}{\((-\infty ;1] \)}{\((-\infty ;0)\)}{\((0;\infty )\)}{}{}}
\End

\Data\ID{9000007603}
\Updated{2020-07-15 03:07:34}
\Level{C}
\SubArea{Rational functions}
\LANGen{
\Question{
Find the domain of the function \(f(x) = 1 + \left | \frac{1}
{2(x-2)}\right |\).}
{\(\mathbb{R}\setminus \{2\}\)}{\(\mathbb{R}\setminus \{ - 2\}\)}{\(\mathbb{R}\setminus \{4\}\)}{\(\mathbb{R}\setminus \{1\}\)}{\(\mathbb{R}\)}{}}
\LANGpl{
\Question{
Znajdź dziedzinę funkcji \(f(x) = 1 + \left | \frac{1}
{2(x-2)}\right |\).}
{\(\mathbb{R}\setminus \{2\}\)}{\(\mathbb{R}\setminus \{ - 2\}\)}{\(\mathbb{R}\setminus \{4\}\)}{\(\mathbb{R}\setminus \{1\}\)}{\(\mathbb{R}\)}{}}
\LANGcs{
\Question{
Určete definiční obor funkce \(f(x) = 1 + \left | \frac{1}
{2(x-2)}\right |\).}
{\(\mathbb{R}\setminus \{2\}\)}{\(\mathbb{R}\setminus \{ - 2\}\)}{\(\mathbb{R}\setminus \{4\}\)}{\(\mathbb{R}\setminus \{1\}\)}{\(\mathbb{R}\)}{}}
\LANGsk{
\Question{
Určte definičný obor funkcie \(f(x) = 1 + \left | \frac{1}
{2(x-2)}\right |\).}
{\(\mathbb{R}\setminus \{2\}\)}{\(\mathbb{R}\setminus \{ - 2\}\)}{\(\mathbb{R}\setminus \{4\}\)}{\(\mathbb{R}\setminus \{1\}\)}{\(\mathbb{R}\)}{}}
\LANGes{
\Question{
Determina el dominio de la función \(f(x) = 1 + \left | \frac{1}
{2(x-2)}\right |\).}
{\(\mathbb{R}\setminus \{2\}\)}{\(\mathbb{R}\setminus \{ - 2\}\)}{\(\mathbb{R}\setminus \{4\}\)}{\(\mathbb{R}\setminus \{1\}\)}{\(\mathbb{R}\)}{}}
\End

\Data\ID{9000007604}
\Updated{2020-07-15 03:07:34}
\Level{C}
\SubArea{Rational functions}
\LANGen{
\Question{
Find the domain of the function \(f(x) = 1 + \left | \frac{1}
{|x|+1}\right |\).}
{\(\mathbb{R}\)}{\(\mathbb{R}\setminus \{ - 1\}\)}{\(\mathbb{R}\setminus \{ - 1;1\}\)}{\(\mathbb{R}\setminus \{ - 1;0;1\}\)}{\(\mathbb{R}\setminus \{1\}\)}{}}
\LANGpl{
\Question{
Znajdź dziedzinę funkcji  \(f(x) = 1 + \left | \frac{1}
{|x|+1}\right |\).}
{\(\mathbb{R}\)}{\(\mathbb{R}\setminus \{ - 1\}\)}{\(\mathbb{R}\setminus \{ - 1;1\}\)}{\(\mathbb{R}\setminus \{ - 1;0;1\}\)}{\(\mathbb{R}\setminus \{1\}\)}{}}
\LANGcs{
\Question{
Určete definiční obor funkce \(f(x) = 1 + \left | \frac{1}
{|x|+1}\right |\).}
{\(\mathbb{R}\)}{\(\mathbb{R}\setminus \{ - 1\}\)}{\(\mathbb{R}\setminus \{ - 1;1\}\)}{\(\mathbb{R}\setminus \{ - 1;0;1\}\)}{\(\mathbb{R}\setminus \{1\}\)}{}}
\LANGsk{
\Question{
Určte definičný obor funkcie \(f(x) = 1 + \left | \frac{1}
{|x|+1}\right |\).}
{\(\mathbb{R}\)}{\(\mathbb{R}\setminus \{ - 1\}\)}{\(\mathbb{R}\setminus \{ - 1;1\}\)}{\(\mathbb{R}\setminus \{ - 1;0;1\}\)}{\(\mathbb{R}\setminus \{1\}\)}{}}
\LANGes{
\Question{
Determina el dominio de la función \(f(x) = 1 + \left | \frac{1}
{|x|+1}\right |\).}
{\(\mathbb{R}\)}{\(\mathbb{R}\setminus \{ - 1\}\)}{\(\mathbb{R}\setminus \{ - 1;1\}\)}{\(\mathbb{R}\setminus \{ - 1;0;1\}\)}{\(\mathbb{R}\setminus \{1\}\)}{}}
\End

\Data\ID{9000007605}
\Updated{2020-07-15 03:07:34}
\Level{C}
\SubArea{Rational functions}
\LANGen{
\Question{
Find the domain of the function \(f(x) = 1 + \left | \frac{1}
{-|x|+1}\right |\).}
{\(\mathbb{R}\setminus \{ - 1;1\}\)}{\(\mathbb{R}\setminus \{ - 1\}\)}{\(\mathbb{R}\setminus \{ - 1;0;1\}\)}{\(\mathbb{R}\setminus \{1\}\)}{\(\mathbb{R}\)}{}}
\LANGpl{
\Question{
Znajdź dziedzinę funkcji  \(f(x) = 1 + \left | \frac{1}
{-|x|+1}\right |\).}
{\(\mathbb{R}\setminus \{ - 1;1\}\)}{\(\mathbb{R}\setminus \{ - 1\}\)}{\(\mathbb{R}\setminus \{ - 1;0;1\}\)}{\(\mathbb{R}\setminus \{1\}\)}{\(\mathbb{R}\)}{}}
\LANGcs{
\Question{
Určete definiční obor funkce \(f(x) = 1 + \left | \frac{1}
{-|x|+1}\right |\).}
{\(\mathbb{R}\setminus \{ - 1;1\}\)}{\(\mathbb{R}\setminus \{ - 1\}\)}{\(\mathbb{R}\setminus \{ - 1;0;1\}\)}{\(\mathbb{R}\setminus \{1\}\)}{\(\mathbb{R}\)}{}}
\LANGsk{
\Question{
Určte definičný obor funkcie \(f(x) = 1 + \left | \frac{1}
{-|x|+1}\right |\).}
{\(\mathbb{R}\setminus \{ - 1;1\}\)}{\(\mathbb{R}\setminus \{ - 1\}\)}{\(\mathbb{R}\setminus \{ - 1;0;1\}\)}{\(\mathbb{R}\setminus \{1\}\)}{\(\mathbb{R}\)}{}}
\LANGes{
\Question{
Determina el dominio de la función \(f(x) = 1 + \left | \frac{1}
{-|x|+1}\right |\).}
{\(\mathbb{R}\setminus \{ - 1;1\}\)}{\(\mathbb{R}\setminus \{ - 1\}\)}{\(\mathbb{R}\setminus \{ - 1;0;1\}\)}{\(\mathbb{R}\setminus \{1\}\)}{\(\mathbb{R}\)}{}}
\End

\Data\ID{9000007608}
\Updated{2020-07-15 03:07:34}
\Level{C}
\SubArea{Rational functions}
\LANGen{
\Question{
Find the range of the function \(f(x) = 1 + \left | \frac{1}
{2(x-2)}\right |\).}
{\((1;\infty )\)}{\(\mathbb{R}\)}{\(\mathbb{R}\setminus \{1\}\)}{\((0;\infty )\)}{\(\mathbb{R}\setminus \{2\}\)}{}}
\LANGpl{
\Question{
Znajdź zakres funkcji  \(f(x) = 1 + \left | \frac{1}
{2(x-2)}\right |\).}
{\((1;\infty )\)}{\(\mathbb{R}\)}{\(\mathbb{R}\setminus \{1\}\)}{\((0;\infty )\)}{\(\mathbb{R}\setminus \{2\}\)}{}}
\LANGcs{
\Question{
Určete obor hodnot funkce \(f(x) = 1 + \left | \frac{1}
{2(x-2)}\right |\).}
{\((1;\infty )\)}{\(\mathbb{R}\)}{\(\mathbb{R}\setminus \{1\}\)}{\((0;\infty )\)}{\(\mathbb{R}\setminus \{2\}\)}{}}
\LANGsk{
\Question{
Určte obor hodnôt funkcie \(f(x) = 1 + \left | \frac{1}
{2(x-2)}\right |\).}
{\((1;\infty )\)}{\(\mathbb{R}\)}{\(\mathbb{R}\setminus \{1\}\)}{\((0;\infty )\)}{\(\mathbb{R}\setminus \{2\}\)}{}}
\LANGes{
\Question{
Determina el rango de la función \(f(x) = 1 + \left | \frac{1}
{2(x-2)}\right |\).}
{\((1;\infty )\)}{\(\mathbb{R}\)}{\(\mathbb{R}\setminus \{1\}\)}{\((0;\infty )\)}{\(\mathbb{R}\setminus \{2\}\)}{}}
\End

\Data\ID{9000007609}
\Updated{2020-07-15 03:07:34}
\Level{C}
\SubArea{Rational functions}
\LANGen{
\Question{
Find the range of the function \(f(x) = 1 + \left | \frac{1}
{|x|+1}\right |\).}
{\((1;2] \)}{\(\mathbb{R}\setminus \{ - 1;1\}\)}{\((0;\infty )\)}{\((1;\infty )\)}{\(\mathbb{R}\)}{}}
\LANGpl{
\Question{
Znajdź zakres funkcji  \(f(x) = 1 + \left | \frac{1}
{|x|+1}\right |\).}
{\((1;2] \)}{\(\mathbb{R}\setminus \{ - 1;1\}\)}{\((0;\infty )\)}{\((1;\infty )\)}{\(\mathbb{R}\)}{}}
\LANGcs{
\Question{
Určete obor hodnot funkce \(f(x) = 1 + \left | \frac{1}
{|x|+1}\right |\).}
{\((1;2\rangle \)}{\(\mathbb{R}\setminus \{ - 1;1\}\)}{\((0;\infty )\)}{\((1;\infty )\)}{\(\mathbb{R}\)}{}}
\LANGsk{
\Question{
Určte obor hodnôt funkcie \(f(x) = 1 + \left | \frac{1}
{|x|+1}\right |\).}
{\((1;2\rangle \)}{\(\mathbb{R}\setminus \{ - 1;1\}\)}{\((0;\infty )\)}{\((1;\infty )\)}{\(\mathbb{R}\)}{}}
\LANGes{
\Question{
Determina el rango de la función  \(f(x) = 1 + \left | \frac{1}
{|x|+1}\right |\).}
{\((1;2] \)}{\(\mathbb{R}\setminus \{ - 1;1\}\)}{\((0;\infty )\)}{\((1;\infty )\)}{\(\mathbb{R}\)}{}}
\End

\Data\ID{9000007610}
\Updated{2020-07-15 03:07:34}
\Level{C}
\SubArea{Rational functions}
\LANGen{
\Question{
Find the range of the function \(f(x) = 1 + \left | \frac{1}
{-|x|+1}\right |\).}
{\((1;\infty )\)}{\(\mathbb{R}\setminus \{ - 1\}\)}{\((0;\infty )\)}{\(\mathbb{R}\setminus \{ - 1;1\}\)}{\(\mathbb{R}\)}{}}
\LANGpl{
\Question{
Znajdź zakres funkcji \(f(x) = 1 + \left | \frac{1}
{-|x|+1}\right |\).}
{\((1;\infty )\)}{\(\mathbb{R}\setminus \{ - 1\}\)}{\((0;\infty )\)}{\(\mathbb{R}\setminus \{ - 1;1\}\)}{\(\mathbb{R}\)}{}}
\LANGcs{
\Question{
Určete obor hodnot funkce \(f(x) = 1 + \left | \frac{1}
{-|x|+1}\right |\).}
{\((1;\infty )\)}{\(\mathbb{R}\setminus \{ - 1\}\)}{\((0;\infty )\)}{\(\mathbb{R}\setminus \{ - 1;1\}\)}{\(\mathbb{R}\)}{}}
\LANGsk{
\Question{
Určte obor hodnôt funkcie \(f(x) = 1 + \left | \frac{1}
{-|x|+1}\right |\).}
{\((1;\infty )\)}{\(\mathbb{R}\setminus \{ - 1\}\)}{\((0;\infty )\)}{\(\mathbb{R}\setminus \{ - 1;1\}\)}{\(\mathbb{R}\)}{}}
\LANGes{
\Question{
Determina el rango de la función  \(f(x) = 1 + \left | \frac{1}
{-|x|+1}\right |\).}
{\((1;\infty )\)}{\(\mathbb{R}\setminus \{ - 1\}\)}{\((0;\infty )\)}{\(\mathbb{R}\setminus \{ - 1;1\}\)}{\(\mathbb{R}\)}{}}
\End

\Data\ID{9000008008}
\Updated{2020-07-15 03:07:34}
\Level{C}
\SubArea{Rational functions}
\LANGen{
\Question{
Given the functions 
\[ 
 \text{$f(x) = -\frac{2}
{x}$ and $g(x)= \frac{k} 
{x}$}
\]
 find the value of the parameter \(k\in \mathbb{R}\setminus \{0\}\) 
which ensures 
\[ 
 g(2) = 2f(-2).
\]}
{\(4\)}{\(2\)}{\(- 1\)}{\(- 2\)}{}{}}
\LANGpl{
\Question{
Dane są funkcje
\[ 
 \text{$f\colon y = -\frac{2}
{x}$ i $g\colon y = \frac{k} 
{x}$}.
\] Wyznacz wartości parametru \(k\in \mathbb{R}\setminus \{0\}\) 
które zapewnia, że 
\[ 
 g(2) = 2f(-2).
\]}
{\(4\)}{\(2\)}{\(- 1\)}{\(- 2\)}{}{}}
\LANGcs{
\Question{
Jsou dány funkce \(f\colon y = -\frac{2}
{x}\) 
a \(g\colon y = \frac{k} 
{x}\). Pro
které \(k\in \mathbb{R}\setminus \{0\}\) 
platí, že \(g(2) = 2f(-2)\)?}
{\(4\)}{\(2\)}{\(- 1\)}{\(- 2\)}{}{}}
\LANGsk{
\Question{
Dané sú funkcie \(f\colon y = -\frac{2}
{x}\) 
a \(g\colon y = \frac{k} 
{x}\). Nájdite hodnotu parametra \(k\in \mathbb{R}\setminus \{0\}\), 
aby platilo \(g(2) = 2f(-2)\).}
{\(4\)}{\(2\)}{\(- 1\)}{\(- 2\)}{}{}}
\LANGes{
\Question{
Dadas las funciones 
\[ 
 \text{$f(x) = -\frac{2}
{x}$ y $g(x)= \frac{k} 
{x}$}
\]
encuentra el valor del parámetro \(k\in \mathbb{R}\setminus \{0\}\) 
que asegura que
\[ 
 g(2) = 2f(-2).
\]}
{\(4\)}{\(2\)}{\(- 1\)}{\(- 2\)}{}{}}
\End

\Data\ID{9000009901}
\Updated{2022-04-23 05:04:17}
\Level{C}
\SubArea{Rational functions}
\IMG{000099_01-figure0.pdf}
\LANGen{
\Question{
The picture shows parts of the graphs of the functions 
\[ 
 \text{$f(x)= \frac{k_{1}} 
 {x} $ and $g(x) = \frac{k_{2}} 
 {x} $.}
\]
Find the relationship between \(k_{1}\) 
and \(k_{2}\)?}
{\(k_{1} > k_{2}\)}{\(k_{1} < k_{2}\)}{\(k_{1} = k_{2}\)}{No conclusion is possible, more of the above possibilities may occur.}{}{}}
\LANGpl{
\Question{
Rysunek przedstawia części wykresów funkcji
\[ 
 \text{$f\colon y = \frac{k_{1}} 
 {x} $ i $g\colon y = \frac{k_{2}} 
 {x} $.}
\]
Jaka jest zależność pomiędzy \(k_{1}\) 
i \(k_{2}\)?}
{\(k_{1} > k_{2}\)}{\(k_{1} < k_{2}\)}{\(k_{1} = k_{2}\)}{Żaden z wniosków nie jest możliwy, istnieje więcej z powyższych możliwości.}{}{}}
\LANGcs{
\Question{
Na obrázku jsou části grafů funkcí
\(f\colon y = \frac{k_{1}} 
 {x} \) a
\(g\colon y = \frac{k_{2}} 
 {x} \).
V jakém vzájemném vztahu jsou oba koeficienty
\(k_{1}\) a
\(k_{2}\)?}
{\(k_{1} > k_{2}\)}{\(k_{1} < k_{2}\)}{\(k_{1} = k_{2}\)}{Vztah mezi \(k_{1}\) 
a \(k_{2}\) 
není možné z obrázku určit.}{}{}}
\LANGsk{
\Question{
Na obrázku sú časti grafov funkcií
\(f\colon y = \frac{k_{1}} 
 {x} \) a
\(g\colon y = \frac{k_{2}} 
 {x} \).
Určte vzájomný vzťah koeficientov
\(k_{1}\) a
\(k_{2}\).}
{\(k_{1} > k_{2}\)}{\(k_{1} < k_{2}\)}{\(k_{1} = k_{2}\)}{Vzťah medzi \(k_{1}\) 
a \(k_{2}\) 
nie je možné z obrázku určiť.}{}{}}
\LANGes{
\Question{
La imagen muestra partes de las gráficas de las funciones.
\[ 
 \text{$f(x)= \frac{k_{1}} 
 {x} $ y $g(x) = \frac{k_{2}} 
 {x} $.}
\]
Encuentra la relación entre \(k_{1}\) 
y \(k_{2}\)}
{\(k_{1} > k_{2}\)}{\(k_{1} < k_{2}\)}{\(k_{1} = k_{2}\)}{La relación entre \(k_{1}\) 
y \(k_{2}\) no se puede determinar a base a la imagen.}{}{}}
\End

\Data\ID{9000009904}
\Updated{2021-04-09 07:04:02}
\Level{C}
\SubArea{Rational functions}
\LANGen{
\Question{
Consider the functions 
\[ 
 \text{$f(x)= \frac{2} 
{x}$ and $g(x) = \frac{2x} 
{x^{2}} $.}
\]
Is the conclusion \(f = g\) 
valid?}
{Yes.}{No.}{Yes, but only on \(\mathbb{R}^{+}\).}{Yes, but only on \(\mathbb{R}^{-}\).}{}{}}
\LANGpl{
\Question{
Rozważ funkcje
\[ 
 \text{$f\colon y = \frac{2} 
{x}$ i $g\colon y = \frac{2x} 
{x^{2}} $.}
\]
Czy wniosek \(f = g\) 
jest zasadny?}
{Tak.}{Nie.}{Tak,  ale tylko dla  \(\mathbb{R}^{+}\).}{Tak, ale tylko dla \(\mathbb{R}^{-}\).}{}{}}
\LANGcs{
\Question{
Jsou dány funkce \(f\colon y = \frac{2} 
{x}\) 
a \(g\colon y = \frac{2x} 
{x^{2}} \). Je možno
tvrdit, že \(f = g\)?}
{Ano.}{Ne.}{Ano, ale pouze pro \(\mathbb{R}^{+}\).}{Ano, ale pouze pro \(\mathbb{R}^{-}\).}{}{}}
\LANGsk{
\Question{
Dané sú funkcie \(f\colon y = \frac{2} 
{x}\) 
a \(g\colon y = \frac{2x} 
{x^{2}} \). Je tvrdenie \(f = g\) pravdivé?}
{Áno.}{Nie.}{Áno, ale len pre \(\mathbb{R}^{+}\).}{Áno, ale len pre \(\mathbb{R}^{-}\).}{}{}}
\LANGes{
\Question{
Considera las funciones
\[ 
 \text{$f(x)= \frac{2} 
{x}$ y $g(x) = \frac{2x} 
{x^{2}} $.}
\]
¿Es la afirmación \(f = g\) 
 válida?}
{Sí.}{No.}{Sí, pero solo en \(\mathbb{R}^{+}\).}{Sí, pero solo en \(\mathbb{R}^{-}\).}{}{}}
\End

\Data\ID{9000009905}
\Updated{2022-04-23 05:04:17}
\Level{C}
\SubArea{Rational functions}
\LANGen{
\Question{
Consider the functions 
\[ 
 \text{$f(x)= \frac{1} 
{2x}$ and $g(x) = \frac{k} 
{x}$.}
\]
Identify the value of the coefficient \(k\) 
which ensures that the graphs of both functions are symmetric about
\(x\)-axis.}
{\(-\frac{1} 
{2}\)}{\(2\)}{\(- 2\)}{\(\frac{1}
{2}\)}{}{}}
\LANGpl{
\Question{
Rozważ funkcje 
\[ 
 \text{$f\colon y = \frac{1} 
{2x}$ i $g\colon y = \frac{k} 
{x}$.}
\]
Wyznacz wartość współczynnika \(k\), który gwarantuje, że wykresy obu funkcji są symetryczne względem osi współrzędnych 
\(x\).}
{\(-\frac{1} 
{2}\)}{\(2\)}{\(- 2\)}{\(\frac{1}
{2}\)}{}{}}
\LANGcs{
\Question{
Jsou dány funkce \(f\colon y = \frac{1} 
{2x}\) a
\(g\colon y = \frac{k} 
{x}\). Jaké hodnoty musí
nabývat koeficient \(k\),
aby grafy obou funkcí byly souměrné podle osy
\(x\)?}
{\(-\frac{1} 
{2}\)}{\(2\)}{\(- 2\)}{\(\frac{1}
{2}\)}{}{}}
\LANGsk{
\Question{
Dané sú funkcie \(f\colon y = \frac{1} 
{2x}\) a
\(g\colon y = \frac{k} 
{x}\). Určte hodnotu parametra \(k\),
aby grafy funkcií \(f\) a \(g\) boli osovo súmerné podľa osi
\(x\).}
{\(-\frac{1} 
{2}\)}{\(2\)}{\(- 2\)}{\(\frac{1}
{2}\)}{}{}}
\LANGes{
\Question{
Considera las funciones
\[ 
 \text{$f(x)= \frac{1} 
{2x}$ y $g(x) = \frac{k} 
{x}$.}
\]
Identifica el valor del coeficiente \(k\)
que asegura que las gráficas de ambas funciones sean simétricas
respecto al eje \(x\).}
{\(-\frac{1} 
{2}\)}{\(2\)}{\(- 2\)}{\(\frac{1}
{2}\)}{}{}}
\End

\Data\ID{9000009906}
\Updated{2021-04-10 09:04:47}
\Level{C}
\SubArea{Rational functions}
\LANGen{
\Question{
Consider a function 
\[ 
 f(x) = \frac{k} 
{x}
\]
with a nonzero real parameter \(k\).
Describe what happens with the function
\(f\) if the
coefficient \(k\) 
changes the sign.}
{The function changes the type of monotonicity on the sets
\(\mathbb{R}^{+}\) and
\(\mathbb{R}^{-}\) 
(either from an increasing function into a decreasing function or vice versa).}{The function changes its parity (from an odd function into an even function or from
an even function into an odd function).}{The domain of the function changes.}{None of the above, both functions have the same parity, monotonicity and domain.}{}{}}
\LANGpl{
\Question{
Rozważ funkcję 
\[ 
 f\colon y = \frac{k} 
{x}
\]
z niezerowym rzeczywistym parametrem \(k\).
Określ co stanie się z funkcją
\(f\) jeśli współczynnik \(k\) 
zmieni znak.}
{Funkcja zmieni rodzaj monotoniczności w zbiorach
\(\mathbb{R}^{+}\) i 
\(\mathbb{R}^{-}\) 
(z funkcji rosnącej na  funkcję malejącą  i odwrotnie).}{Funkcja zmieni parzystość (z funkcji nieparzystej na parzystą i odwrotnie).}{Zmieni się dziedzina funkcji.}{Żadne z powyższych. Obydwie funkcje mają tę samą parzystość, monotoniczność i dziedzinę.}{}{}}
\LANGcs{
\Question{
Je dána funkce 
\[f\colon y = \frac{k} {x}\]
s nenulovým reálným parametrem \(k\). Popište, jaký vliv má 
změna znaménka koeficientu \(k\) na průběh funkce.}
{Funkce se změní v \(\mathbb{R}^{+}\) 
i v \(\mathbb{R}^{-}\) z rostoucí na klesající, nebo naopak.}{Funkce se změní z liché na sudou, nebo naopak.}{Změní se definiční obor funkce.}{Změna znaménka koeficientu \(k\) nemá vliv na sudost-lichost, obor hodnot, ani
monotónnost funkce.}{}{}}
\LANGsk{
\Question{
Daná je funkcia \[f\colon y = \frac{k} {x}\]
pričom \(k\in \mathbb{R}\setminus \{0\}\). Popíšte, aký vplyv má 
zmena znamienka koeficientu \(k\) na priebeh funkcie.}
{Funkcia sa zmení v \(\mathbb{R}^{+}\) 
a v \(\mathbb{R}^{-}\) z rastúcej na klesajúcu, alebo naopak.}{Funkcia sa zmení z nepárnej na párnu, alebo naopak.}{Zmení sa definičný obor funkcie.}{Zmena znamienka koeficientu \(k\) nemá vplyv na paritu, obor hodnôt, ani
monotónnosť funkcie.}{}{}}
\LANGes{
\Question{
Considera la función
\[ 
 f(x) = \frac{k} 
{x}
\]
con un parámetro real distinto de cero \(k\).
Describe lo que sucede con la función
\(f\) si se cambia el signo del parámetro \(k\).}
{La función cambia el tipo de monotonía en los conjuntos
\(\mathbb{R}^{+}\) y
\(\mathbb{R}^{-}\) 
( de una función creciente a una función decreciente o viceversa).}{La función cambia su simetría (de una función impar a una función par o de
una función par a una función impar).}{El dominio de la función cambia.}{Ninguna de las anteriores respuestas, ambas funciones tienen la misma paridad, monotonía y dominio.}{}{}}
\End

\Data\ID{9000009907}
\Updated{2021-04-10 09:04:23}
\Level{C}
\SubArea{Rational functions}
\LANGen{
\Question{
Consider a function 
\[ 
 f(x) = \frac{k} 
{x}
\]
with a nonzero real parameter \(k\).
Suppose that the value of the coefficient
\(k\) changes, but
the sign of \(k\) 
remains the same. Describe which of the properties of
\(f\) is
changed.}
{None of the above, both functions have the same parity, monotonicity and range.}{The function changes its parity (from an odd function into an even function or from
en even function into an odd function).}{The range of the function changes.}{The function changes the type of monotonicity on the sets
\(\mathbb{R}^{+}\) and
\(\mathbb{R}^{-}\) 
(either from an increasing function into a decreasing function or vice versa).}{}{}}
\LANGpl{
\Question{
Rozważ funkcję
\[ 
 f\colon y = \frac{k} 
{x}
\]
z niezerowym rzeczywistym parametrem \(k\).
Przypuśćmy, że wartość współczynnika 
\(k\) zmienia się, ale 
znak liczby \(k\) 
pozostaje taki sam. Określ, która z właściwości funkcji 
\(f\) się zmieni.}
{Żadne z powyższych. Obydwie funkcje mają tę samą monotonność, zakres i parzystość.}{Funkcja zmienia swoja parzystość  (z funkcji nieparzystej na parzystą i odwrotnie).}{Zakres funkcji się zmienia.}{Funkcja zmienia rodzaj monotoniczność w zbiorze 
\(\mathbb{R}^{+}\) i
\(\mathbb{R}^{-}\) 
(z funkcji rosnącej na malejącą i odwrotnie).}{}{}}
\LANGcs{
\Question{
Je dána funkce \[f\colon y = \frac{k} {x}\] 
s reálným nenulovým parametrem \(k\). Popište, jaký vliv má na průběh funkce 
změna velikosti koeficientu \(k\) (při zachování znaménka).}
{Změna velikosti koeficientu \(k\) nemá vliv na sudost-lichost, obor hodnot, ani
monotónnost funkce.}{Funkce se změní z liché na sudou, nebo naopak.}{Změní se obor hodnot funkce.}{Funkce se změní v \(\mathbb{R}^{+}\) 
i v \(\mathbb{R}^{-}\) z rostoucí na klesající, nebo naopak.}{}{}}
\LANGsk{
\Question{
Daná je funkcia \[f\colon y = \frac{k} {x}\]
pričom \(k\in \mathbb{R}\setminus \{0\}\). Popíšte, aký vplyv má 
zmena veľkosti koeficientu \(k\) (pri zachovaní znamienka) na priebeh funkcie.}
{Zmena veľkosti koeficientu \(k\) nemá vplyv na paritu, obor hodnôt, ani monotónnosť funkcie.}{Funkcia sa zmení z nepárnej na párnu, alebo naopak.}{Zmení sa definičný obor funkcie.}{Funkcia sa zmení v \(\mathbb{R}^{+}\) 
a v \(\mathbb{R}^{-}\) z rastúcej na klesajúcu, alebo naopak.}{}{}}
\LANGes{
\Question{
Considera la función
\[ 
 f(x) = \frac{k} 
{x}
\]
con un parámetro real distinto de cero \(k\).
Imagina que el valor del coeficiente 
\(k\) cambia, pero el signo del \(k\) 
sigue siendo el mismo. Describe cuál de las propiedades de
la función \(f\) ha
cambiado.}
{Ninguna de las anteriores respuestas, ambas funciones tienen la misma paridad, monotonía y dominio.}{La función cambia su simetría (de una función impar a una función par o de una función par a una función impar).}{El rango de la función cambia.}{La función cambia el tipo de monotonía en los conjuntos
\(\mathbb{R}^{+}\) y
\(\mathbb{R}^{-}\) 
 (o de una función creciente a una función decreciente o viceversa).}{}{}}
\End

\Data\ID{9000025802}
\Updated{2020-07-15 03:07:34}
\Level{C}
\SubArea{Rational functions}
\LANGen{
\Question{
Find all intersections of the graph of the following function with
\(x\)-axis.
\[ 
 f(x) = \frac{x^{2} + x - 2} 
 {x + 1}
\]}
{\(X_{1} = [-2;0]\),
\(X_{2} = [1;0]\)}{\(X = [0;0]\)}{\(X_{1} = [-2;0]\),
\(X_{2} = [-1;0]\),
\(X_{3} = [1;0]\)}{\(X = [-1;0]\)}{}{}}
\LANGpl{
\Question{
Znajdź wszystkie punkty przecięcia wykresu podanej funkcji z osią 
\(x\) układu współrzędnych.
\[ 
 f\colon y = \frac{x^{2} + x - 2} 
 {x + 1}
\]}
{\(X_{1} = [-2;0]\),
\(X_{2} = [1;0]\)}{\(X = [0;0]\)}{\(X_{1} = [-2;0]\),
\(X_{2} = [-1;0]\),
\(X_{3} = [1;0]\)}{\(X = [-1;0]\)}{}{}}
\LANGcs{
\Question{
Určete všechny společné body osy
\(x\) a grafu
funkce \(f\colon y = \frac{x^{2}+x-2} 
 {x+1} \).}
{\(X_{1} = [-2;0]\),
\(X_{2} = [1;0]\)}{\(X = [0;0]\)}{\(X_{1} = [-2;0]\),
\(X_{2} = [-1;0]\),
\(X_{3} = [1;0]\)}{\(X = [-1;0]\)}{}{}}
\LANGsk{
\Question{
Určte všetky spoločné body osy 
\(x\) a grafu funkcie
\[ 
 f\colon y = \frac{x^{2} + x - 2} 
 {x + 1}.
\]}
{\(X_{1} = [-2;0]\),
\(X_{2} = [1;0]\)}{\(X = [0;0]\)}{\(X_{1} = [-2;0]\),
\(X_{2} = [-1;0]\),
\(X_{3} = [1;0]\)}{\(X = [-1;0]\)}{}{}}
\LANGes{
\Question{
Determina todas las intersecciones de la gráfica de la siguiente función con el eje
\(x\).
\[ 
 f(x) = \frac{x^{2} + x - 2} 
 {x + 1}
\]}
{\(X_{1} = [-2;0]\),
\(X_{2} = [1;0]\)}{\(X = [0;0]\)}{\(X_{1} = [-2;0]\),
\(X_{2} = [-1;0]\),
\(X_{3} = [1;0]\)}{\(X = [-1;0]\)}{}{}}
\End

\Data\ID{9000025803}
\Updated{2020-07-15 03:07:34}
\Level{C}
\SubArea{Rational functions}
\LANGen{
\Question{
Find all intersections of the graph of the following function with
\(x\)-axis.
\[ 
 f(x) = \frac{2x + 1} 
{x^{2} - x - 6}
\]}
{\(X = \left [-\frac{1}
{2};0\right ]\)}{\(X = \left [-\frac{1}
{6};0\right ]\)}{\(X_{1} = [-2;0]\),
\(X_{2} = [3;0]\)}{\(X_{1} = [-2;0]\),
\(X_{2} = \left [-\frac{1}
{2};0\right ]\),
\(X_{3} = [3;0]\)}{}{}}
\LANGpl{
\Question{
Znajdź wszystkie punkty przecięcia wykresu podanej funkcji z osią 
\(x\) układu współrzędnych.
\[ 
 f\colon y = \frac{2x + 1} 
{x^{2} - x - 6}
\]}
{\(X = \left [-\frac{1}
{2};0\right ]\)}{\(X = \left [-\frac{1}
{6};0\right ]\)}{\(X_{1} = [-2;0]\),
\(X_{2} = [3;0]\)}{\(X_{1} = [-2;0]\),
\(X_{2} = \left [-\frac{1}
{2};0\right ]\),
\(X_{3} = [3;0]\)}{}{}}
\LANGcs{
\Question{
Určete všechny společné body osy
\(x\) a grafu
funkce \[f\colon y = \frac{2x+1} 
{x^{2}-x-6}.\]}
{\(X = \left [-\frac{1}
{2};0\right ]\)}{\(X = \left [-\frac{1}
{6};0\right ]\)}{\(X_{1} = [-2;0]\),
\(X_{2} = [3;0]\)}{\(X_{1} = [-2;0]\),
\(X_{2} = \left [-\frac{1}
{2};0\right ]\),
\(X_{3} = [3;0]\)}{}{}}
\LANGsk{
\Question{
Určte všetky spoločné body osy 
\(x\) a grafu funkcie
\[ 
 f\colon y = \frac{2x + 1} 
{x^{2} - x - 6}.
\]}
{\(X = \left [-\frac{1}
{2};0\right ]\)}{\(X = \left [-\frac{1}
{6};0\right ]\)}{\(X_{1} = [-2;0]\),
\(X_{2} = [3;0]\)}{\(X_{1} = [-2;0]\),
\(X_{2} = \left [-\frac{1}
{2};0\right ]\),
\(X_{3} = [3;0]\)}{}{}}
\LANGes{
\Question{
Determina todas las intersecciones de la gráfica de la siguiente función con el eje \(x\).
\[ 
 f(x) = \frac{2x + 1} 
{x^{2} - x - 6}
\]}
{\(X = \left [-\frac{1}
{2};0\right ]\)}{\(X = \left [-\frac{1}
{6};0\right ]\)}{\(X_{1} = [-2;0]\),
\(X_{2} = [3;0]\)}{\(X_{1} = [-2;0]\),
\(X_{2} = \left [-\frac{1}
{2};0\right ]\),
\(X_{3} = [3;0]\)}{}{}}
\End

\Data\ID{9000025806}
\Updated{2020-07-15 03:07:34}
\Level{C}
\SubArea{Rational functions}
\LANGen{
\Question{
In the following list identify a true statement on the function
\(f\).
\[ 
 f(x)= \frac{(3x - 1)(2 - x)} 
 {x + 2}
\]}
{\(f(x) > 0 \iff x\in (-\infty ;-2)\cup \left (\frac{1}
{3};2\right )\)}{\(f(x) > 0 \iff x\in \left (-2; \frac{1} 
{3}\right )\cup (2;\infty )\)}{\(f(x) > 0 \iff x\in \left (-\infty ; \frac{1} 
{3}\right )\cup (2;\infty )\)}{\(f(x) > 0 \iff x\in \left (-\infty ; \frac{1} 
{3}\right )\)}{}{}}
\LANGpl{
\Question{
Które z poniższych stwierdzeń dotyczących funkcji
\(f\) jest prawdziwe? 
\[ 
 f\colon y = \frac{(3x - 1)(2 - x)} 
 {x + 2}
\]}
{\(f(x) > 0 \iff x\in (-\infty ;-2)\cup \left (\frac{1}
{3};2\right )\)}{\(f(x) > 0 \iff x\in \left (-2; \frac{1} 
{3}\right )\cup (2;\infty )\)}{\(f(x) > 0 \iff x\in \left (-\infty ; \frac{1} 
{3}\right )\cup (2;\infty )\)}{\(f(x) > 0 \iff x\in \left (-\infty ; \frac{1} 
{3}\right )\)}{}{}}
\LANGcs{
\Question{
Který z následujících výroků o funkci
\(f\colon y = \frac{(3x-1)(2-x)} 
 {x+2} \) je
pravdivý?}
{\(f(x) > 0 \iff x\in (-\infty ;-2)\cup \left (\frac{1}
{3};2\right )\)}{\(f(x) > 0 \iff x\in \left (-2; \frac{1} 
{3}\right )\cup (2;\infty )\)}{\(f(x) > 0 \iff x\in \left (-\infty ; \frac{1} 
{3}\right )\cup (2;\infty )\)}{\(f(x) > 0 \iff x\in \left (-\infty ; \frac{1} 
{3}\right )\)}{}{}}
\LANGsk{
\Question{
Ktorý z nasledujúcich výrokov o funkcii
\(f\) je pravdivý?
\[ 
 f\colon y = \frac{(3x - 1)(2 - x)} 
 {x + 2}
\]}
{\(f(x) > 0 \iff x\in (-\infty ;-2)\cup \left (\frac{1}
{3};2\right )\)}{\(f(x) > 0 \iff x\in \left (-2; \frac{1} 
{3}\right )\cup (2;\infty )\)}{\(f(x) > 0 \iff x\in \left (-\infty ; \frac{1} 
{3}\right )\cup (2;\infty )\)}{\(f(x) > 0 \iff x\in \left (-\infty ; \frac{1} 
{3}\right )\)}{}{}}
\LANGes{
\Question{
En la siguiente lista, identifica una proposición verdadera sobre la función
\(f\).
\[ 
 f(x)= \frac{(3x - 1)(2 - x)} 
 {x + 2}
\]}
{\(f(x) > 0 \iff x\in (-\infty ;-2)\cup \left (\frac{1}
{3};2\right )\)}{\(f(x) > 0 \iff x\in \left (-2; \frac{1} 
{3}\right )\cup (2;\infty )\)}{\(f(x) > 0 \iff x\in \left (-\infty ; \frac{1} 
{3}\right )\cup (2;\infty )\)}{\(f(x) > 0 \iff x\in \left (-\infty ; \frac{1} 
{3}\right )\)}{}{}}
\End

\Data\ID{9000025808}
\Updated{2020-07-15 03:07:34}
\Level{C}
\SubArea{Rational functions}
\LANGen{
\Question{
In the following list identify a true statement on the function
\(f\).
\[ 
 f(x) = \frac{(x - 1)(x + 2)} 
{(2x + 1)(3 - 2x)}
\]}
{\(f(x) > 0 \iff x\in \left (-2;-\frac{1}
{2}\right )\cup \left (1; \frac{3} 
{2}\right )\)}{\(f(x) > 0 \iff x\in (-\infty ;-2)\cup \left (-\frac{1}
{2};1\right )\cup \left (\frac{3}
{2};\infty \right )\)}{\(f(x) > 0 \iff x\in (-\infty ;-2)\cup (1;\infty )\)}{\(f(x) > 0 \iff x\in \left (-2; \frac{3} 
{2}\right )\)}{}{}}
\LANGpl{
\Question{
Które z poniższych stwierdzeń dotyczących funkcji 
\(f\) jest prawdziwe?
\[ 
 f\colon y = \frac{(x - 1)(x + 2)} 
{(2x + 1)(3 - 2x)}
\]}
{\(f(x) > 0 \iff x\in \left (-2;-\frac{1}
{2}\right )\cup \left (1; \frac{3} 
{2}\right )\)}{\(f(x) > 0 \iff x\in (-\infty ;-2)\cup \left (-\frac{1}
{2};1\right )\cup \left (\frac{3}
{2};\infty \right )\)}{\(f(x) > 0 \iff x\in (-\infty ;-2)\cup (1;\infty )\)}{\(f(x) > 0 \iff x\in \left (-2; \frac{3} 
{2}\right )\)}{}{}}
\LANGcs{
\Question{
Který z následujících výroků o funkci \( f \) je
pravdivý?
\[f\colon y = \frac{(x-1)(x+2)} 
{(2x+1)(3-2x)}\]}
{\(f(x) > 0 \iff x\in \left (-2;-\frac{1}
{2}\right )\cup \left (1; \frac{3} 
{2}\right )\)}{\(f(x) > 0 \iff x\in (-\infty ;-2)\cup \left (-\frac{1}
{2};1\right )\cup \left (\frac{3}
{2};\infty \right )\)}{\(f(x) > 0 \iff x\in (-\infty ;-2)\cup (1;\infty )\)}{\(f(x) > 0 \iff x\in \left (-2; \frac{3} 
{2}\right )\)}{}{}}
\LANGsk{
\Question{
Ktorý z nasledujúcich výrokov o funkcii
\(f\) je pravdivý?
\[ 
 f\colon y = \frac{(x - 1)(x + 2)} 
{(2x + 1)(3 - 2x)}
\]}
{\(f(x) > 0 \iff x\in \left (-2;-\frac{1}
{2}\right )\cup \left (1; \frac{3} 
{2}\right )\)}{\(f(x) > 0 \iff x\in (-\infty ;-2)\cup \left (-\frac{1}
{2};1\right )\cup \left (\frac{3}
{2};\infty \right )\)}{\(f(x) > 0 \iff x\in (-\infty ;-2)\cup (1;\infty )\)}{\(f(x) > 0 \iff x\in \left (-2; \frac{3} 
{2}\right )\)}{}{}}
\LANGes{
\Question{
En la siguiente lista, identifica una proposición verdadera sobre la función
\(f\).
\[ 
 f(x) = \frac{(x - 1)(x + 2)} 
{(2x + 1)(3 - 2x)}
\]}
{\(f(x) > 0 \iff x\in \left (-2;-\frac{1}
{2}\right )\cup \left (1; \frac{3} 
{2}\right )\)}{\(f(x) > 0 \iff x\in (-\infty ;-2)\cup \left (-\frac{1}
{2};1\right )\cup \left (\frac{3}
{2};\infty \right )\)}{\(f(x) > 0 \iff x\in (-\infty ;-2)\cup (1;\infty )\)}{\(f(x) > 0 \iff x\in \left (-2; \frac{3} 
{2}\right )\)}{}{}}
\End

\Data\ID{9000025809}
\Updated{2020-07-15 03:07:34}
\Level{C}
\SubArea{Rational functions}
\LANGen{
\Question{
In the following list identify a true statement on the function
\(f\).
\[ 
 f(x)= \frac{(6x - 1)} 
{(x - 2)(3x + 1)}
\]}
{\(f(x)\geq 0 \iff x\in \left (-\frac{1}
{3}; \frac{1} 
{6}\right ] \cup (2;\infty )\)}{\(f(x)\geq 0 \iff x\in \left (-\frac{1}
{3}; \frac{1} 
{6}\right )\cup (2;\infty )\)}{\(f(x)\geq 0 \iff x\in \left (-\infty ;-\frac{1}
{3}\right )\cup \left [ \frac{1}
{6};2\right )\)}{\(f(x)\geq 0 \iff x\in \left [ -\frac{1}
{3}; \frac{1} 
{6}\right ] \cup (2;\infty )\)}{}{}}
\LANGpl{
\Question{
Które z poniższych stwierdzeń dotyczących funkcji 
\(f\) jest prawdziwe?
\[ 
 f\colon y = \frac{(6x - 1)} 
{(x - 2)(3x + 1)}
\]}
{\(f(x)\geq 0 \iff x\in \left (-\frac{1}
{3}; \frac{1} 
{6}\right ] \cup (2;\infty )\)}{\(f(x)\geq 0 \iff x\in \left (-\frac{1}
{3}; \frac{1} 
{6}\right )\cup (2;\infty )\)}{\(f(x)\geq 0 \iff x\in \left (-\infty ;-\frac{1}
{3}\right )\cup \left [ \frac{1}
{6};2\right )\)}{\(f(x)\geq 0 \iff x\in \left [ -\frac{1}
{3}; \frac{1} 
{6}\right ] \cup (2;\infty )\)}{}{}}
\LANGcs{
\Question{
Který z následujících výroků o funkci \( f \) je pravdivý?
\[f\colon y = \frac{(6x-1)} 
{(x-2)(3x+1)}\]}
{\(f(x)\geq 0 \iff x\in \left (-\frac{1}
{3}; \frac{1} 
{6}\right \rangle \cup (2;\infty )\)}{\(f(x)\geq 0 \iff x\in \left (-\frac{1}
{3}; \frac{1} 
{6}\right )\cup (2;\infty )\)}{\(f(x)\geq 0 \iff x\in \left (-\infty ;-\frac{1}
{3}\right )\cup \left \langle \frac{1}
{6};2\right )\)}{\(f(x)\geq 0 \iff x\in \left \langle -\frac{1}
{3}; \frac{1} 
{6}\right \rangle \cup (2;\infty )\)}{}{}}
\LANGsk{
\Question{
Ktorý z nasledujúcich výrokov o funkcii
\(f\) je pravdivý?
\[ 
 f\colon y = \frac{(6x - 1)} 
{(x - 2)(3x + 1)}
\]}
{\(f(x)\geq 0 \iff x\in \left (-\frac{1}
{3}; \frac{1} 
{6}\right ] \cup (2;\infty )\)}{\(f(x)\geq 0 \iff x\in \left (-\frac{1}
{3}; \frac{1} 
{6}\right )\cup (2;\infty )\)}{\(f(x)\geq 0 \iff x\in \left (-\infty ;-\frac{1}
{3}\right )\cup \left [ \frac{1}
{6};2\right )\)}{\(f(x)\geq 0 \iff x\in \left [ -\frac{1}
{3}; \frac{1} 
{6}\right ] \cup (2;\infty )\)}{}{}}
\LANGes{
\Question{
En la siguiente lista, identifica una proposición verdadera sobre la función
\(f\).
\[ 
 f(x)= \frac{(6x - 1)} 
{(x - 2)(3x + 1)}
\]}
{\(f(x)\geq 0 \iff x\in \left (-\frac{1}
{3}; \frac{1} 
{6}\right ] \cup (2;\infty )\)}{\(f(x)\geq 0 \iff x\in \left (-\frac{1}
{3}; \frac{1} 
{6}\right )\cup (2;\infty )\)}{\(f(x)\geq 0 \iff x\in \left (-\infty ;-\frac{1}
{3}\right )\cup \left [ \frac{1}
{6};2\right )\)}{\(f(x)\geq 0 \iff x\in \left [ -\frac{1}
{3}; \frac{1} 
{6}\right ] \cup (2;\infty )\)}{}{}}
\End

\Data\ID{9000025810}
\Updated{2020-07-15 03:07:34}
\Level{C}
\SubArea{Rational functions}
\LANGen{
\Question{
In the following list identify a true statement on the function
\(f\).
\[ 
 f(x) = \frac{(x - 2)(3 - x)} 
{(2x - 1)(3x - 1)}
\]}
{\(f(x)\geq 0 \iff x\in \left (\frac{1}
{3}; \frac{1} 
{2}\right )\cup [ 2;3] \)}{\(f(x)\geq 0 \iff x\in \left [ \frac{1}
{3}; \frac{1} 
{2}\right ] \cup [ 2;3] \)}{\(f(x)\geq 0 \iff x\in \left (-\infty ; \frac{1} 
{3}\right )\cup \left [ \frac{1}
{2};2\right ] \cup [ 3;\infty )\)}{\(f(x)\geq 0 \iff x\in \left (\frac{1}
{3}; \frac{1} 
{2}\right )\cup (2;3)\)}{}{}}
\LANGpl{
\Question{
Które z poniższych stwierdzeń dotyczących funkcji 
\(f\) jest prawdziwe?
\[ 
 f\colon y = \frac{(x - 2)(3 - x)} 
{(2x - 1)(3x - 1)}
\]}
{\(f(x)\geq 0 \iff x\in \left (\frac{1}
{3}; \frac{1} 
{2}\right )\cup [ 2;3] \)}{\(f(x)\geq 0 \iff x\in \left [ \frac{1}
{3}; \frac{1} 
{2}\right ] \cup [ 2;3] \)}{\(f(x)\geq 0 \iff x\in \left (-\infty ; \frac{1} 
{3}\right )\cup \left [ \frac{1}
{2};2\right ] \cup [ 3;\infty )\)}{\(f(x)\geq 0 \iff x\in \left (\frac{1}
{3}; \frac{1} 
{2}\right )\cup (2;3)\)}{}{}}
\LANGcs{
\Question{
Který z následujících výroků o funkci \( f \) je pravdivý?
\[f\colon y = \frac{(x-2)(3-x)} 
{(2x-1)(3x-1)}\]}
{\(f(x)\geq 0 \iff x\in \left (\frac{1}
{3}; \frac{1} 
{2}\right )\cup \langle 2;3\rangle \)}{\(f(x)\geq 0 \iff x\in \left \langle \frac{1}
{3}; \frac{1} 
{2}\right \rangle \cup \langle 2;3\rangle \)}{\(f(x)\geq 0 \iff x\in \left (-\infty ; \frac{1} 
{3}\right )\cup \left \langle \frac{1}
{2};2\right \rangle \cup \langle 3;\infty )\)}{\(f(x)\geq 0 \iff x\in \left (\frac{1}
{3}; \frac{1} 
{2}\right )\cup (2;3)\)}{}{}}
\LANGsk{
\Question{
Ktorý z nasledujúcich výrokov o funkcii
\(f\) je pravdivý?
\[ 
 f\colon y = \frac{(x - 2)(3 - x)} 
{(2x - 1)(3x - 1)}
\]}
{\(f(x)\geq 0 \iff x\in \left (\frac{1}
{3}; \frac{1} 
{2}\right )\cup [ 2;3] \)}{\(f(x)\geq 0 \iff x\in \left [ \frac{1}
{3}; \frac{1} 
{2}\right ] \cup [ 2;3] \)}{\(f(x)\geq 0 \iff x\in \left (-\infty ; \frac{1} 
{3}\right )\cup \left [ \frac{1}
{2};2\right ] \cup [ 3;\infty )\)}{\(f(x)\geq 0 \iff x\in \left (\frac{1}
{3}; \frac{1} 
{2}\right )\cup (2;3)\)}{}{}}
\LANGes{
\Question{
En la siguiente lista, identifica una proposición verdadera sobre la función
\(f\).
\[ 
 f(x) = \frac{(x - 2)(3 - x)} 
{(2x - 1)(3x - 1)}
\]}
{\(f(x)\geq 0 \iff x\in \left (\frac{1}
{3}; \frac{1} 
{2}\right )\cup [ 2;3] \)}{\(f(x)\geq 0 \iff x\in \left [ \frac{1}
{3}; \frac{1} 
{2}\right ] \cup [ 2;3] \)}{\(f(x)\geq 0 \iff x\in \left (-\infty ; \frac{1} 
{3}\right )\cup \left [ \frac{1}
{2};2\right ] \cup [ 3;\infty )\)}{\(f(x)\geq 0 \iff x\in \left (\frac{1}
{3}; \frac{1} 
{2}\right )\cup (2;3)\)}{}{}}
\End
